\documentclass[UTF8,a4paper,10pt]{ctexart}
\usepackage{amsmath,amssymb,geometry,listings,xcolor,hyperref,enumitem,fancyhdr,needspace}
\geometry{margin=1.6cm,headheight=15pt}
\hypersetup{hidelinks,pdftitle={Otomachi Una C++ 模板库XCPC 考场代码手册}}
\setlist[itemize]{leftmargin=1.6em,itemsep=0.25em,topsep=0.3em,parsep=0pt}
\sloppy\setlength{\emergencystretch}{3em}
\setlength{\parindent}{0pt}\setlength{\parskip}{0.4em}
\pagestyle{fancy}\fancyhf{}\fancyhead[L]{Otomachi Una 模板库}\fancyhead[R]{XCPC 考场代码手册}\fancyfoot[C]{\thepage}
\lstset{language=C++,basicstyle=\ttfamily\fontsize{9}{10.5}\selectfont,keywordstyle=\color{black},commentstyle=\color{black!70},frame=none,numbers=left,numberstyle=\tiny\color{black!50},numbersep=5pt,xleftmargin=1.4em,rulecolor=\color{black!25},breaklines=true,columns=fullflexible,keepspaces=true,tabsize=2,showstringspaces=false}
\newcommand{\TemplUsage}[1]{\textbf{什么时候用}\par #1\par}
\newcommand{\TemplApi}[1]{\textbf{调用顺序与接口}\par #1\par}
\newcommand{\TemplNote}[1]{\textbf{使用边界}\par #1\par}
\title{Otomachi Una C++ 模板库\\XCPC 考场代码手册}
\author{Junlin Ye}
\date{2026-09-23}
\begin{document}
\begin{titlepage}\maketitle\thispagestyle{empty}
\vspace{1cm}
收录当前模板库的全部 72 份主模板。每节给出调用方式、最小示例和完整代码；模板有依赖时，先抄依赖，再抄本节代码。标题中的路径用于对应源码，实际抄写以本手册代码为准。

\textbf{考场版改动：}SA 使用计数排序倍增，构建 $O(n\log n)$；poly-fast、poly-new 和 MTT 使用普通二进制位逆序迭代 NTT，去掉大段预制根表和展开蝶形。其余算法保留原实现。源码目录中的高效版本不受影响。

\textbf{选择多项式：}常规卷积及 ln/exp 用 poly-fast；需要开方、带余除法或多点求值选 poly-new；非 NTT 目标模数选 MTT。poly-old 为旧接口，保留备查。各节代码完整，模整数定义也分别保留，使用时只能选其中一份。

\textbf{编译与调试：}GNU C++17。basic/src.cpp 是含头文件和 main 的底稿，其他代码抄在 main 前。LOCAL 控制 Otomachi\_Una.in/out，DEBUG 控制 cerr、debug、debugl。不要把全部模板拼成一个程序。

\textbf{抄写提示：}行号不属于代码，排版换行不代表源码换行。长模板会跨页，按页码接续。清空对象使用 set/setN，几何与点分治用 output 调试；FastIO 仅适用于 Linux 的普通文件输入，额外头文件见对应节。

\vfill
使用顺序：按目录选模板 → 核对下标和依赖 → 抄写代码 → 跑最小示例 → 替换题目逻辑。
\end{titlepage}
{\small\tableofcontents}\clearpage
\section{基础与数值类型}
\Needspace{10\baselineskip}
\subsection{\texttt{basic/src.cpp}}
\textbf{公共底稿}

\TemplUsage{提供类型别名、调试宏和其他模板依赖的辅助函数。自带标准头文件和 main，可直接作为提交程序的底稿。}
\TemplApi{\begin{itemize}
\item 直接复制本文件，在 \texttt{main} 前粘贴算法，并在其中填写输入输出与求解逻辑。
\item \texttt{debug(...) / debugl(tag,...)} 接受 1 至 10 个参数；编译时加 \texttt{-DDEBUG} 才输出，关闭时参数不求值。
\item \texttt{chkmin/chkmax} 更新最值并返回是否修改。\texttt{first\_bit/last\_bit} 返回最低/最高置位的位置，参数必须非零。
\item \texttt{inf<int>}、\texttt{inf<ll>} 已特化；其他类型默认值为零，不能当无穷大。
\end{itemize}}
\textbf{最小示例（放入 main）：}
\begin{lstlisting}
int a=7; chkmin(a,3);
assert(a==3 && first_bit(12)==2 && last_bit(12)==3);
debug(a);
cerr << "diagnostic\n";
\end{lstlisting}
\TemplNote{定义 \texttt{LOCAL} 时从 \texttt{Otomachi\_Una.in} 读入并写入 \texttt{Otomachi\_Una.out}，打开失败返回 1；否则使用标准输入输出。\texttt{DEBUG} 独立控制调试输出。未定义时通过宏屏蔽裸写的 \texttt{cerr}，不要写 \texttt{std::cerr}。向量输入前必须先分配非空长度；向量输出带方括号并截断长向量，正式答案应自己循环输出。}

\Needspace{6\baselineskip}\textbf{完整代码}
\begin{lstlisting}[firstnumber=1]
#include<bits/stdc++.h>
using namespace std;
using ll=long long;
using ull=unsigned long long;
// head file
#ifndef DEBUG
	#define cerr for(;false;) cerr
#endif
template<typename T>
inline constexpr T inf=T{};
template<>
inline constexpr int inf<int> = 1010000000;
template<>
inline constexpr long long inf<long long> = 2020000000000000000LL;
template<typename T,size_t S>
istream& operator >>(istream &o,array<T,S> &I){
	assert(S>0);
	for(int i=0;i<S;i++) o>>I[i];
	return o;
}
template<typename T,size_t S>
ostream& operator <<(ostream &o,const array<T,S> &I){
	o<<"[";
	for(int i=0;i<S;i++){
		if(i==50){o<<"...(total "<<S<<" elements)";break;}
		o<<I[i];
		if(i+1<I.size()) o<<",";
	}
	o<<"]";
	return o;
}
template<typename T>
istream& operator >>(istream &o,vector<T> &I){
	assert(!I.empty());
	for(int i=0;i<I.size();i++) o>>I[i];
	return o;
}
template<typename T>
ostream& operator <<(ostream &o,const vector<T> &I){
	o<<"[";
	for(int i=0;i<I.size();i++){
		if(i==50){o<<"...(total "<<I.size()<<" elements)";break;}
		o<<I[i];
		if(i+1<I.size()) o<<",";
	}
	o<<"]";
	return o;
}

#define DEBUG_1(x) #x"="<<(x)
#define DEBUG_2(x,y) DEBUG_1(x)<<", "<<DEBUG_1(y)
#define DEBUG_3(x,y,z) DEBUG_2(x,y)<<", "<<DEBUG_1(z)
#define DEBUG_4(x,y,z,w) DEBUG_3(x,y,z)<<", "<<DEBUG_1(w)
#define DEBUG_5(x,y,z,w,a) DEBUG_4(x,y,z,w)<<", "<<DEBUG_1(a)
#define DEBUG_6(x,y,z,w,a,b) DEBUG_5(x,y,z,w,a)<<", "<<DEBUG_1(b)
#define DEBUG_7(x,y,z,w,a,b,c) DEBUG_6(x,y,z,w,a,b)<<", "<<DEBUG_1(c)
#define DEBUG_8(x,y,z,w,a,b,c,d) DEBUG_7(x,y,z,w,a,b,c)<<", "<<DEBUG_1(d)
#define DEBUG_9(x,y,z,w,a,b,c,d,e) DEBUG_8(x,y,z,w,a,b,c,d)<<", "<<DEBUG_1(e)
#define DEBUG_10(x,y,z,w,a,b,c,d,e,f) DEBUG_9(x,y,z,w,a,b,c,d,e)<<", "<<DEBUG_1(f)
#define GET_MACRO(_1,_2,_3,_4,_5,_6,_7,_8,_9,_10,NAME,...) NAME
#ifdef DEBUG
#define debug(...) do { \
	cerr<<"(line "<< __LINE__ <<"): "; \
	cerr<<GET_MACRO(__VA_ARGS__,DEBUG_10,DEBUG_9,DEBUG_8, \
						   DEBUG_7,DEBUG_6,DEBUG_5,DEBUG_4, \
						   DEBUG_3,DEBUG_2,DEBUG_1)(__VA_ARGS__); \
	cerr<<endl; \
} while(0)
#define debugl(s,...) do { \
	cerr<<"(line "<< __LINE__ <<"): <"<<s<<"> "; \
	cerr<<GET_MACRO(__VA_ARGS__,DEBUG_10,DEBUG_9,DEBUG_8, \
						   DEBUG_7,DEBUG_6,DEBUG_5,DEBUG_4, \
						   DEBUG_3,DEBUG_2,DEBUG_1)(__VA_ARGS__); \
	cerr<<endl; \
} while(0)
#else
#define debug(...) ((void)0)
#define debugl(...) ((void)0)
#endif
inline int last_bit(ull x){assert(x);return 63-__builtin_clzll(x);}
inline int first_bit(ull x){assert(x);return __builtin_ctzll(x);}
inline int popc(ull x){return __builtin_popcountll(x);}
template<typename A,typename B> inline bool chkmax(A &x,const B &y){if(x<y){x=y;return true;}return false;}
template<typename A,typename B> inline bool chkmin(A &x,const B &y){if(x>y){x=y;return true;}return false;}
int main(){
	#ifdef LOCAL
		if(!freopen("Otomachi_Una.in","r",stdin)) return 1;
		if(!freopen("Otomachi_Una.out","w",stdout)) return 1;
	#endif
	ios::sync_with_stdio(false);cin.tie(0);cout.tie(0);
	return 0;
}
\end{lstlisting}
\par\vspace{1em}

\Needspace{10\baselineskip}
\subsection{\texttt{basic/ST table.cpp}}
\textbf{静态区间最值}

\TemplUsage{数组不修改、需要多次查询闭区间最大值时使用。预处理 O(n log n)，查询 O(1)，空间 O(n log n)。}
\TemplApi{\begin{itemize}
\item \texttt{ST\_table t; t.setN(n)} 分配下标 1..n，要求 n>=1。
\item 填入 \texttt{t.a[0][i]} 后调用 \texttt{t.build()}。
\item \texttt{t.ask(l,r)} 返回最大值；要维护最小值，先把 \texttt{chk} 改成 \texttt{min}。
\end{itemize}}
\textbf{最小示例（放入 main）：}
\begin{lstlisting}
ST_table t; t.setN(3);
t.a[0][1]=4; t.a[0][2]=1; t.a[0][3]=8;
t.build(); assert(t.ask(1,2)==4);
\end{lstlisting}
\TemplNote{修改原数组后必须重新 build。文件末尾已有全局对象 \texttt{S}，可以直接用它，也可以另建局部对象。}

\Needspace{6\baselineskip}\textbf{完整代码}
\begin{lstlisting}[firstnumber=1]
//
// template for ST table
// version 1.0 (Last Update Jun 23rd, 2026)
//
// usage:
//   S.setN(n); fill S.a[0][1..n]; S.build();
//   int mx = S.ask(l, r);  // static RMQ, tweak chk() for min
//
struct ST_table{
int n=0;
vector<vector<int>> a;
void setN(int _n){
	assert(1<=_n&&_n<INT_MAX);n=_n;
	a.assign(__lg(n)+1,vector<int>(n+1));
}
inline int chk(const int &x,const int &y){return max(x,y);} // compare
void build(){
	assert(n>0);
	for(int i=1;i<(int)a.size();i++) for(int j=1;j<=n-(1<<i)+1;j++) 
		a[i][j]=chk(a[i-1][j],a[i-1][j+(1<<(i-1))]);
}
int ask(int l,int r){
	assert(1<=l&&l<=r&&r<=n);
	int k=__lg(r-l+1);
	return chk(a[k][l],a[k][r-(1<<k)+1]);
}
}S;
\end{lstlisting}\Needspace{10\baselineskip}
\subsection{\texttt{basic/ST table (linear).cpp}}
\textbf{四毛子 RMQ：线性预处理与空间}

\TemplUsage{静态数组最值。采用块内单调栈位掩码、块间 ST 表；在机器字位运算为 O(1) 的模型下，预处理及空间 O(n)，查询 O(1)。}
\TemplApi{\texttt{Linear\_RMQ<int> t; t.build(a)} 接受从 0 开始的 vector，内部保存副本。\texttt{ask(l,r)} 查询闭区间最小值，\texttt{pos(l,r)} 返回最小值下标；相等取最左。\texttt{Linear\_RMQ<int,true>} 改为最大值，其他接口相同。}
\textbf{最小示例（放入 main）：}
\begin{lstlisting}
vector<int> a{4,1,1,3};
Linear_RMQ<int> mn; mn.build(a);
assert(mn.ask(0,3)==1 && mn.pos(0,3)==1);
Linear_RMQ<int,true> mx; mx.build(a);
assert(mx.ask(1,3)==3);
\end{lstlisting}
\TemplNote{块长 B=max(1,floor(log2 n))，不是固定常数；块数 m=ceil(n/B)，块间表 O(m log m)=O(n)。每个元素存一份块前缀的单调栈掩码，一次清除低位和 ctz 即可查询块内最左最值。n<INT\_MAX，块长最多 30，使用 unsigned 位运算。可 build 空数组，但不能查询；修改数组后须重新 build。set() 清空。不需要普通 ST 表或笛卡尔树作为依赖。}

\Needspace{6\baselineskip}\textbf{完整代码}
\begin{lstlisting}[firstnumber=1]
////////////////////////////////////////////////////////////////
//
// template for linear RMQ
//
// usage:
//   Linear_RMQ<int> t; t.build(a); t.ask(l,r); t.pos(l,r);
//   0-indexed, inclusive; use Linear_RMQ<int,true> for maximum.
//   O(n) build / space, O(1) query on the word RAM; ties choose the leftmost.
//
////////////////////////////////////////////////////////////////
template<typename T,bool MAX=false> struct Linear_RMQ{
private:
int n=0,B=1;
vector<T> a;
vector<unsigned> mask;
vector<vector<int>> st;
bool better(int x,int y)const{if constexpr(MAX)return a[y]<a[x];else return a[x]<a[y];}
int chk(int x,int y)const{return better(y,x)?y:x;}
int small(int l,int r)const{
	unsigned s=mask[r]&(~0u<<(l%B));
	return r/B*B+__builtin_ctz(s);
}
public:
void set(){n=0;B=1;a.clear();mask.clear();st.clear();}
void build(const vector<T> &v){
	assert(v.size()<(size_t)INT_MAX);a=v;n=a.size();mask.resize(n);st.clear();
	if(!n){B=1;return;}
	// B grows with log(n); the block sparse table has only O(n) entries.
	B=max(1,31-__builtin_clz((unsigned)n));
	int m=(n-1)/B+1;st.resize(32-__builtin_clz((unsigned)m));st[0].resize(m);
	unsigned s=0;
	for(int i=0;i<n;i++){
		if(i%B==0)s=0;
		while(s){
			int j=31-__builtin_clz(s);
			if(!better(i,i/B*B+j))break;
			s^=1u<<j;
		}
		mask[i]=s|=1u<<(i%B);
		st[0][i/B]=i/B*B+__builtin_ctz(s);
	}
	for(int k=1;k<(int)st.size();k++){
		st[k].resize(m-(1<<k)+1);
		for(int i=0;i<(int)st[k].size();i++)st[k][i]=chk(st[k-1][i],st[k-1][i+(1<<(k-1))]);
	}
}
int pos(int l,int r)const{
	assert(0<=l&&l<=r&&r<n);
	int x=l/B,y=r/B;if(x==y)return small(l,r);
	int ans=small(l,(x+1)*B-1);
	if(x+1<y){
		int k=31-__builtin_clz((unsigned)(y-x-1));
		ans=chk(ans,chk(st[k][x+1],st[k][y-(1<<k)]));
	}
	return chk(ans,small(y*B,r));
}
T ask(int l,int r)const{return a[pos(l,r)];}
};
// end for basic/ST table (linear).cpp
/////////////////////////
\end{lstlisting}

\par\vspace{1em}

\Needspace{10\baselineskip}
\subsection{\texttt{basic/DSU.cpp}}
\textbf{普通并查集}

\TemplUsage{维护集合连通关系，路径压缩配合按大小合并。操作均摊近似 O(1)，空间 O(n)。}
\TemplApi{\begin{itemize}
\item \texttt{DSU d; d.setN(n)} 创建 1..n 的独立集合。
\item \texttt{find(x)} 返回代表元；\texttt{same(x,y)} 判断同集合。
\item \texttt{merge(x,y)} 成功返回新代表元，已经连通返回 0。\texttt{siz[find(x)]} 为集合大小。
\end{itemize}}
\textbf{最小示例（放入 main）：}
\begin{lstlisting}
DSU d; d.setN(4);
d.merge(1,2); assert(d.same(1,2));
assert(d.siz[d.find(1)]==2 && !d.same(1,3));
\end{lstlisting}
\TemplNote{再次 setN 会清空关系。普通版不能回滚；需要撤销时选择 unionfind。}

\Needspace{6\baselineskip}\textbf{完整代码}
\begin{lstlisting}[firstnumber=1]
//
// template for Disjoint Set Union (DSU)
//
// usage:
//   DSU dsu; dsu.setN(n); dsu.merge(u, v);
//   dsu.same(u, v); int r = dsu.find(u);  // merge returns new root or 0
//
struct DSU{
int n=0;
vector<int> fa,siz;
void setN(int _n){
	assert(0<=_n&&_n<INT_MAX);n=_n;
	fa.resize(n+1);iota(fa.begin(),fa.end(),0);siz.assign(n+1,1);siz[0]=0;
}
inline int find(int x){
	assert(1<=x&&x<=n);
	while(x^fa[x])x=fa[x]=fa[fa[x]];return x;
}
inline int merge(int x,int y){
	assert(1<=min(x,y)&&max(x,y)<=n);
	x=find(x),y=find(y);
	if(x==y) return 0;
	if(siz[x]<siz[y]) swap(x,y);
	siz[x]+=siz[y],fa[y]=x;return x;
}
inline bool same(int x,int y){
	assert(1<=min(x,y)&&max(x,y)<=n);
	return find(x)==find(y);
}
};
\end{lstlisting}
\par\vspace{1em}

\Needspace{10\baselineskip}
\subsection{\texttt{basic/int128\_IO.cpp}}
\textbf{128 位整数输入输出}

\TemplUsage{为有符号 \texttt{\_\_int128} 提供十进制流输入输出，仍然是定长整数。}
\TemplApi{\begin{itemize}
\item 类型别名 \texttt{lll} 等于 \texttt{\_\_int128}。
\item 粘贴后可直接使用 \texttt{cin >> x}、\texttt{cout << x}。
\end{itemize}}
\textbf{最小示例（放入 main）：}
\begin{lstlisting}
lll x=1;
x*=1000000000000000000LL; x*=100;
ostringstream out; out << x;
assert(out.str()=="100000000000000000000");
\end{lstlisting}
\TemplNote{需要 GCC/Clang 的 \_\_int128 支持；输入和运算结果都必须在有符号 128 位范围内。乘法之前就要提升类型，例如 \texttt{(\_\_int128)a*b}。}

\Needspace{6\baselineskip}\textbf{完整代码}
\begin{lstlisting}[firstnumber=1]
//
// template for __int128 I/O
//
// usage:
//   using lll = __int128; lll x; cin >> x; cout << x;
//   GCC/Clang only; decimal I/O
//
using lll=__int128;
istream& operator >>(istream &o,lll &x){
	string s;if(!(o>>s))return o;
	bool neg=s[0]=='-';size_t i=(neg||s[0]=='+');
	assert(i<s.size());__uint128_t v=0,lim=((__uint128_t)1<<127)-!neg;
	for(;i<s.size();i++){
		assert('0'<=s[i]&&s[i]<='9');unsigned d=s[i]-'0';
		assert(v<=(lim-d)/10);v=v*10+d;
	}
	x=neg?(lll)(__uint128_t(0)-v):(lll)v;
	return o;
}
ostream& operator <<(ostream &o,lll x){
	array<int,40> a;int n=0;
	if(x==0){o<<0;return o;}
	__uint128_t v=x;
	if(x<0){o<<"-";v=__uint128_t(0)-v;}
	while(v) a[++n]=v%10,v/=10;
	for(int i=n;i>=1;i--) o<<a[i];
	return o;
}
\end{lstlisting}
\par\vspace{1em}

\Needspace{10\baselineskip}
\subsection{\texttt{basic/BigInt.cpp}}
\textbf{高精度有符号整数}

\TemplUsage{十进制大整数，vector 按 10\textasciicircum{}9 分块。加减 O(n)，普通乘法及长除法最坏 O(nm)，n、m 是分块数。}
\TemplApi{\begin{itemize}
\item \texttt{BigInt x=123} 从 long long 构造；更长整数用流读入。
\item 支持 \texttt{+ - * / \%}、复合赋值和比较。
\item \texttt{a.divmod(b)} 同时返回商和余数，避免分别执行除法与取模。
\end{itemize}}
\textbf{最小示例（放入 main）：}
\begin{lstlisting}
BigInt a,b;
istringstream in("-12345678901234567890 100");
in >> a >> b;
auto [q,r]=a.divmod(b);
assert(q*b+r==a);
ostringstream out; out << q << ' ' << r;
assert(out.str()=="-123456789012345678 -90");
\end{lstlisting}
\TemplNote{除法向零取整，余数与被除数同号；除数不能为零。这不是 FFT 高精度乘法，也不支持小数。}

\Needspace{6\baselineskip}\textbf{完整代码}
\begin{lstlisting}[firstnumber=1]
struct BigInt{
	static constexpr int base=1000000000;
	vector<int> a;
	int sign=1;
	BigInt(long long x=0){
		unsigned long long y=x;
		if(x<0)sign=-1,y=0-y;
		for(;y;y/=base)a.push_back(y%base);
	}
	void trim(){while(!a.empty()&&!a.back())a.pop_back();if(a.empty())sign=1;}
	explicit operator bool()const{return !a.empty();}
	int cmp_abs(const BigInt &b)const{
		if(a.size()!=b.a.size())return a.size()<b.a.size()?-1:1;
		for(int i=(int)a.size()-1;i>=0;i--)if(a[i]!=b.a[i])return a[i]<b.a[i]?-1:1;
		return 0;
	}
	int cmp(const BigInt &b)const{return sign!=b.sign?(sign<b.sign?-1:1):sign*cmp_abs(b);}
	bool operator <(const BigInt &b)const{return cmp(b)<0;}
	bool operator >(const BigInt &b)const{return cmp(b)>0;}
	bool operator <=(const BigInt &b)const{return cmp(b)<=0;}
	bool operator >=(const BigInt &b)const{return cmp(b)>=0;}
	bool operator ==(const BigInt &b)const{return sign==b.sign&&a==b.a;}
	bool operator !=(const BigInt &b)const{return !(*this==b);}
	BigInt operator -()const{BigInt b=*this;if(b)b.sign=-b.sign;return b;}
	BigInt operator +()const{return *this;}
	BigInt operator +(const BigInt &b)const{
		if(!b)return *this;
		if(!*this)return b;
		if(sign!=b.sign)return *this-(-b);
		BigInt c;c.sign=sign;c.a.resize(max(a.size(),b.a.size()));
		int carry=0;
		for(int i=0;i<(int)c.a.size();i++){
			long long x=(long long)carry+(i<(int)a.size()?a[i]:0)+(i<(int)b.a.size()?b.a[i]:0);
			c.a[i]=x%base;carry=x/base;
		}
		if(carry)c.a.push_back(carry);
		return c;
	}
	BigInt operator -(const BigInt &b)const{
		if(!b)return *this;
		if(!*this)return -b;
		if(sign!=b.sign)return *this+(-b);
		if(cmp_abs(b)<0)return -(b-*this);
		BigInt c=*this;int carry=0;
		for(int i=0;i<(int)a.size();i++){
			int x=a[i]-carry-(i<(int)b.a.size()?b.a[i]:0);
			carry=x<0;c.a[i]=x+(carry?base:0);
		}
		c.trim();return c;
	}
	BigInt mul(int x)const{
		assert(0<=x&&x<base);
		BigInt c;c.sign=sign;c.a.resize(a.size());long long carry=0;
		for(int i=0;i<(int)a.size();i++){
			long long y=1ll*a[i]*x+carry;c.a[i]=y%base;carry=y/base;
		}
		if(carry)c.a.push_back(carry);
		c.trim();return c;
	}
	BigInt operator *(const BigInt &b)const{
		BigInt c;if(!*this||!b)return c;
		c.sign=sign*b.sign;c.a.resize(a.size()+b.a.size());
		for(int i=0;i<(int)a.size();i++){
			long long carry=0;
			for(int j=0;j<(int)b.a.size();j++){
				long long x=c.a[i+j]+1ll*a[i]*b.a[j]+carry;
				c.a[i+j]=x%base;carry=x/base;
			}
			c.a[i+b.a.size()]=carry;
		}
		c.trim();return c;
	}
	// Division truncates towards zero; remainder has the dividend's sign.
	pair<BigInt,BigInt> divmod(const BigInt &b)const{
		assert(b);
		if(cmp_abs(b)<0)return {0,*this};
		if(b.a.size()==1){
			BigInt q=*this;long long rem=0;
			for(int i=(int)a.size()-1;i>=0;i--){
				long long x=rem*base+a[i];q.a[i]=x/b.a[0];rem=x%b.a[0];
			}
			q.sign=sign*b.sign;q.trim();return {q,BigInt(rem*sign)};
		}
		int norm=base/(b.a.back()+1);
		BigInt x=mul(norm),y=b.mul(norm),q,r;x.sign=y.sign=1;
		q.a.resize(x.a.size());
		for(int i=(int)x.a.size()-1;i>=0;i--){
			r.a.insert(r.a.begin(),x.a[i]);r.trim();
			int n=y.a.size();
			long long hi=r.a.size()>(size_t)n?r.a[n]:0;
			long long lo=r.a.size()>=(size_t)n?r.a[n-1]:0;
			int d=min((hi*base+lo)/y.a.back(),(long long)base-1);
			r=r-y.mul(d);
			while(r.sign<0)r=r+y,--d;
			q.a[i]=d;
		}
		long long carry=0;
		for(int i=(int)r.a.size()-1;i>=0;i--){
			long long x=r.a[i]+carry*base;r.a[i]=x/norm;carry=x%norm;
		}
		q.sign=sign*b.sign;r.sign=sign;q.trim();r.trim();return {q,r};
	}
	BigInt operator /(const BigInt &b)const{return divmod(b).first;}
	BigInt operator %(const BigInt &b)const{return divmod(b).second;}
	BigInt& operator +=(const BigInt &b){return *this=*this+b;}
	BigInt& operator -=(const BigInt &b){return *this=*this-b;}
	BigInt& operator *=(const BigInt &b){return *this=*this*b;}
	BigInt& operator /=(const BigInt &b){return *this=*this/b;}
	BigInt& operator %=(const BigInt &b){return *this=*this%b;}
	friend istream& operator >>(istream &in,BigInt &x){
		string s;if(!(in>>s))return in;
		int l=(s[0]=='+'||s[0]=='-');
		if(l==(int)s.size()){in.setstate(ios::failbit);return in;}
		for(int i=l;i<(int)s.size();i++)if(s[i]<'0'||s[i]>'9'){
			in.setstate(ios::failbit);return in;
		}
		BigInt y;y.sign=s[0]=='-'?-1:1;y.a.reserve((s.size()-l+8)/9);
		for(int r=s.size();r>l;r-=9){
			int v=0;for(int i=max(l,r-9);i<r;i++)v=v*10+s[i]-'0';
			y.a.push_back(v);
		}
		y.trim();x=move(y);return in;
	}
	friend ostream& operator <<(ostream &out,const BigInt &x){
		if(!x)return out<<'0';
		string s=x.sign<0?"-":"";s+=to_string(x.a.back());
		for(int i=(int)x.a.size()-2;i>=0;i--){
			string t=to_string(x.a[i]);s.append(9-t.size(),'0');s+=t;
		}
		return out<<s;
	}
};
\end{lstlisting}
\par\vspace{1em}

\Needspace{10\baselineskip}
\subsection{\texttt{basic/FastIO.cpp}}
\textbf{普通文件上的 mmap 快读快写}

\TemplUsage{适合 Linux 下大规模整数输入，输入必须是可映射的普通文件。}
\TemplApi{\begin{itemize}
\item 额外包含 \texttt{sys/stat.h}、\texttt{sys/mman.h}、\texttt{unistd.h}。
\item 重定向输入输出之后调用 \texttt{input.set(stdin)} 和 \texttt{output.set(stdout)}。
\item 整数使用 \texttt{input >> x}、\texttt{output << x}；\texttt{output.flush()} 主动刷新。
\end{itemize}}
\textbf{最小示例（放入 main）：}
\begin{lstlisting}
// Input file: 2 10 20
input.set(stdin); output.set(stdout);
int n; input >> n; ll sum=0;
for(int i=0;i<n;i++){ll x; input >> x; sum+=x;}
output << sum << '\n'; output.flush(); // 30
\end{lstlisting}
\TemplNote{不适用于交互题、终端或管道输入，也不能原生运行于 Windows。不读浮点和字符串。不要和 cin 混读同一个输入源。本手册此例因当前检查环境为 Windows，未执行 Linux 运行验证。}

\Needspace{6\baselineskip}\textbf{完整代码}
\begin{lstlisting}[firstnumber=1]
//
// template for FastIO (mmap)
// requires sys/stat.h, sys/mman.h, unistd.h
//
// usage:
//   input.set(stdin); output.set(stdout);  // after freopen
//   input >> x; output << x;  // integers only, Linux/WSL
//
// ---------------------------------------------------------------
// Fast mmap‑based reader / buffered writer (from the __yzlf code)
// ---------------------------------------------------------------
using u8  = unsigned char;
using u16 = unsigned short;
using u32 = unsigned;
using i64 = long long;
using u64 = unsigned long long;

constexpr std::size_t BUF_DEF_SIZE	   = 262144;
constexpr std::size_t BUF_FLUSH_THRESHOLD = 32;
constexpr u64 E16 = 10000000000000000ULL;
constexpr u64 E12 = 1000000000000ULL;
constexpr u64 E8  = 100000000ULL;
constexpr u64 E4  = 10000ULL;

struct _io_t {
	std::array<u8,1 << 15> t_i;		  // two‑digit decode table
	std::array<u32,10000> t_o;			// four ASCII digits packed in u32
	_io_t() : t_i{}, t_o{} {
		t_i.fill(u8(-1));
		for (int i = 0; i < 10; ++i)
			for (int j = 0; j < 10; ++j)
				t_i[0x3030 + 256 * j + i] = j + 10 * i;
		int idx = 0;
		for (int e0 = (48 << 0); e0 < (58 << 0); e0 += (1 << 0))
			for (int e1 = (48 << 8); e1 < (58 << 8); e1 += (1 << 8))
				for (int e2 = (48 << 16); e2 < (58 << 16); e2 += (1 << 16))
					for (int e3 = (48 << 24); e3 < (58 << 24); e3 += (1 << 24))
						t_o[idx++] = e0 ^ e1 ^ e2 ^ e3;
	}
	void get(char* s, u32 p) const { memcpy(s, &t_o[p], sizeof(u32)); }
};
_io_t _iot;

struct Reader {
	char* bg=nullptr;
	char* p=nullptr;
	char* ed=nullptr;
	std::size_t sz=0;
	Reader() = default;
	Reader(const Reader&)=delete;
	Reader& operator=(const Reader&)=delete;
	void set(FILE* f = stdin) {
		assert(f);
		if(bg)munmap(bg,sz);
		int fd=fileno(f);struct stat st;
		int res=fstat(fd,&st);assert(res==0&&S_ISREG(st.st_mode));
		sz=st.st_size;bg=p=ed=nullptr;
		if(!sz)return;
		bg=static_cast<char*>(mmap(nullptr,sz,PROT_READ,MAP_PRIVATE,fd,0));
		assert(bg!=MAP_FAILED);madvise(bg,sz,MADV_SEQUENTIAL);p=bg;ed=bg+sz;
	}

	~Reader(){if(bg)munmap(bg,sz);}

	inline void skip_space() { while(p!=ed&&*p<=' ')++p; }

	template <typename T>
	Reader& operator>>(T& x) {
		skip_space();assert(p!=ed);bool neg=false;
		if(*p=='+')++p;
		else if constexpr(std::is_signed_v<T>){neg=(*p=='-');p+=neg;}
		assert(p!=ed&&'0'<=*p&&*p<='9');
		using U=std::make_unsigned_t<T>;U v=*p++-'0';
		while(ed-p>=2){
			u16 z;memcpy(&z,p,sizeof(z));
			u32 y=z<_iot.t_i.size()?_iot.t_i[z]:255;
			if(y>99)break;
			v=v*100+y;p+=2;
		}
		if(p!=ed&&'0'<=*p&&*p<='9')v=v*10+(*p++-'0');
		x=neg?T(U(0)-v):T(v);return *this;
	}
};

struct Writer {
	std::vector<char> buf;
	char* bg=nullptr;char* p=nullptr;char* ed=nullptr;FILE* f=nullptr;
	Writer() = default;
	Writer(const Writer&)=delete;
	Writer& operator=(const Writer&)=delete;

	void set(FILE* fi = stdout, std::size_t sz = BUF_DEF_SIZE) {
		assert(fi&&sz>BUF_FLUSH_THRESHOLD);if(bg)flush();
		f = fi;buf.resize(sz);
		bg = buf.data();
		p = bg;
		ed = bg + sz - BUF_FLUSH_THRESHOLD;
	}
	~Writer() { if(bg)flush(); }
	inline void flush() { assert(bg&&f);fwrite(bg, 1, p - bg, f); p = bg; }
	inline void chk() { if (p > ed) flush(); }

	inline void put_block(u32 x) { _iot.get(p, x); p += 4; }

	inline void put2(u32 x) {
		if (x > 9) { _iot.get(p, x * 100); p += 2; }
		else *p++ = char('0' + x);
	}
	inline void put4(u32 x) {
		if (x > 99) { if (x > 999) put_block(x); else { _iot.get(p, x * 10); p += 3; } }
		else put2(x);
	}
	inline void write_u32(u32 x) {
		if (x >= E8) { put2(x / E8); x %= E8; put_block(x / E4); put_block(x % E4); }
		else if (x >= E4) { put4(x / E4); put_block(x % E4); }
		else put4(x);
		chk();
	}
	inline void write_u64(u64 x) {
		if (x >= E8) {
			u64 q0 = x / E8, r0 = x % E8;
			if (x >= E16) { u64 q1 = q0 / E8, r1 = q0 % E8; put4(q1); put_block(r1 / E4); put_block(r1 % E4); }
			else if (x >= E12) { put4(q0 / E4); put_block(q0 % E4); }
			else put4(q0);
			put_block(r0 / E4); put_block(r0 % E4);
		} else if (x >= E4) { put4(x / E4); put_block(x % E4); }
		else put4(static_cast<u32>(x));
		chk();
	}

	Writer& operator<<(u32 x) { write_u32(x); return *this; }
	Writer& operator<<(u64 x) { write_u64(x); return *this; }
	Writer& operator<<(int x) { if (x < 0) { *p++ = '-'; write_u32(-static_cast<u32>(x)); } else write_u32(static_cast<u32>(x)); return *this; }
	Writer& operator<<(long long x) { if (x < 0) { *p++ = '-'; write_u64(-static_cast<u64>(x)); } else write_u64(static_cast<u64>(x)); return *this; }
	Writer& operator<<(char c) { *p++ = c; chk(); return *this; }
};
inline Reader input;
inline Writer output;
\end{lstlisting}
\par\vspace{1em}

\Needspace{10\baselineskip}
\subsection{\texttt{basic/Coordinate Compression.cpp}}
\textbf{离散化并原地回写}

\TemplUsage{先收集所有会用到的值，再排序去重，把已登记变量改写成 1..k 的排名。时间 O(n log n)。}
\TemplApi{\begin{itemize}
\item \texttt{compress<T> c; c.set()} 清空。\texttt{c.pb(x)} 要求可修改左值，会保存 x 的地址。
\item \texttt{c.build()} 原地改写 x；\texttt{c.id(original)} 找已登记原值的排名。
\item \texttt{c[i]} 由排名恢复原值；\texttt{range(l,r)} 返回落在原值闭区间中的排名范围，左端大于右端表示空。
\end{itemize}}
\textbf{最小示例（放入 main）：}
\begin{lstlisting}
compress<int> c; c.set();
int a=30,b=10,d=30;
c.pb(a); c.pb(b); c.pb(d); c.build();
assert(a==2 && b==1 && d==2 && c[2]==30);
assert(c.id(30)==2);
\end{lstlisting}
\TemplNote{pb 后、build 前变量必须存活且地址稳定：不要登记即将离开作用域的局部变量，也不要让所属 vector 扩容。\texttt{id} 参数是原值，不是已改写后的排名。range 当前端点参数为 int。}

\Needspace{6\baselineskip}\textbf{完整代码}
\begin{lstlisting}[firstnumber=1]
//
// template for Coordinate Compression
// version 2.0 (Last Update Jul 3rd, 2026)
//
// usage:
//   compress<int> cc; cc.set(); cc.pb(x); cc.build();
//   cc.id(x); cc[i]; cc.range(l,r);  // call set() each problem
//
template<typename T>
struct compress{
vector<T> I;
vector<T*> II;
// make compress values to 1 ~ size
void set(){I.clear();II.clear();}
void pb(T &x){I.push_back(x);II.push_back(&x);}
void pb(const T &x){I.push_back(x);}
int id(T x)const{
	auto it=lower_bound(I.begin(),I.end(),x);
	assert(it!=I.end()&&*it==x);
	return it-I.begin()+1;
}
void build(){
	sort(I.begin(),I.end());
	I.erase(unique(I.begin(),I.end()),I.end());
	vector<int> a; a.reserve(II.size());
	for(auto p:II)a.push_back(id(*p));
	for(int i=0;i<(int)II.size();i++)*II[i]=a[i];
	II.clear();
}
int size()const{return I.size();}
array<int,2> range(const T &l,const T &r)const{
	int a=lower_bound(I.begin(),I.end(),l)-I.begin();
	int b=upper_bound(I.begin(),I.end(),r)-I.begin();
	return {a+1,b};
}
T operator [](int x)const{
	assert(1<=x&&x<=I.size());
	return I[x-1];
}
};
// !!!!!!!!!! before each use, use cc.set() to clear !!!!!!!
// !!!!!!!!!! 如果元素被插入，cc.build() 会自动重编号 !!!!!!!!!!!!!
\end{lstlisting}
\par\vspace{1em}

\Needspace{10\baselineskip}
\subsection{\texttt{basic/ModInt.cpp}}
\textbf{固定模数及组合数}

\TemplUsage{默认模数 998244353。模数写在源码中的 \texttt{MOD}；阶乘表按需扩容。}
\TemplApi{\begin{itemize}
\item \texttt{mint} 支持四则、\texttt{pow(k)}、\texttt{inv()}；除数必须可逆。
\item \texttt{init(n)} 保证 fac、ifac、inv 覆盖到 n，组合数调用会自动预处理。
\item \texttt{C(k,n)} 和 \texttt{binom(n,k)} 都表示从 n 个中选 k 个；这里 C 的参数顺序与常见记法相反。
\end{itemize}}
\textbf{最小示例（放入 main）：}
\begin{lstlisting}
init(10);
assert(C(2,5).x==10 && binom(5,2).x==10);
mint a=3; assert((a*a.inv()).x==1);
assert(a.pow(4).x==81);
\end{lstlisting}
\TemplNote{阶乘方案用于质数模数且 n<MOD；比较时可读 \texttt{.x}，此版未定义 mint 的 ==。输入运算符直接读 \texttt{.x}，读负数或未归一化数据时先读 long long 再构造 mint。避免 \texttt{pow(LLONG\_MIN)}。多项式文件已内置本类，不要重复粘贴。}

\Needspace{6\baselineskip}\textbf{完整代码}
\begin{lstlisting}[firstnumber=1]
template <unsigned _M> struct ModInt{
static constexpr unsigned MOD=_M;
static_assert(1<MOD&&MOD<=INT_MAX);
unsigned x;
constexpr ModInt():x(0){}
constexpr ModInt(unsigned y):x(y%MOD){}
constexpr ModInt(int y):x((y%=static_cast<int>(MOD))<0?y+MOD:y){}
constexpr ModInt(unsigned long long y):x(y%MOD){}
constexpr ModInt(long long y):x((y%=static_cast<long long>(MOD))<0?y+MOD:y){}
ModInt& operator +=(const ModInt &a){x=(((x+=a.x)>=MOD)?x-MOD:x);return *this;}
ModInt& operator -=(const ModInt &a){x=(((x-=a.x)>=MOD)?x+MOD:x);return *this;}
ModInt& operator *=(const ModInt &a){x=static_cast<unsigned long long>(x)*a.x%MOD;return *this;}
ModInt& operator /=(const ModInt &a){return (*this)*=a.inv();}
bool operator ==(const ModInt &a)const{return x==a.x;}
bool operator !=(const ModInt &a)const{return x!=a.x;}
explicit operator bool() const { return x != 0; }
bool operator!() const { return x == 0; }
ModInt inv() const {
	unsigned a=MOD,b=x;ll y=0,z=1;
	while(b){
		const unsigned q=a/b,c=a-q*b;
		a=b,b=c;
		const ll w=y-(ll)q*z;
		y=z,z=w;
	}
	assert(a==1);return ModInt(y);
}
ModInt pow(long long k) const {
	unsigned long long e=k;
	ModInt a=*this,b=1;
	if(k<0){a=a.inv();e=0-e;}
	for(;e;e>>=1){if(e&1)b*=a;a*=a;}
	return b;
}
ModInt& operator++() {*this+=1;return *this;}
ModInt& operator--() {*this-=1;return *this;}
ModInt operator--(int){ModInt t=*this;--*this;return t;}
ModInt operator++(int){ModInt t=*this;++*this;return t;}
ModInt operator +()const{return *this;}
ModInt operator -()const{ModInt a;a.x=(x?MOD-x:0u);return a;}
ModInt operator +(const ModInt &a)const{return ModInt(*this)+=a;}
ModInt operator -(const ModInt &a)const{return ModInt(*this)-=a;}
ModInt operator *(const ModInt &a)const{return ModInt(*this)*=a;}
ModInt operator /(const ModInt &a)const{return ModInt(*this)/=a;}

template<typename T> ModInt friend operator +(T a,const ModInt &b){return ModInt(a)+=b;}
template<typename T> ModInt friend operator -(T a,const ModInt &b){return ModInt(a)-=b;}
template<typename T> ModInt friend operator *(T a,const ModInt &b){return ModInt(a)*=b;}
template<typename T> ModInt friend operator /(T a,const ModInt &b){return ModInt(a)/=b;}
friend istream& operator >> (istream &i,ModInt &x){ll y;if(i>>y)x=ModInt(y);return i;}
friend ostream& operator << (ostream &o,const ModInt &x){o<<x.x;return o;}
};
// Basic function, fac, ifac, binom
// init(n) prepares fac / ifac / inv through n
const int MOD=998244353;
using mint=ModInt<MOD>;
vector<mint> fac{1},ifac{1},inv{0,1};
void init(int n=0){
	assert(0<=n&&n<MOD);
	int m=fac.size();if(n<m)return;
	n=max(n,(int)min(2ll*m,(ll)MOD-1));
	fac.resize(n+1);ifac.resize(n+1);inv.resize(max(n+1,2));
	for(int i=m;i<=n;i++){
		fac[i]=fac[i-1]*i;
		if(i>1)inv[i]=-(MOD/i)*inv[MOD%i];
		ifac[i]=ifac[i-1]*inv[i];
	}
}
inline mint C(int x,int y){
	// choose x from y
	if(x<0||y<x)return 0;
	init(y);return fac[y]*ifac[x]*ifac[y-x];
}
inline mint binom(int y,int x){return C(x,y);}

inline mint Apple_in_Box(int n,int m){
	// n apples into m boxes, each box must be non-empty
	return C(m-1,n-1);
}
inline mint Apple_in_Box(int n,int m,int k){
	// n apples into m boxes, 
	// the first k boxes must be non-empty (others can be empty)
	assert(0<=k&&k<=m);
	return Apple_in_Box(n+(m-k),m);
}
\end{lstlisting}
\par\vspace{1em}

\Needspace{10\baselineskip}
\subsection{\texttt{basic/ModInt-nonstaticMod.cpp}}
\textbf{运行时模数及组合数}

\TemplUsage{同一份程序需要选定不同模数时使用。设置模数后再创建该轮使用的值。}
\TemplApi{\begin{itemize}
\item \texttt{mint::setM(p)} 修改模数。
\item \texttt{init\_fact(n)} 重建 0..n 阶乘、逆阶乘，要求 n! 可逆。
\item \texttt{C(k,n)} / \texttt{binom(n,k)} 查询前必须已预处理到 n。
\end{itemize}}
\textbf{最小示例（放入 main）：}
\begin{lstlisting}
mint::setM(101); init_fact(10);
assert(C(2,5).x==10);
mint a=100; assert((a+2).x==1);
\end{lstlisting}
\TemplNote{修改模数不会转换已有 mint 变量；丢弃旧值并重建阶乘。与固定模数版、各多项式版本的 mint / fac 等全局名字冲突。}

\Needspace{6\baselineskip}\textbf{完整代码}
\begin{lstlisting}[firstnumber=1]
//
// template for ModInt (Mod isn't assign)
// version 1.0 (Last Update Jun 16th, 2026)
//
// usage:
//   mint::setM(mod); init_fact(n); mint x; cin >> x;
//
struct mint {
	static unsigned M;
	static unsigned long long NEG_INV_M;
	static void setM(unsigned m) {
		assert(1<m&&m<=INT_MAX);M = m;
		NEG_INV_M = -1ULL / M;
	}
	unsigned x;
	mint() : x(0U) {}
	mint(unsigned x_) : x(x_ % M) {}
	mint(unsigned long long x_) : x(x_ % M) {}
	mint(int x_) : x(((x_ %= static_cast<int>(M)) < 0) ? (x_ + static_cast<int>(M)) : x_) {}
	mint(long long x_) : x(((x_ %= static_cast<long long>(M)) < 0) ? (x_ + static_cast<long long>(M)) : x_) {}
	mint& operator+=(const mint &a) {
		x = ((x += a.x) >= M) ? (x - M) : x;
		return *this;
	}
	mint& operator-=(const mint &a) {
		x = ((x -= a.x) >= M) ? (x + M) : x;
		return *this;
	}
	mint& operator*=(const mint &a) {
		const unsigned long long y = static_cast<unsigned long long>(x) * a.x;
		const unsigned long long q = static_cast<unsigned long long>(
			(static_cast<unsigned __int128>(NEG_INV_M) * y) >> 64
		);
		const unsigned long long r = y - M * q;
		x = r - M * (r >= M);
		return *this;
	}
	mint& operator/=(const mint &a) { return (*this *= a.inv()); }
	mint pow(long long e) const {
		unsigned long long k=e;
		mint a=*this,b=1;
		if(e<0){a=a.inv();k=0-k;}
		for(;k;k>>=1){if(k&1)b*=a;a*=a;}
		return b;
	}
	mint inv() const {
		unsigned a = M, b = x;
		ll y = 0, z = 1;
		for (; b; ) {
			const unsigned q = a / b;
			const unsigned c = a - q * b;
			a = b; b = c;
			const ll w = y - (ll)q * z;
			y = z; z = w;
		}
		assert(a == 1U);
		return mint(y);
	}
	mint operator+() const { return *this; }
	mint operator-() const { mint a; a.x = x ? (M - x) : 0U; return a;}
	mint operator+(const mint &a) const { return (mint(*this) += a); }
	mint operator-(const mint &a) const { return (mint(*this) -= a); }
	mint operator*(const mint &a) const { return (mint(*this) *= a); }
	mint operator/(const mint &a) const { return (mint(*this) /= a); }
	mint& operator++() { *this += 1; return *this; }
	mint& operator--() { *this -= 1; return *this; }
	bool operator!() const { return x == 0; }
	explicit operator bool() const { return x != 0; }
	bool operator==(const mint &a) const { return (x == a.x); }
	bool operator!=(const mint &a) const { return (x != a.x); }
	template<class T> friend mint operator+(T a, const mint &b) { return (mint(a) += b); }
	template<class T> friend mint operator-(T a, const mint &b) { return (mint(a) -= b); }
	template<class T> friend mint operator*(T a, const mint &b) { return (mint(a) *= b); }
	template<class T> friend mint operator/(T a, const mint &b) { return (mint(a) /= b); }
	friend istream& operator>>(istream &i, mint &a) { ll x;if(i>>x)a=mint(x);return i; }
	friend ostream& operator<<(ostream &o, const mint &a) { return o << a.x; }
};
unsigned mint::M = 998244353;
unsigned long long mint::NEG_INV_M = -1ULL / mint::M;

// !!!!!!!!!!!Use mint::setM(<new mod>)!!!!!!!!!!

vector<mint> fac,ifac,inv;
unsigned fact_mod=0;
void init_fact(int n=0){
	assert(0<=n&&(unsigned)n<mint::M);
	fac.assign(n+1,1);ifac.resize(n+1);inv.assign(n+1,0);
	for(int i=1;i<=n;i++)fac[i]=fac[i-1]*i;
	ifac[n]=fac[n].inv();
	for(int i=n;i>=1;i--)ifac[i-1]=ifac[i]*i;
	for(int i=1;i<=n;i++)inv[i]=ifac[i]*fac[i-1];
	fact_mod=mint::M;
}
inline mint C(int x,int y){
	if(x<0||y<x)return 0;
	assert(fact_mod==mint::M&&y<(int)fac.size());
	return fac[y]*ifac[x]*ifac[y-x];
}
inline mint binom(int y,int x){return C(x,y);}
\end{lstlisting}
\par\vspace{1em}

\Needspace{10\baselineskip}
\subsection{\texttt{basic/roll back saver (int).cpp}}
\textbf{记录并撤销变量修改}

\TemplUsage{保存变量的地址和旧值，用版本号撤销一段修改。每次保存和撤销一次记录都是 O(1)。}
\TemplApi{\begin{itemize}
\item 修改前先 \texttt{saver.save(x)}。
\item \texttt{saver.version()} 返回当前记录数，\texttt{saver.roll\_back(t)} 逆序恢复到该记录数。
\item \texttt{saver.set()} 仅丢弃历史，不恢复变量。
\end{itemize}}
\textbf{最小示例（放入 main）：}
\begin{lstlisting}
saver.set(); int x=7;
auto t=saver.version(); saver.save(x); x=20;
saver.roll_back(t); assert(x==7);
\end{lstlisting}
\TemplNote{保存期间变量地址必须保持有效，不能让包含它的 vector 扩容。每个版本号只能用于当前历史；两种 saver 不能同时粘贴。}

\Needspace{6\baselineskip}\textbf{完整代码}
\begin{lstlisting}[firstnumber=1]
//
// template for roll back saver
//
// usage:
//   saver.set(); saver.save(x); int t=saver.version(); saver.roll_back(t);
//   int only; pair with DSU (with undo)
//   saved elements must stay at the same address until rollback
//
struct rollback_saver{
vector<pair<int*,int>> buf;
void set(){buf.clear();}
void save(int &x){buf.push_back(make_pair(&x,x));}
int version(){return buf.size();}
void roll_back(int t){
	assert(0<=t&&t<=version());
	while(buf.size()>t){
		auto [a,b]=buf.back();buf.pop_back();
		*a=b;
	}
}
};
rollback_saver saver;
\end{lstlisting}
\par\vspace{1em}

\Needspace{10\baselineskip}
\subsection{\texttt{basic/roll back saver (int or ll).cpp}}
\textbf{记录并撤销变量修改}

\TemplUsage{保存变量的地址和旧值，用版本号撤销一段修改。每次保存和撤销一次记录都是 O(1)。}
\TemplApi{\begin{itemize}
\item 修改前先 \texttt{saver.save(x)}。
\item \texttt{saver.version()} 返回当前记录数，\texttt{saver.roll\_back(t)} 逆序恢复到该记录数。
\item \texttt{saver.set()} 仅丢弃历史，不恢复变量。
\end{itemize}}
\textbf{最小示例（放入 main）：}
\begin{lstlisting}
saver.set(); ll x=7;
auto t=saver.version(); saver.save(x); x=20;
saver.roll_back(t); assert(x==7);
\end{lstlisting}
\TemplNote{保存期间变量地址必须保持有效，不能让包含它的 vector 扩容。每个版本号只能用于当前历史；两种 saver 不能同时粘贴。}

\Needspace{6\baselineskip}\textbf{完整代码}
\begin{lstlisting}[firstnumber=1]
//
// template for roll back saver
//
// usage:
//   saver.set(); saver.save(x); saver.save(y); saver.roll_back(t);
//   supports int and ll; do not paste with (int) saver
//   saved elements must stay at the same address until rollback
//
using ll=long long;
struct rollback_saver{
vector<pair<int*,int>> buf;
vector<pair<ll*,ll>> buf1;
int saved_ele=0;
vector<char> R=vector<char>(1);
void set(){saved_ele=0;buf.clear();buf1.clear();R.assign(1,false);}
void save(int &x){buf.push_back(make_pair(&x,x));R.push_back(false);++saved_ele;}
void save(ll &x){buf1.push_back(make_pair(&x,x));R.push_back(true);++saved_ele;}
int version(){return saved_ele;}
void roll_back(int t){
	assert(0<=t&&t<=version());
	while(saved_ele>t){
		if(R[saved_ele]){
			auto [a,b]=buf1.back();buf1.pop_back();
			*a=b;
		}else{
			auto [a,b]=buf.back();buf.pop_back();
			*a=b;
		}
		saved_ele--;R.pop_back();
	}
}
};
rollback_saver saver;
\end{lstlisting}
\par\vspace{1em}

\Needspace{10\baselineskip}
\subsection{\texttt{basic/DSU (with undo).cpp}}
\textbf{可撤销并查集}

\TemplUsage{适合线段树分治或搜索中加入连边后回退。按大小合并、不做路径压缩，单次查找 O(log n)。}
\TemplApi{\begin{itemize}
\item 先粘贴一种 saver，再粘贴本文件。
\item \texttt{unionfind d; d.setN(n)} 初始化；\texttt{find/merge} 与普通版类似。
\item 通过全局 \texttt{saver.version/saver.roll\_back} 撤销 merge，而不是调用对象成员。
\end{itemize}}
\textbf{额外依赖（按此顺序粘贴）：}\texttt{basic/roll back saver (int).cpp}。

\textbf{最小示例（放入 main）：}
\begin{lstlisting}
saver.set(); unionfind d; d.setN(3);
auto t=saver.version(); d.merge(1,2);
assert(d.find(1)==d.find(2));
saver.roll_back(t); assert(d.find(1)!=d.find(2));
\end{lstlisting}
\TemplNote{setN 前先清空旧历史，不能在旧历史仍引用数组时重新分配。不要自行加入路径压缩。}

\Needspace{6\baselineskip}\textbf{完整代码}
\begin{lstlisting}[firstnumber=1]
//
// template for unionfind (accept undo, O(nlogn))
// saver.version 1.0 (Last Update Jun 28th, 2026)
//
// usage:
//   Paste roll back saver (int).cpp first.
//   unionfind uf; uf.setN(n); int t=saver.version(); uf.merge(u,v); saver.roll_back(t);
//
struct unionfind{
int n=0;
vector<int> fa,siz;
inline void setN(int _n){
	// saved addresses must not survive setN()
	assert(saver.version()==0);
	assert(0<=_n&&_n<INT_MAX);n=_n;
	fa.resize(n+1);iota(fa.begin(),fa.end(),0);siz.assign(n+1,1);siz[0]=0;
}
inline int find(int x){assert(1<=x&&x<=n);while(x^fa[x])x=fa[x];return x;}
inline int merge(int x,int y){
	assert(1<=min(x,y)&&max(x,y)<=n);
	x=find(x),y=find(y);
	if(x==y) return 0;
	if(siz[x]<siz[y]) swap(x,y);
	saver.save(siz[x]);saver.save(fa[y]);
	siz[x]+=siz[y];fa[y]=x;
	return x;
}
};
\end{lstlisting}
\par\vspace{1em}

\clearpage
\section{数据结构}
\Needspace{10\baselineskip}
\subsection{\texttt{ds/BIT.cpp}}
\textbf{自定义下标的树状数组}

\TemplUsage{同时提供单点加的 BIT 和区间加的双 BIT。初始化 O(n)，单次操作 O(log n)，空间 O(n)。}
\TemplApi{两种结构都用 \texttt{setRange(L,R)} 设置闭区间，\texttt{setN(n)} 等价于 \texttt{setRange(1,n)}。BIT 的 \texttt{add(i,v)} 单点加，\texttt{ask(r)} 返回 [L,r] 的和，\texttt{ask(l,r)} 返回区间和。rars 的 \texttt{add(l,r,v)} 区间加，\texttt{ask(l,r)} 查区间和。}
\textbf{最小示例（放入 main）：}
\begin{lstlisting}
BIT b; b.setRange(-2,2); b.add(0,7);
assert(b.ask(0)==7 && b.ask(-1,1)==7);
rars t; t.setRange(-2,2); t.add(-1,1,3);
assert(t.ask(-2,0)==6);
b.setN(3); assert(b.ask(3)==0);
\end{lstlisting}
\TemplNote{下标 0 或负数只要在设定区间内就合法；内部会转成正下标。ask(L-1) 返回零，add(R+1,v) 是差分哨兵，不计入查询。极端坐标的 L-1/R+1 用 ll 计算。rars 的加权差分使用相对位置，不依赖坐标绝对值；值、加权积与答案须在 ll 范围内。set() 清空；重新 setRange/setN 丢弃旧值。文件末尾保留全局 rars 对象 T。}

\Needspace{6\baselineskip}\textbf{完整代码}
\begin{lstlisting}[firstnumber=1]
////////////////////////////////////////////////////////////////
//
// template for BIT
//
// usage:
//   BIT: setRange(l,r); add(i,v); ask(r); ask(l,r).
//   setN(n) = setRange(1,n); add(vr+1,v) is a difference sentinel.
//   rars: setRange(l,r); add(l,r,v); ask(l,r).
//
////////////////////////////////////////////////////////////////
struct BIT{
private:
struct info{
	ll x=0;
	info(){}
	info(ll _x):x(_x){}
	info& operator +=(const info &y){x+=y.x;return *this;}
	info operator +(const info &y)const{return info(*this)+=y;}
};
int vl=1,vr=0,n=0;vector<info> a;
info query(int x)const{
	info ans;for(;x;x-=x&-x)ans+=a[x];return ans;
}
public:
void set(){vl=1;vr=n=0;a.clear();}
void setRange(int l,int r){
	ll m=(ll)r-l+1;assert(0<=m&&m<INT_MAX-1);
	vl=l;vr=r;n=m;a.assign(n+2,info());
}
void setN(int n){setRange(1,n);}
void add(ll x,ll v){
	assert(!a.empty()&&vl<=x&&x<=(ll)vr+1);
	for(ll i=x-vl+1;i<=n+1;i+=i&-i)a[i]+=info(v);
}
ll ask(ll x)const{assert((ll)vl-1<=x&&x<=vr);return query(x-vl+1).x;}
ll ask(int l,int r)const{assert(vl<=l&&l<=r&&r<=vr);return ask(r)-ask((ll)l-1);}
};
struct rars{
int vl=1,vr=0;BIT T0,T1;
void set(){vl=1;vr=0;T0.set();T1.set();}
void setRange(int l,int r){vl=l;vr=r;T0.setRange(l,r);T1.setRange(l,r);}
void setN(int n){setRange(1,n);}
void add(int l,int r,ll x){
	assert(vl<=l&&l<=r&&r<=vr);
	ll a=(ll)l-vl+1,b=(ll)r-vl+2;
	T0.add(l,x);T0.add((ll)r+1,-x);
	T1.add(l,a*x);T1.add((ll)r+1,-b*x);
}
ll query(ll x)const{return (x-vl+2)*T0.ask(x)-T1.ask(x);}
ll ask(int l,int r)const{assert(vl<=l&&l<=r&&r<=vr);return query(r)-query((ll)l-1);}
}T;
// end for ds/BIT.cpp
/////////////////////////
\end{lstlisting}
\Needspace{10\baselineskip}
\subsection{\texttt{ds/BIT/Sum.cpp}}
\textbf{树状数组：和特化}

\TemplUsage{初始化 O(n)，单点修改和前缀查询 O(log n)，空间 O(n)。}
\TemplApi{\texttt{Sum\_BIT t; t.setRange(L,R)}；\texttt{setN(n)} 调用 \texttt{setRange(1,n)}。\texttt{add(i,v)} 单点加；\texttt{ask(r)} 查询 [L,r] 的和；\texttt{ask(l,r)} 查询区间和。}
\textbf{最小示例（放入 main）：}
\begin{lstlisting}
Sum_BIT t; t.setRange(-2,2); t.add(0,5); t.add(1,-3);
assert(t.ask(0)==5 && t.ask(0,1)==2);
\end{lstlisting}
\TemplNote{初始值全零；允许 add(R+1,v) 差分哨兵及 ask(L-1) 空前缀。支持负坐标，查询不会修改树。 三个类名不同，可组合粘贴；set() 清空。}

\Needspace{6\baselineskip}\textbf{完整代码}
\begin{lstlisting}[firstnumber=1]
////////////////////////////////////////////////////////////////
//
// template for BIT (sum)
//
// usage:
//   Sum_BIT t; t.setRange(l,r); // setN(n) = setRange(1,n)
//   add(i,v); ask(r); ask(l,r); initial values are zero.
//   add(vr+1,v) is allowed as a difference sentinel.
//
////////////////////////////////////////////////////////////////
struct Sum_BIT{
private:
int vl=1,vr=0,n=0;vector<ll> a;
public:
void set(){vl=1;vr=n=0;a.clear();}
void setRange(int l,int r){
	ll m=(ll)r-l+1;assert(0<=m&&m<INT_MAX-1);
	vl=l;vr=r;n=m;a.assign(n+2,0);
}
void setN(int n){setRange(1,n);}
void add(ll x,ll v){
	assert(!a.empty()&&vl<=x&&x<=(ll)vr+1);
	for(ll i=x-vl+1;i<=n+1;i+=i&-i)a[i]+=v;
}
ll ask(ll x)const{
	assert((ll)vl-1<=x&&x<=vr);ll ans=0;
	for(int i=x-vl+1;i;i-=i&-i)ans+=a[i];
	return ans;
}
ll ask(int l,int r)const{assert(vl<=l&&l<=r&&r<=vr);return ask(r)-ask((ll)l-1);}
};
// end for ds/BIT/Sum.cpp
/////////////////////////
\end{lstlisting}
\Needspace{10\baselineskip}
\subsection{\texttt{ds/BIT/Max.cpp}}
\textbf{树状数组：最大值特化}

\TemplUsage{初始化 O(n)，单点修改和前缀查询 O(log n)，空间 O(n)。}
\TemplApi{\texttt{Max\_BIT t; t.setRange(L,R)}；\texttt{setN(n)} 调用 \texttt{setRange(1,n)}。\texttt{upd(i,v)} 把第 i 项取最大值，并非赋值；\texttt{ask(r)} 查询 [L,r] 的最大值。}
\textbf{最小示例（放入 main）：}
\begin{lstlisting}
Max_BIT t; t.setRange(-2,2); t.upd(0,5); t.upd(0,3);
assert(t.ask(1)==5);
\end{lstlisting}
\TemplNote{初值和空前缀返回 LLONG\_MIN。只支持单点 chmax；不能把已有大值减小。没有 ask(l,r)，最值不能像和一样用两个前缀相减；任意点赋值或一般区间最值请用线段树。 三个类名不同，可组合粘贴；set() 清空。}

\Needspace{6\baselineskip}\textbf{完整代码}
\begin{lstlisting}[firstnumber=1]
////////////////////////////////////////////////////////////////
//
// template for BIT (max)
//
// usage:
//   Max_BIT t; t.setRange(l,r); // setN(n) = setRange(1,n)
//   upd(i,v) is chmax, not assignment; ask(r) queries [vl,r].
//   Initial / empty-prefix value: LLONG_MIN. No inverse for interval queries.
//
////////////////////////////////////////////////////////////////
struct Max_BIT{
private:
int vl=1,vr=0,n=0;vector<ll> a;
public:
void set(){vl=1;vr=n=0;a.clear();}
void setRange(int l,int r){
	ll m=(ll)r-l+1;assert(0<=m&&m<INT_MAX-1);
	vl=l;vr=r;n=m;a.assign(n+2,LLONG_MIN);
}
void setN(int n){setRange(1,n);}
void upd(ll x,ll v){
	assert(vl<=x&&x<=vr);
	for(ll i=x-vl+1;i<=n;i+=i&-i)a[i]=max(a[i],v);
}
ll ask(ll x)const{
	assert((ll)vl-1<=x&&x<=vr);ll ans=LLONG_MIN;
	for(int i=x-vl+1;i;i-=i&-i)ans=max(ans,a[i]);
	return ans;
}

};
// end for ds/BIT/Max.cpp
/////////////////////////
\end{lstlisting}
\Needspace{10\baselineskip}
\subsection{\texttt{ds/BIT/Min.cpp}}
\textbf{树状数组：最小值特化}

\TemplUsage{初始化 O(n)，单点修改和前缀查询 O(log n)，空间 O(n)。}
\TemplApi{\texttt{Min\_BIT t; t.setRange(L,R)}；\texttt{setN(n)} 调用 \texttt{setRange(1,n)}。\texttt{upd(i,v)} 把第 i 项取最小值，并非赋值；\texttt{ask(r)} 查询 [L,r] 的最小值。}
\textbf{最小示例（放入 main）：}
\begin{lstlisting}
Min_BIT t; t.setRange(-2,2); t.upd(0,5); t.upd(0,3);
assert(t.ask(1)==3);
\end{lstlisting}
\TemplNote{初值和空前缀返回 LLONG\_MAX。只支持单点 chmin；不能把已有小值增大。没有 ask(l,r)，最值不能像和一样用两个前缀相减；任意点赋值或一般区间最值请用线段树。 三个类名不同，可组合粘贴；set() 清空。}

\Needspace{6\baselineskip}\textbf{完整代码}
\begin{lstlisting}[firstnumber=1]
////////////////////////////////////////////////////////////////
//
// template for BIT (min)
//
// usage:
//   Min_BIT t; t.setRange(l,r); // setN(n) = setRange(1,n)
//   upd(i,v) is chmin, not assignment; ask(r) queries [vl,r].
//   Initial / empty-prefix value: LLONG_MAX. No inverse for interval queries.
//
////////////////////////////////////////////////////////////////
struct Min_BIT{
private:
int vl=1,vr=0,n=0;vector<ll> a;
public:
void set(){vl=1;vr=n=0;a.clear();}
void setRange(int l,int r){
	ll m=(ll)r-l+1;assert(0<=m&&m<INT_MAX-1);
	vl=l;vr=r;n=m;a.assign(n+2,LLONG_MAX);
}
void setN(int n){setRange(1,n);}
void upd(ll x,ll v){
	assert(vl<=x&&x<=vr);
	for(ll i=x-vl+1;i<=n;i+=i&-i)a[i]=min(a[i],v);
}
ll ask(ll x)const{
	assert((ll)vl-1<=x&&x<=vr);ll ans=LLONG_MAX;
	for(int i=x-vl+1;i;i-=i&-i)ans=min(ans,a[i]);
	return ans;
}

};
// end for ds/BIT/Min.cpp
/////////////////////////
\end{lstlisting}

\par\vspace{1em}

\Needspace{10\baselineskip}
\subsection{\texttt{ds/Segment Tree/AddSum.cpp}}
\textbf{区间加、区间和}

\TemplUsage{默认闭区间下标 1..n、初值全零，初始化 O(n)，单次修改查询 O(log n)。}
\TemplApi{\begin{itemize}
\item \texttt{setN(n)} 或 \texttt{set(l,r)} 指定有效闭区间。
\item \texttt{upd(i,v)} 是赋值，\texttt{add(i,v)} 是加法；\texttt{add(l,r,v)} 区间加。
\item \texttt{ask(l,r)} 查询，\texttt{ask(i)} 查单点。
\item \texttt{find\_first(k)} 从左累计达到 k 的位置；\texttt{find\_last(k)} 从右累计达到 k 的位置。
\end{itemize}}
\textbf{最小示例（放入 main）：}
\begin{lstlisting}
AddSum_segt t; t.setN(3);
t.upd(2,5); t.add(1,2,2);
assert(t.ask(1,3)==9);
\end{lstlisting}
\TemplNote{这是区间加的特化实现；要更换信息或标记类型，使用外面的两份通用模板。set/setN 直接初始化全零节点，不递归建树。区间长度由递归参数计算，不存入节点；ask 和二分携带祖先标记，不下传、不改树。使用二分查询时所有点值必须非负，且 1<=k<=总和；k 参数为 ll。三个类名不同，可以组合粘贴。}

\Needspace{6\baselineskip}\textbf{完整代码}
\begin{lstlisting}[firstnumber=1]
////////////////////////////////////////////////////////////////
//
// template for Segment Tree — 区间加 / 区间和 / 前缀和二分
// version 1.1 (Last Update Jun 27th, 2026)
//
// usage:
//   AddSum_segt st; st.setN(n);  // or set(l,r), 1-indexed
//   upd/add point; add(l,r,v) range; ask; find_first/last
//
////////////////////////////////////////////////////////////////
struct AddSum_segt{
private:
struct node{ll s=0,a=0;};
int vl=1,vr=0;
vector<node> tr;
inline void pushup(int id){tr[id].s=tr[id<<1].s+tr[id<<1|1].s;}
inline void push(int id,ll x,int len){
	tr[id].s+=x*len;
	if(len>1)tr[id].a+=x;
}
inline void pushdown(int id,int l,int r){
	ll x=tr[id].a;if(!x)return;
	int mid=l+((ll)r-l)/2;
	push(id<<1,x,mid-l+1);push(id<<1|1,x,r-mid);
	tr[id].a=0;
}
void update(int x,ll v,bool assign,int id,int l,int r){
	if(l==r){if(assign)tr[id].s=v;else tr[id].s+=v;tr[id].a=0;return;}
	int mid=l+((ll)r-l)/2;pushdown(id,l,r);
	if(x<=mid)update(x,v,assign,id<<1,l,mid);
	else update(x,v,assign,id<<1|1,mid+1,r);
	pushup(id);
}
void modify(int ql,int qr,ll x,int id,int l,int r){
	if(ql<=l&&r<=qr){push(id,x,r-l+1);return;}
	int mid=l+((ll)r-l)/2;pushdown(id,l,r);
	if(ql<=mid)modify(ql,qr,x,id<<1,l,mid);
	if(mid<qr)modify(ql,qr,x,id<<1|1,mid+1,r);
	pushup(id);
}
ll query(int ql,int qr,int id,int l,int r,ll a)const{
	// carry ancestor tags; queries do not change the tree
	if(ql<=l&&r<=qr)return tr[id].s+a*(r-l+1);
	int mid=l+((ll)r-l)/2;a+=tr[id].a;
	if(qr<=mid)return query(ql,qr,id<<1,l,mid,a);
	if(mid<ql)return query(ql,qr,id<<1|1,mid+1,r,a);
	ll x=query(ql,qr,id<<1,l,mid,a),y=query(ql,qr,id<<1|1,mid+1,r,a);
	return x+y;
}
int find(ll x,bool rev)const{
	int id=1,l=vl,r=vr;ll a=0;
	while(l<r){
		int mid=l+((ll)r-l)/2;a+=tr[id].a;
		ll s=tr[id<<1|rev].s+a*(rev?r-mid:mid-l+1);
		bool right=rev;
		if(x>s){x-=s;right=!right;}
		id=id<<1|right;
		if(right)l=mid+1;else r=mid;
	}
	return l;
}
public:
void set(){vl=1;vr=0;tr.clear();}
void set(int l,int r){
	ll n=(ll)r-l+1;assert(1<=n&&n<=INT_MAX/4);
	vl=l;vr=r;tr.assign(4*n,node());
}
void setN(int n){set(1,n);}
void upd(int x,ll v){assert(vl<=x&&x<=vr);update(x,v,true,1,vl,vr);}
void add(int x,ll v){assert(vl<=x&&x<=vr);update(x,v,false,1,vl,vr);}
void add(int l,int r,ll v){
	assert(vl<=l&&(ll)l<=1ll+r&&r<=vr);
	if(l<=r)modify(l,r,v,1,vl,vr);
}
ll ask(int l,int r)const{assert(vl<=l&&l<=r&&r<=vr);return query(l,r,1,vl,vr,0);}
ll ask(int x)const{return ask(x,x);}
// all leaf values must be non-negative
int find_first(ll x)const{assert(!tr.empty()&&0<x&&x<=tr[1].s);return find(x,false);}
int find_last(ll x)const{assert(!tr.empty()&&0<x&&x<=tr[1].s);return find(x,true);}
};
// end for ds/Segment Tree/AddSum.cpp
/////////////////////////
\end{lstlisting}
\par\vspace{1em}

\Needspace{10\baselineskip}
\subsection{\texttt{ds/Segment Tree/MaxAdd.cpp}}
\textbf{区间加、区间最大值}

\TemplUsage{默认闭区间下标 1..n、初值全零，初始化 O(n)，单次修改查询 O(log n)。}
\TemplApi{\begin{itemize}
\item \texttt{setN(n)} 或 \texttt{set(l,r)} 指定有效闭区间。
\item \texttt{upd(i,v)} 是赋值，\texttt{add(i,v)} 是加法；\texttt{add(l,r,v)} 区间加。
\item \texttt{ask(l,r)} 查询，\texttt{ask(i)} 查单点。
\end{itemize}}
\textbf{最小示例（放入 main）：}
\begin{lstlisting}
MaxAdd_segt t; t.setN(3);
t.upd(2,5); t.add(1,2,2);
assert(t.ask(1,3)==7);
\end{lstlisting}
\TemplNote{这是区间加的特化实现；要更换信息或标记类型，使用外面的两份通用模板。set/setN 直接初始化全零节点，不递归建树。若初始数组不是全零，应逐点 upd；负数数组不需要改初始化，只要把所有位置正确赋值。三个类名不同，可以组合粘贴。}

\Needspace{6\baselineskip}\textbf{完整代码}
\begin{lstlisting}[firstnumber=1]
////////////////////////////////////////////////////////////////
//
// template for Segment Tree — 区间加 / 区间最大
// version 1.0 (Last Update Jun 27th, 2026)
//
// usage:
//   MaxAdd_segt st; st.setN(n); upd/add; ask(l,r) for range max
//   default leaf value 0
//
////////////////////////////////////////////////////////////////
struct MaxAdd_segt{
private:
struct node{ll s=0,a=0;};
int vl=1,vr=0;
vector<node> tr;
inline void pushup(int id){tr[id].s=max(tr[id<<1].s,tr[id<<1|1].s);}
inline void push(int id,ll x){
	tr[id].s+=x;
	tr[id].a+=x;
}
inline void pushdown(int id){
	ll x=tr[id].a;if(!x)return;
	push(id<<1,x);push(id<<1|1,x);
	tr[id].a=0;
}
void update(int x,ll v,bool assign,int id,int l,int r){
	if(l==r){if(assign)tr[id].s=v;else tr[id].s+=v;tr[id].a=0;return;}
	int mid=l+((ll)r-l)/2;pushdown(id);
	if(x<=mid)update(x,v,assign,id<<1,l,mid);
	else update(x,v,assign,id<<1|1,mid+1,r);
	pushup(id);
}
void modify(int ql,int qr,ll x,int id,int l,int r){
	if(ql<=l&&r<=qr){push(id,x);return;}
	int mid=l+((ll)r-l)/2;pushdown(id);
	if(ql<=mid)modify(ql,qr,x,id<<1,l,mid);
	if(mid<qr)modify(ql,qr,x,id<<1|1,mid+1,r);
	pushup(id);
}
ll query(int ql,int qr,int id,int l,int r){
	if(ql<=l&&r<=qr)return tr[id].s;
	int mid=l+((ll)r-l)/2;pushdown(id);
	if(qr<=mid)return query(ql,qr,id<<1,l,mid);
	if(mid<ql)return query(ql,qr,id<<1|1,mid+1,r);
	ll x=query(ql,qr,id<<1,l,mid),y=query(ql,qr,id<<1|1,mid+1,r);
	return max(x,y);
}

public:
void set(){vl=1;vr=0;tr.clear();}
void set(int l,int r){
	ll n=(ll)r-l+1;assert(1<=n&&n<=INT_MAX/4);
	vl=l;vr=r;tr.assign(4*n,node());
}
void setN(int n){set(1,n);}
void upd(int x,ll v){assert(vl<=x&&x<=vr);update(x,v,true,1,vl,vr);}
void add(int x,ll v){assert(vl<=x&&x<=vr);update(x,v,false,1,vl,vr);}
void add(int l,int r,ll v){
	assert(vl<=l&&(ll)l<=1ll+r&&r<=vr);
	if(l<=r)modify(l,r,v,1,vl,vr);
}
ll ask(int l,int r){assert(vl<=l&&l<=r&&r<=vr);return query(l,r,1,vl,vr);}
ll ask(int x){return ask(x,x);}

};
// end for ds/Segment Tree/MaxAdd.cpp
/////////////////////////
\end{lstlisting}
\par\vspace{1em}

\Needspace{10\baselineskip}
\subsection{\texttt{ds/Segment Tree/MinAdd.cpp}}
\textbf{区间加、区间最小值}

\TemplUsage{默认闭区间下标 1..n、初值全零，初始化 O(n)，单次修改查询 O(log n)。}
\TemplApi{\begin{itemize}
\item \texttt{setN(n)} 或 \texttt{set(l,r)} 指定有效闭区间。
\item \texttt{upd(i,v)} 是赋值，\texttt{add(i,v)} 是加法；\texttt{add(l,r,v)} 区间加。
\item \texttt{ask(l,r)} 查询，\texttt{ask(i)} 查单点。
\end{itemize}}
\textbf{最小示例（放入 main）：}
\begin{lstlisting}
MinAdd_segt t; t.setN(3);
t.upd(2,5); t.add(1,2,2);
assert(t.ask(1,3)==0);
\end{lstlisting}
\TemplNote{这是区间加的特化实现；要更换信息或标记类型，使用外面的两份通用模板。set/setN 直接初始化全零节点，不递归建树。若初始数组不是全零，应逐点 upd；负数数组不需要改初始化，只要把所有位置正确赋值。三个类名不同，可以组合粘贴。}

\Needspace{6\baselineskip}\textbf{完整代码}
\begin{lstlisting}[firstnumber=1]
////////////////////////////////////////////////////////////////
//
// template for Segment Tree — 区间加 / 区间最小
// version 1.0 (Last Update Jun 27th, 2026)
//
// usage:
//   MinAdd_segt st; st.setN(n); upd/add; ask(l,r) for range min
//   needs chkmin; set leaf INF in build for min queries
//
////////////////////////////////////////////////////////////////
struct MinAdd_segt{
private:
struct node{ll s=0,a=0;};
int vl=1,vr=0;
vector<node> tr;
inline void pushup(int id){tr[id].s=min(tr[id<<1].s,tr[id<<1|1].s);}
inline void push(int id,ll x){
	tr[id].s+=x;
	tr[id].a+=x;
}
inline void pushdown(int id){
	ll x=tr[id].a;if(!x)return;
	push(id<<1,x);push(id<<1|1,x);
	tr[id].a=0;
}
void update(int x,ll v,bool assign,int id,int l,int r){
	if(l==r){if(assign)tr[id].s=v;else tr[id].s+=v;tr[id].a=0;return;}
	int mid=l+((ll)r-l)/2;pushdown(id);
	if(x<=mid)update(x,v,assign,id<<1,l,mid);
	else update(x,v,assign,id<<1|1,mid+1,r);
	pushup(id);
}
void modify(int ql,int qr,ll x,int id,int l,int r){
	if(ql<=l&&r<=qr){push(id,x);return;}
	int mid=l+((ll)r-l)/2;pushdown(id);
	if(ql<=mid)modify(ql,qr,x,id<<1,l,mid);
	if(mid<qr)modify(ql,qr,x,id<<1|1,mid+1,r);
	pushup(id);
}
ll query(int ql,int qr,int id,int l,int r){
	if(ql<=l&&r<=qr)return tr[id].s;
	int mid=l+((ll)r-l)/2;pushdown(id);
	if(qr<=mid)return query(ql,qr,id<<1,l,mid);
	if(mid<ql)return query(ql,qr,id<<1|1,mid+1,r);
	ll x=query(ql,qr,id<<1,l,mid),y=query(ql,qr,id<<1|1,mid+1,r);
	return min(x,y);
}

public:
void set(){vl=1;vr=0;tr.clear();}
void set(int l,int r){
	ll n=(ll)r-l+1;assert(1<=n&&n<=INT_MAX/4);
	vl=l;vr=r;tr.assign(4*n,node());
}
void setN(int n){set(1,n);}
void upd(int x,ll v){assert(vl<=x&&x<=vr);update(x,v,true,1,vl,vr);}
void add(int x,ll v){assert(vl<=x&&x<=vr);update(x,v,false,1,vl,vr);}
void add(int l,int r,ll v){
	assert(vl<=l&&(ll)l<=1ll+r&&r<=vr);
	if(l<=r)modify(l,r,v,1,vl,vr);
}
ll ask(int l,int r){assert(vl<=l&&l<=r&&r<=vr);return query(l,r,1,vl,vr);}
ll ask(int x){return ask(x,x);}

};
// end for ds/Segment Tree/MinAdd.cpp
/////////////////////////
\end{lstlisting}
\par\vspace{1em}

\Needspace{10\baselineskip}
\subsection{\texttt{ds/Segment Tree(with lazytag).cpp}}
\textbf{通用线段树（带懒标记）}

\TemplUsage{默认维护区间和，下标 1..n，初值全零。递归部分只处理 info，按题目改写信息与标记的运算。}
\TemplApi{\begin{itemize}
\item \texttt{segt t; t.setN(n)} 初始化。
\item \texttt{upd(i,v)} 单点赋值；\texttt{ask(l,r)} 区间和，\texttt{ask(i)} 单点值。
\item \texttt{add(i,v)} 单点加，\texttt{add(l,r,v)} 区间加。
\end{itemize}}
\textbf{最小示例（放入 main）：}
\begin{lstlisting}
segt t; t.setN(3); t.upd(2,5);
t.add(1,2,2); assert(t.ask(1,3)==9);
\end{lstlisting}
\TemplNote{info + info 按左右顺序合并；upd 在叶子整体替换 info，单点 add 使用 info(delta,0)，通过 info += info 更新叶子。初值在 build 的 info(0,1) 处修改。参数转换、提取答案只改公开接口。USE\_COMPARE 默认开启；自定义信息不支持减法与比较时，删除文件顶部的 define 行，对应运算和二分函数就不编译。默认二分要求所有点值非负、目标为正且不超过总和。set/setN 清空重建。两版同名 segt，只选一份。 lazt += lazt 的顺序是旧标记后接新标记；info += lazt 负责作用标记，区间长度等元数据由 info 保存。不能仅因支持单点加，就直接认为相同统计能支持区间加。}

\Needspace{6\baselineskip}\textbf{完整代码}
\begin{lstlisting}[firstnumber=1]
////////////////////////////////////////////////////////////////
//
// template for Segment Tree
// version 1.1 (Last Update Jun 27th, 2026)
//
// usage:
//   segt st; st.setN(n);  // or set(l,r), 1-indexed
//   upd/add point; add(l,r,v) range; ask; find_first/last
//
////////////////////////////////////////////////////////////////
// remove this line when info has no subtraction / comparison
#define USE_COMPARE
struct segt{
private:
//////////////////////////////////////////////////////////////
// basic operation
struct lazt{
	// modify here
	ll a;
	bool empty()const{return a==0;}
	void set(){a=0;}
	lazt():a(0){}
	lazt(ll _a):a(_a){}
	inline lazt& operator +=(const lazt &y){
		a+=y.a;
		return *this;
	}
	inline lazt operator +(const lazt &y)const{
		return lazt(*this)+=y;
	}
};
struct info{
	// modify here
	ll s,len;
	info():s(0),len(0){}
	info(ll _s,ll _len):s(_s),len(_len){}
	inline info& operator +=(const info &y){
		s+=y.s;len+=y.len;
		return *this;
	}
	inline info operator +(const info &y)const{
		return info(*this)+=y;
	}
	inline info& operator +=(const lazt &y){
		s+=len*y.a;
		return *this;
	}
	inline info operator +(const lazt &y)const{
		return info(*this)+=y;
	}
#ifdef USE_COMPARE
	// for find_first or find_last
	inline info& operator -=(const info &y){
		s-=y.s;len-=y.len;
		return *this;
	}
	inline info operator -(const info &y)const{
		return info(*this)-=y;
	}
	inline bool operator <= (const info &a)const{
		return s<=a.s;
	}
#endif
};
/////////////////////////////////////////////////////////////////
// basic segment tree template, private operations
int vl=1,vr=0;
vector<info> val;
vector<lazt> laz;
inline void push(int id,const lazt &x){
	laz[id]+=x;val[id]+=x;
}
inline void pushdown(int id){
	if(!laz[id].empty()) push(id<<1,laz[id]),push(id<<1|1,laz[id]),laz[id].set();
}
inline void pushup(int id){
	val[id]=val[id<<1]+val[id<<1|1];
}
void build(int id,int l,int r){
	laz[id]=lazt();
	if(l==r){
		val[id]=info(0,1); // please modify this init value
		return;
	}
	int mid=l+((ll)r-l)/2;
	build(id<<1,l,mid);build(id<<1|1,mid+1,r);
	pushup(id);
}
void update(int x,const info &v,int id,int l,int r){
	if(l==r){
		assert(x==l);
		val[id]=v;laz[id].set();return;
	}
	int mid=l+((ll)r-l)/2;
	pushdown(id);
	if(x<=mid) update(x,v,id<<1,l,mid);
	else update(x,v,id<<1|1,mid+1,r);
	pushup(id);
}
void modify(int x,const info &v,int id,int l,int r){
	if(l==r){
		assert(x==l);
		val[id]+=v;return;
	}
	int mid=l+((ll)r-l)/2;
	pushdown(id);
	if(x<=mid) modify(x,v,id<<1,l,mid);
	else modify(x,v,id<<1|1,mid+1,r);
	pushup(id);
}
void modify(int ql,int qr,const lazt &v,int id,int l,int r){
	if(ql<=l&&r<=qr){push(id,v);return;}
	int mid=l+((ll)r-l)/2;
	pushdown(id);
	if(ql<=mid) modify(ql,qr,v,id<<1,l,mid);
	if(mid<qr) modify(ql,qr,v,id<<1|1,mid+1,r);
	pushup(id);
}
info query(int ql,int qr,int id,int l,int r){
	if(ql<=l&&r<=qr) return val[id];
	int mid=l+((ll)r-l)/2;
	pushdown(id);
	if(qr<=mid) return query(ql,qr,id<<1,l,mid);
	if(mid<ql) return query(ql,qr,id<<1|1,mid+1,r);
	return query(ql,qr,id<<1,l,mid)+query(ql,qr,id<<1|1,mid+1,r);
}
#ifdef USE_COMPARE
int find_first(info x,int id,int l,int r){
	// if no compare betweeen infos, delete this
	if(l==r) return l;
	int mid=l+((ll)r-l)/2;
	pushdown(id);
	return x<=val[id<<1]?find_first(x,id<<1,l,mid):find_first(x-val[id<<1],id<<1|1,mid+1,r);
}
int find_last(info x,int id,int l,int r){
	// if no compare betweeen infos, delete this
	if(l==r) return l;
	int mid=l+((ll)r-l)/2;
	pushdown(id);
	return x<=val[id<<1|1]?find_last(x,id<<1|1,mid+1,r):find_last(x-val[id<<1|1],id<<1,l,mid);
}
#endif
/////////////////////////////////////////////////////////////
// public operations, remember to setN()/set()
// convert arguments to info / lazt, and extract the answer here
public:
void set(){vl=1;vr=0;val.clear();laz.clear();}
void set(int _l,int _r){
	ll n=(ll)_r-_l+1;assert(1<=n&&n<=INT_MAX/4);
	vl=_l;vr=_r;val.assign(4*n,info());
	laz.assign(4*n,lazt());
	build(1,vl,vr);
}
void setN(int _n){
	set(1,_n);
}
void upd(int x,ll a){
	assert(vl<=x&&x<=vr);
	update(x,info(a,1),1,vl,vr);
}
void add(int x,ll a){
	assert(vl<=x&&x<=vr);
	modify(x,info(a,0),1,vl,vr);
}
void add(int ql,int qr,ll a){
	assert(vl<=ql&&(ll)ql<=1ll+qr&&qr<=vr);
	if(ql<=qr) modify(ql,qr,lazt(a),1,vl,vr);
}
ll ask(int ql,int qr){
	assert(vl<=ql&&ql<=qr&&qr<=vr);
	return query(ql,qr,1,vl,vr).s;
}
ll ask(int x){
	assert(vl<=x&&x<=vr);
	// query a[x]
	return query(x,x,1,vl,vr).s;
}
#ifdef USE_COMPARE
int find_first(ll x){
	assert(!val.empty()&&x>0&&info(x,0)<=val[1]);
	return find_first(info(x,0),1,vl,vr);
}
int find_last(ll x){
	assert(!val.empty()&&x>0&&info(x,0)<=val[1]);
	return find_last(info(x,0),1,vl,vr);
}
#endif
};

#undef USE_COMPARE
// end for ds/Segment Tree(with lazytag).cpp
/////////////////////////
\end{lstlisting}
\par\vspace{1em}

\Needspace{10\baselineskip}
\subsection{\texttt{ds/Segment Tree(without lazytag).cpp}}
\textbf{通用线段树（不带懒标记）}

\TemplUsage{默认维护区间和，下标 1..n，初值全零。递归部分只处理 info，按题目改写信息与标记的运算。}
\TemplApi{\begin{itemize}
\item \texttt{segt t; t.setN(n)} 初始化。
\item \texttt{upd(i,v)} 单点赋值；\texttt{ask(l,r)} 区间和，\texttt{ask(i)} 单点值。
\item \texttt{add(i,v)} 单点加，不支持区间加。
\end{itemize}}
\textbf{最小示例（放入 main）：}
\begin{lstlisting}
segt t; t.setN(3); t.upd(2,5);
t.add(2,2); assert(t.ask(1,3)==7);
\end{lstlisting}
\TemplNote{info + info 按左右顺序合并；upd 在叶子整体替换 info，单点 add 使用 info(delta,0)，通过 info += info 更新叶子。初值在 build 的 info(0,1) 处修改。参数转换、提取答案只改公开接口。USE\_COMPARE 默认开启；自定义信息不支持减法与比较时，删除文件顶部的 define 行，对应运算和二分函数就不编译。默认二分要求所有点值非负、目标为正且不超过总和。set/setN 清空重建。两版同名 segt，只选一份。}


\textbf{改写示例：单点加、区间最大子段和（非空）}
删除顶部的 \texttt{\#define USE\_COMPARE}，把 info 替换为下面这段；区间 ask 的返回字段从 s 改为 mx，其他部分不改。len=0 的参数表示单点增量，len=1 表示叶子；合并的左右区间均非空。
\begin{lstlisting}
struct info{
	ll s,pre,suf,mx,len;
	info(ll x=0,ll n=0):s(x),pre(x),suf(x),mx(x),len(n){}
	inline info& operator +=(const info &y){
		if(!y.len){ // point add: y is info(delta,0)
			s+=y.s;pre=suf=mx=s;return *this;
		}
		mx=max({mx,y.mx,suf+y.pre});
		pre=max(pre,s+y.pre);suf=max(y.suf,y.s+suf);
		s+=y.s;len+=y.len;return *this;
	}
	inline info operator +(const info &y)const{return info(*this)+=y;}
};
\end{lstlisting}
\Needspace{6\baselineskip}\textbf{完整代码}
\begin{lstlisting}[firstnumber=1]
////////////////////////////////////////////////////////////////
//
// template for Segment Tree
// version 1.1 (Last Update Jun 27th, 2026)
//
// usage:
//   segt st; st.setN(n);  // or set(l,r), 1-indexed
//   upd/add point; ask(l,r); find_first/last
//
////////////////////////////////////////////////////////////////
// remove this line when info has no subtraction / comparison
#define USE_COMPARE
struct segt{
private:
//////////////////////////////////////////////////////////////
// basic operation
struct info{
	// modify here
	ll s,len;
	info():s(0),len(0){}
	info(ll _s,ll _len):s(_s),len(_len){}
	inline info& operator +=(const info &y){
		s+=y.s;len+=y.len;
		return *this;
	}
	inline info operator +(const info &y)const{
		return info(*this)+=y;
	}
#ifdef USE_COMPARE
	// for find_first or find_last
	inline info& operator -=(const info &y){
		s-=y.s;len-=y.len;
		return *this;
	}
	inline info operator -(const info &y)const{
		return info(*this)-=y;
	}
	inline bool operator <= (const info &a)const{
		return s<=a.s;
	}
#endif
};
/////////////////////////////////////////////////////////////////
// basic segment tree template, private operations
int vl=1,vr=0;
vector<info> val;
inline void pushup(int id){
	val[id]=val[id<<1]+val[id<<1|1];
}
void build(int id,int l,int r){
	if(l==r){
		val[id]=info(0,1); // please modify this init value
		return;
	}
	int mid=l+((ll)r-l)/2;
	build(id<<1,l,mid);build(id<<1|1,mid+1,r);
	pushup(id);
}
void update(int x,const info &v,int id,int l,int r){
	if(l==r){
		assert(x==l);
		val[id]=v;return;
	}
	int mid=l+((ll)r-l)/2;
	if(x<=mid) update(x,v,id<<1,l,mid);
	else update(x,v,id<<1|1,mid+1,r);
	pushup(id);
}
void modify(int x,const info &v,int id,int l,int r){
	if(l==r){
		assert(x==l);
		val[id]+=v;return;
	}
	int mid=l+((ll)r-l)/2;
	if(x<=mid) modify(x,v,id<<1,l,mid);
	else modify(x,v,id<<1|1,mid+1,r);
	pushup(id);
}
info query(int ql,int qr,int id,int l,int r){
	if(ql<=l&&r<=qr) return val[id];
	int mid=l+((ll)r-l)/2;
	if(qr<=mid) return query(ql,qr,id<<1,l,mid);
	if(mid<ql) return query(ql,qr,id<<1|1,mid+1,r);
	return query(ql,qr,id<<1,l,mid)+query(ql,qr,id<<1|1,mid+1,r);
}
#ifdef USE_COMPARE
int find_first(info x,int id,int l,int r){
	// if no compare betweeen infos, delete this
	if(l==r) return l;
	int mid=l+((ll)r-l)/2;
	return x<=val[id<<1]?find_first(x,id<<1,l,mid):find_first(x-val[id<<1],id<<1|1,mid+1,r);
}
int find_last(info x,int id,int l,int r){
	// if no compare betweeen infos, delete this
	if(l==r) return l;
	int mid=l+((ll)r-l)/2;
	return x<=val[id<<1|1]?find_last(x,id<<1|1,mid+1,r):find_last(x-val[id<<1|1],id<<1,l,mid);
}
#endif
/////////////////////////////////////////////////////////////
// public operations, remember to setN()/set()
// convert arguments to info / lazt, and extract the answer here
public:
void set(){vl=1;vr=0;val.clear();}
void set(int _l,int _r){
	ll n=(ll)_r-_l+1;assert(1<=n&&n<=INT_MAX/4);
	vl=_l;vr=_r;val.assign(4*n,info());
	build(1,vl,vr);
}
void setN(int _n){
	set(1,_n);
}
void upd(int x,ll a){
	assert(vl<=x&&x<=vr);
	update(x,info(a,1),1,vl,vr);
}
void add(int x,ll a){
	assert(vl<=x&&x<=vr);
	modify(x,info(a,0),1,vl,vr);
}
ll ask(int ql,int qr){
	assert(vl<=ql&&ql<=qr&&qr<=vr);
	return query(ql,qr,1,vl,vr).s;
}
ll ask(int x){
	assert(vl<=x&&x<=vr);
	// query a[x]
	return query(x,x,1,vl,vr).s;
}
#ifdef USE_COMPARE
int find_first(ll x){
	assert(!val.empty()&&x>0&&info(x,0)<=val[1]);
	return find_first(info(x,0),1,vl,vr);
}
int find_last(ll x){
	assert(!val.empty()&&x>0&&info(x,0)<=val[1]);
	return find_last(info(x,0),1,vl,vr);
}
#endif
};

#undef USE_COMPARE
// end for ds/Segment Tree(without lazytag).cpp
/////////////////////////
\end{lstlisting}
\par\vspace{1em}

\Needspace{10\baselineskip}
\subsection{\texttt{ds/zkw.cpp}}
\textbf{非递归线段树框架}

\TemplUsage{单点赋值和闭区间聚合，当前 info 维护区间和与最大前缀和（允许空前缀）。单次 O(log n)。}
\TemplApi{\begin{itemize}
\item \texttt{zkw\_segt t; t.setN(n)} 下标为 0..n-1。
\item \texttt{upd(i,a,b)} 设置该叶 info 的 s=a、mx=b。
\item 合并两段时 s=L.s+R.s，mx=max(L.mx,L.s+R.mx)；ask 返回最终 mx。
\end{itemize}}
\textbf{最小示例（放入 main）：}
\begin{lstlisting}
zkw_segt t; t.setN(3);
t.upd(0,-4,0); t.upd(1,9,9); t.upd(2,-2,0);
assert(t.ask(0,2)==5);
\end{lstlisting}
\TemplNote{普通数组的叶子设为 s=a[i]、mx=max(0,a[i])。若 s 全为零且 mx 非负，也可用于区间最大值。其他统计应重写 info 的合并和单位元；合并必须满足结合律。下标从 0 开始。}

\Needspace{6\baselineskip}\textbf{完整代码}
\begin{lstlisting}[firstnumber=1]
//
// template for ZKW segment tree
// version 1.0 (Last Update Jun 13th, 2026)
//
// usage:
//   zkw_segt z; z.setN(n); z.upd(i, s, mx); z.ask(l, r);
//   0-indexed; leaf mx=max(0ll,s); ask returns maximum prefix sum
//
struct zkw_segt{
// information
struct info{
	// sum and maximum prefix sum, including the empty prefix
	ll s,mx;
	info():s(0),mx(0){}
	info(ll _s,ll _mx):s(_s),mx(_mx){}
	inline info& operator +=(const info &y){
		ll os=s;
		s+=y.s;
		mx=max(mx,os+y.mx);
		return *this;
	}
	inline info operator +(const info &y)const{
		return info(*this)+=y;
	}
};
// basic segment tree template, private operations
private:
int n=0,B=0;vector<info> a;
void modify(int x,info v){
	x+=B+1;
	a[x]=v;
	while(x>1) x>>=1,a[x]=a[x<<1]+a[x<<1|1]; 
}
info query(int l,int r){
	l+=B,r+=B+2;
	info ql,qr;
	while(l^r^1){
		if(~l&1) ql=ql+a[l^1];
		if(r&1) qr=a[r^1]+qr;
		l>>=1,r>>=1;
	}
	return ql+qr;
}
// basic segment tree template, public operations
// please transfer (int,int) to info here
public:
void setN(int _n){
	assert(1<=_n&&_n<(1<<29));n=_n;B=1;
	while(B<n+2)B*=2;
	a.assign(2*B,info());
}
void upd(int x,ll a,ll b){
	assert(0<=x&&x<n);
	modify(x,info(a,b));
}
ll ask(int l,int r){
	assert(0<=l&&l<=r&&r<n);
	return query(l,r).mx;
}
};
\end{lstlisting}
\par\vspace{1em}

\Needspace{10\baselineskip}
\subsection{\texttt{ds/Dynamic Segment Tree.cpp}}
\textbf{大值域稀疏单点统计}

\TemplUsage{在整数键 -10\textasciicircum{}9..10\textasciicircum{}9 上做单点加、区间和；节点带键交换以压缩空间。单次 O(log V)。}
\TemplApi{\begin{itemize}
\item \texttt{dynamic\_segt t; t.set()} 清空；\texttt{set(m)} 可预留节点。
\item \texttt{add(x,v)} 给键 x 的权值加 v；\texttt{ask(l,r)} 求键区间权值和。
\end{itemize}}
\textbf{最小示例（放入 main）：}
\begin{lstlisting}
dynamic_segt t; t.set();
t.add(-100,3); t.add(200,4);
assert(t.ask(-100,0)==3);
\end{lstlisting}
\TemplNote{空键权值为零。与主席树不同：不保留历史版本，也不支持复用旧根。文件有全局对象 T。}

\Needspace{6\baselineskip}\textbf{完整代码}
\begin{lstlisting}[firstnumber=1]
//
// template for Dynamic Segment Tree
// this code accept single-point modify, range query, cost memory O(N)
// version 1.0 (Last Update Jun 16th, 2026)
//
// usage:
//   dynamic_segt T; T.set(); T.add(pos, val); T.ask(l, r);
//   sparse segt on value domain [L,R]
//
struct dynamic_segt{
private:
static const int L=-1e9,R=1e9;
int tot=0;
struct info{
	ll s;
	info():s(0){}
	info(ll _s):s(_s){}
	inline info& operator +=(const info &x){s+=x.s;return *this;}
	inline info operator +(const info &x)const{return info(*this)+=x;}
};
struct node{
	int p=0,ls=0,rs=0;
	info self,val;
};
vector<node> tr=vector<node>(1);
int rt=0;
int new_node(int p,info x){
	assert(tr.size()<INT_MAX);tr.push_back({p,0,0,x,x});
	return ++tot;
}
inline void pushup(int id){
	tr[id].val=tr[tr[id].ls].val+tr[id].self+tr[tr[id].rs].val;
}
int add(int p,info v,int id,int l,int r){
	if(!id)return new_node(p,v);
	if(tr[id].p==p){tr[id].self+=v;pushup(id);return id;}
	int mid=l+((ll)r-l)/2;
	if(p<=mid){
		if(p>tr[id].p)swap(tr[id].p,p),swap(tr[id].self,v);
		int x=add(p,v,tr[id].ls,l,mid);tr[id].ls=x;
	}else{
		if(p<tr[id].p)swap(tr[id].p,p),swap(tr[id].self,v);
		int x=add(p,v,tr[id].rs,mid+1,r);tr[id].rs=x;
	}
	pushup(id);return id;
}
void query(int ql,int qr,int id,int l,int r,info &X){
	if(max(ql,l)>min(r,qr)||!id) return;
	if(ql<=l&&r<=qr){X+=tr[id].val;return;}
	int mid=l+((ll)r-l)/2;
	query(ql,qr,tr[id].ls,l,mid,X);
	if(ql<=tr[id].p&&tr[id].p<=qr) X+=tr[id].self;
	query(ql,qr,tr[id].rs,mid+1,r,X);
}
public:
void set(int m=0){assert(m>=0);tot=rt=0;tr.assign(1,node());tr.reserve((size_t)m+1);}
void add(int p,ll s){
	assert(L<=p&&p<=R);
	rt=add(p,info(s),rt,L,R);
}
ll ask(int l,int r){
	assert(L<=l&&l<=r&&r<=R);
	info X;
	query(l,r,rt,L,R,X);
	return X.s;
}
}T;
\end{lstlisting}
\par\vspace{1em}

\Needspace{10\baselineskip}
\subsection{\texttt{ds/Segment Tree Merge.cpp}}
\textbf{可合并的多棵线段树}

\TemplUsage{维护多棵互不共享节点的树，值域 0..n。当前权值设计用于统计所有键 y>=x 的 (y-x)*权值。}
\TemplApi{\begin{itemize}
\item \texttt{segt\_merge t; t.setN(n)} 初始化池，外部根变量设为 0。
\item \texttt{add(rt,y,w)} 给键 y 加权 w，根按引用更新。
\item \texttt{ask(rt,x)} 返回 sum\_\{y>=x\}(y-x)*w[y]。
\item \texttt{merge(a,b)} 把 b 合入 a 并把 b 置零。
\end{itemize}}
\textbf{最小示例（放入 main）：}
\begin{lstlisting}
segt_merge t; t.setN(10); int a=0,b=0;
t.add(a,5,2); t.add(b,8,1); t.merge(a,b);
assert(b==0 && t.ask(a,3)==9);
\end{lstlisting}
\TemplNote{这是破坏性合并，不是可持久化。合并后不能再使用旧根副本。单点操作 O(log V)，合并成本取决于重叠节点数；若想查询普通权值和，要改 info/query 的对外接口。}

\Needspace{6\baselineskip}\textbf{完整代码}
\begin{lstlisting}[firstnumber=1]
//
// template for Segment Tree Merge
// version 1.0 (Last Update Jun 25th, 2026)
//
// usage:
//   segt_merge sm; sm.setN(n); int rt=0;
//   sm.add(rt,i,v); sm.ask(rt,x); sm.merge(a,b);
//
struct segt_merge{
private:
struct info{
	// modify here
	ll s,s1;
	info():s(0),s1(0){}
	info(ll _s,ll _s1):s(_s),s1(_s1){}
	inline info& operator +=(const info &y){
		s+=y.s;s1+=y.s1;
		return *this;
	}
	inline info operator +(const info &y)const{
		return info(*this)+=y;
	}
};
vector<info> val=vector<info>(1);
vector<int> ls{0},rs{0};int n=0,tot=0;
inline void pushup(int id){val[id]=val[ls[id]]+val[rs[id]];}
int add(int x,info v,int id,int l,int r){
	if(!id){
		assert(tot<INT_MAX-1);id=++tot;
		val.emplace_back();ls.push_back(0);rs.push_back(0);
	}
	if(l==r){val[id]+=v;return id;}
	int mid=l+(r-l)/2;
	if(x<=mid){int t=add(x,v,ls[id],l,mid);ls[id]=t;}
	else {int t=add(x,v,rs[id],mid+1,r);rs[id]=t;}
	pushup(id);return id;
}
info query(int ql,int qr,int id,int l,int r){
	if(!id||max(l,ql)>min(r,qr)) return info(); 
	if(ql<=l&&r<=qr) return val[id];
	int mid=l+(r-l)/2;
	return query(ql,qr,ls[id],l,mid)+query(ql,qr,rs[id],mid+1,r);
}
void merge(int &u,int &v,int l,int r){
	if(!u||!v){u|=v;return;}
	if(l==r){val[u]+=val[v];return;}
	int mid=l+(r-l)/2;
	merge(ls[u],ls[v],l,mid);
	merge(rs[u],rs[v],mid+1,r);
	pushup(u);
}
public:
void setN(int _n){
	assert(_n>=0);n=_n;tot=0;
	val.assign(1,info());ls.assign(1,0);rs.assign(1,0);
}
void add(int &id,int x,int o){
	assert(0<=x&&x<=n);
	assert(0<=id&&id<=tot);
	id=add(x,info(o,1ll*x*o),id,0,n);
}
ll ask(int &id,int x){
	assert(0<=id&&id<=tot&&0<=x&&x<=n);
	auto it=query(x,n,id,0,n);
	return it.s1-it.s*x;
}
void merge(int &x,int &y){
	assert(0<=x&&x<=tot&&0<=y&&y<=tot&&(&x!=&y)&&(!x||x!=y));
	merge(x,y,0,n);
	y=0;
}
};
\end{lstlisting}
\par\vspace{1em}

\Needspace{10\baselineskip}
\subsection{\texttt{ds/lichao.cpp}}
\textbf{线段李超树}

\TemplUsage{插入只在一段整数横坐标上有效的直线，查询某点的最小值。}
\TemplApi{\begin{itemize}
\item \texttt{lichao t; t.set()} 清空。
\item \texttt{add(line\{k,b\},l,r)} 插入 y=kx+b；省略 l,r 表示整个默认值域。
\item \texttt{ask(x)} 返回最小纵坐标，没有覆盖返回 2e18。
\end{itemize}}
\textbf{最小示例（放入 main）：}
\begin{lstlisting}
lichao t; t.set();
t.add({2,3},0,10); t.add({-1,8},0,10);
assert(t.ask(2)==6);
\end{lstlisting}
\TemplNote{横坐标及斜率限制在 [-10\textasciicircum{}9,10\textasciicircum{}9]，有效端点的纵坐标需在 [-2e18,2e18]。整条线插入 O(log V)，线段插入 O(log\textasciicircum{}2 V)，查询 O(log V)。}

\Needspace{6\baselineskip}\textbf{完整代码}
\begin{lstlisting}[firstnumber=1]
//
// template for lichao segment tree
// version 1.2 (Last Update Jun 13th, 2026)
//
// usage:
//   lichao lc; lc.set(); lc.add({k,b}); lc.add({k,b}, l, r);
//   ll y = lc.ask(x);  // min y=kx+b, INF if empty
//
struct line{
	ll k,b;
	ll f(ll x){return k*x+b;}
	// please, check whether use __int128
};
struct lichao{
// you may add y=kb+x where x \in [l,r]
// you can query x, answer the **min** y (if no such, return inf = 2e18)
static const ll inf=2e18,V=1e9;
vector<int> ls{0},rs{0},p{0};int tot=0,cnt=0,rt=0;
vector<line> L{{0,inf}};
// if you don't use range add, this can up to n
inline ll gv(const int &id,const ll &v){return L[id].f(v);}
void set(){
	ls.assign(1,0);rs.assign(1,0);p.assign(1,0);
	L.assign(1,{0,inf});tot=cnt=rt=0;
}
int new_node(){
	assert(tot<INT_MAX-1);ls.push_back(0);rs.push_back(0);p.push_back(0);
	return ++tot;
}
int add(int x,int id,int l,int r){
	if(!id)id=new_node();
	if(l==r){
		if(gv(x,l)<gv(p[id],l))p[id]=x;
		return id;
	}
	int mid=l+((ll)r-l)/2;
	if(gv(x,mid)<gv(p[id],mid))swap(x,p[id]);
	if(gv(x,l)<gv(p[id],l)){int t=add(x,ls[id],l,mid);ls[id]=t;}
	if(gv(x,r)<gv(p[id],r)){int t=add(x,rs[id],mid+1,r);rs[id]=t;}
	return id;
}
int upd(int x,int ql,int qr,int id,int l,int r){
	if(max(ql,l)>min(r,qr))return id;
	if(!id)id=new_node();
	if(ql<=l&&r<=qr)return add(x,id,l,r);
	int mid=l+((ll)r-l)/2;
	int a=upd(x,ql,qr,ls[id],l,mid),b=upd(x,ql,qr,rs[id],mid+1,r);
	ls[id]=a;rs[id]=b;return id;
}
ll ask(int x,int id,int l,int r){
	if(!id) return inf;
	ll ans=gv(p[id],x);
	if(l==r) return ans;
	int mid=l+((ll)r-l)/2;
	if(x<=mid) ans=min(ans,ask(x,ls[id],l,mid));
	else ans=min(ans,ask(x,rs[id],mid+1,r));
	return ans;
}


void add(line x,int l=-V,int r=V){
	assert(-V<=l&&l<=r&&r<=V);
	assert(-V<=x.k&&x.k<=V);
	assert((__int128)x.k*l+x.b>=-inf&&(__int128)x.k*l+x.b<=inf);
	assert((__int128)x.k*r+x.b>=-inf&&(__int128)x.k*r+x.b<=inf);
	assert(cnt<INT_MAX-1);L.push_back(x);++cnt;
	rt=upd(cnt,l,r,rt,-V,V);
}
ll ask(int x){
	assert(-V<=x&&x<=V);
	return ask(x,rt,-V,V);
}
};
\end{lstlisting}
\par\vspace{1em}

\Needspace{10\baselineskip}
\subsection{\texttt{ds/ODT.cpp}}
\textbf{有序区间集合}

\TemplUsage{维护若干不相交的常值闭区间。适合区间覆盖以及依赖当前分段的统计。}
\TemplApi{\begin{itemize}
\item \texttt{ODT t; t.set()} 清空；\texttt{add(l,r,c)} 只插入一个区间，不会自动清除旧区间。
\item 赋值应先 \texttt{erase(l,r)} 再 \texttt{add(l,r,c)}；erase 返回被覆盖的旧段。
\item \texttt{extract(l,r)} 返回当前区间片段，每段字段为 l,r,c。
\end{itemize}}
\textbf{最小示例（放入 main）：}
\begin{lstlisting}
ODT t; t.set(); t.add(1,5,0);
t.erase(2,4); t.add(2,4,7);
auto a=t.extract(2,3);
assert(a.size()==1 && a[0].c==7);
\end{lstlisting}
\TemplNote{extract 会切开边界，可能增加内部段数；没有自动合并相邻同色段。一次操作可能遍历很多段，不保证对任意输入都是对数复杂度。}

\Needspace{6\baselineskip}\textbf{完整代码}
\begin{lstlisting}[firstnumber=1]
//
// template for ODT
// Chtholly is very cute!
// version 1.0 (Last Update Jun 16th, 2026)
//
// usage:
//   ODT odt; odt.set(); odt.erase(l,r); odt.add(l,r,c);
//   odt.extract(l,r);  // assign color on [l,r]
//
struct seg{
	int l,r,c;
	seg():l(0),r(0),c(0){}
	seg(int _l,int _r,int _c):l(_l),r(_r),c(_c){}
	inline bool operator < (const seg &x) const {return l<x.l;}
};
struct ODT{
private:
std::set<seg> S;
void cut(int y){
	auto it=S.upper_bound(seg(y,0,0));
	if(it==S.begin()) return;
	auto [l,r,c]=*(--it);
	if(l<y&&y<=r) S.erase(it),S.insert(seg(l,y-1,c)),S.insert(seg(y,r,c));
}
public:
void set(){S.clear();}
vector<seg> erase(int l,int r){
	// erase segments from [l,r]
	assert(l<=r);
	cut(l);if(r<INT_MAX)cut(r+1);
	auto it=S.lower_bound(seg(l,0,0));
	vector<seg> I;
	while(1){
		if(it==S.end()||it->l>r) break;
		I.push_back(*it);
		it=S.erase(it);
	}
	return I;
}
vector<seg> extract(int l,int r){
	// extract segments from [l,r]
	// before you use this function, pay attention is the TC correct?
	assert(l<=r);
	cut(l);if(r<INT_MAX)cut(r+1);
	auto it=S.lower_bound(seg(l,0,0));
	vector<seg> I;
	while(1){
		if(it==S.end()||it->l>r) break;
		I.push_back(*it);
		it++;
	}
	return I;
}
void add(int l,int r,int c){
	assert(l<=r);
	S.insert(seg(l,r,c));
}
};
// !!!!!!!!!! before add segments, erase the range first!!!!!!!!!!
\end{lstlisting}
\par\vspace{1em}

\Needspace{10\baselineskip}
\subsection{\texttt{ds/FastSet.cpp}}
\textbf{有界整数集合}

\TemplUsage{用分层位集维护集合成员及前驱后继。修改和前驱后继 O(log\_64 V)，成员判断 O(1)。}
\TemplApi{\begin{itemize}
\item \texttt{FastSet s; s.setN(v)} 允许存储 0..v，包含 v。
\item \texttt{insert/erase} 幂等，\texttt{s[x]} 判断是否存在。
\item \texttt{prev(x)} 为最大 <=x 的元素，\texttt{next(x)} 为最小 >=x 的元素；没有则返回 -1。
\item \texttt{front/back} 返回最小/最大元素；\texttt{to\_vector()} 扫描整个值域输出元素。
\end{itemize}}
\textbf{最小示例（放入 main）：}
\begin{lstlisting}
FastSet s; s.setN(100);
s.insert(3); s.insert(100);
assert(s.next(3)==3 && s.prev(99)==3);
assert(s.next(101)==-1);
\end{lstlisting}
\TemplNote{值域太大时不适合使用，空间跟 V 相关而非当前元素个数。需要 src 的位运算和最值辅助函数。}

\Needspace{6\baselineskip}\textbf{完整代码}
\begin{lstlisting}[firstnumber=1]
//
// template for fast set, insert/erase a non-negative integer in O(log_w(V))
// query a integer exist in O(1) 
// version 1.2 (Last Update Jun 13th, 2026)
//
// usage:
//   FastSet fs; fs.setN(n); fs.insert(x); fs.erase(x);
//   fs[x]; fs.prev/next; needs chkmin/chkmax/first_bit/last_bit
//
struct FastSet{
vector<vector<ull>> a;int V=0,B=0;
void setN(int _V){
	assert(0<=_V&&_V<INT_MAX);V=_V+1;a.clear();
	for(int n=V;;){
		n=(n-1)/64+1;a.emplace_back(n,0);
		if(n==1)break;
	}
	B=(int)a.size()-1;
}
inline bool operator [](int x)const{assert(0<=x&&x<V);return a[0][x>>6]>>(x&63)&1;}
vector<int> to_vector()const{
	vector<int> I;
	if(a.empty())return I;
	for(int i=0;i<(int)a[0].size();i++)
		for(ull x=a[0][i];x;x&=x-1)I.push_back(i*64+first_bit(x));
	return I;
}
void insert(int x){
	assert(0<=x&&x<V);
	if((*this)[x]) return;
	for(int i=0;i<=B;i++){
		a[i][x>>6]|=(1ull<<(x&63));
		x>>=6;
	}
}
void erase(int x){
	assert(0<=x&&x<V);
	if(!(*this)[x]) return;
	for(int i=0;i<=B;i++){
		a[i][x>>6]&=~(1ull<<(x&63));
		if(a[i][x>>6]) break;
		x>>=6;
	}
}
int prev(int x){
	// find the first element <= x, if no, return -1
	if(x<0||a.empty())return -1;
	chkmin(x,V-1);
	for(int h=0;h<=B;h++){
		if(x<0||(x>>6)>=(int)a[h].size()) break;
		ull f=a[h][x>>6]<<(63-(x&63));
		if(!f){
			x=(x>>6)-1;
			continue;
		}
		x-=63-last_bit(f);
		for(int g=h-1;g>=0;g--){
			x<<=6;
			x+=last_bit(a[g][x>>6]);
		}
		return x;
	}
	return -1;
}
int next(int x){
	// find the first element >=x, if no, return -1
	if(x>=V||a.empty())return -1;
	chkmax(x,0);
	for(int h=0;h<=B;h++){
		if(x<0||(x>>6)>=(int)a[h].size()) break;
		ull f=a[h][x>>6]>>(x&63);
		if(!f){
			x=(x>>6)+1;
			continue;
		}
		x+=first_bit(f);
		for(int g=h-1;g>=0;g--){
			x<<=6;
			x+=first_bit(a[g][x>>6]);
		}
		return x;
	}
	return -1;
}
inline int front(){return next(0);}
inline int back(){return prev(V-1);}
};
\end{lstlisting}
\par\vspace{1em}

\Needspace{10\baselineskip}
\subsection{\texttt{ds/FastODT.cpp}}
\textbf{位集加速的区间覆盖}

\TemplUsage{借助 FastSet 存储区间右端点。下标 1..n，适合值域固定的区间覆盖。}
\TemplApi{\begin{itemize}
\item 先粘贴 FastSet，再粘贴本文件。
\item \texttt{Fast\_ODT t; t.setN(n)} 分配空间；先 insert(1,n,c0) 明确初始颜色。
\item \texttt{insert(l,r,c)} 覆盖区间并返回旧的 \{L,R,col\} 段；\texttt{extract(l,r)} 获取当前片段。
\end{itemize}}
\textbf{额外依赖（按此顺序粘贴）：}\texttt{ds/FastSet.cpp}。

\textbf{最小示例（放入 main）：}
\begin{lstlisting}
Fast_ODT t; t.setN(5); t.insert(1,5,0);
t.insert(2,4,7); auto a=t.extract(2,3);
assert(a.size()==1 && a[0][2]==7);
\end{lstlisting}
\TemplNote{extract 会分裂边界。不保证任意操作序列的高效性，成本仍取决于访问的段数。}

\Needspace{6\baselineskip}\textbf{完整代码}
\begin{lstlisting}[firstnumber=1]
//
// template for Fast ODT, require Fast Set
// Chtholly is very cute!
// version 1.0 (Last Update Jun 16th, 2026)
//
// usage:
//   Fast_ODT fot; fot.setN(n); fot.insert(l,r,c); fot.extract(l,r);
//   requires FastSet.cpp
//
struct Fast_ODT{
private:
vector<int> a;int n=0;
FastSet S;
void cut(int y){
	if(S[y]) return;
	a[y]=a[S.next(y)];
	S.insert(y);
}
public:
void setN(int _n){assert(1<=_n&&_n<INT_MAX);n=_n;a.assign(n+1,0);S.setN(n);S.insert(n);a[n]=-1;}

vector<array<int,3>> extract(int l,int r){
	// extract segments from [l,r]
	// before you use this function, pay attention is the TC correct?
	assert(1<=l&&l<=r&&r<=n);
	cut(l-1);cut(r);
	vector<array<int,3>> I;
	for(int i=l,j;i<=r;){
		j=S.next(i);
		I.push_back({i,j,a[j]});
		i=j+1;
	}
	return I;
}
vector<array<int,3>> insert(int l,int r,int c){
	// erase segments from [l,r]
	// return vector for segs, [l,r,c]
	assert(1<=l&&l<=r&&r<=n);
	cut(l-1);cut(r);
	vector<array<int,3>> I;
	for(int i=l,j;i<=r;){
		j=S.next(i);
		S.erase(j);
		I.push_back({i,j,a[j]});
		i=j+1;
	}
	a[r]=c;S.insert(r);
	return I;
}
};
\end{lstlisting}
\par\vspace{1em}

\Needspace{10\baselineskip}
\subsection{\texttt{ds/FHQ.cpp}}
\textbf{按键或按序列位置分裂的 Treap}

\TemplUsage{随机平衡二叉树。外部维护根编号 rt，空树为 0，普通 split/merge 期望 O(log n)。}
\TemplApi{\begin{itemize}
\item \texttt{FHQ t; t.set()} 清空；\texttt{newnode(x)} 建一个 key=x 的节点。
\item \texttt{split\_by\_key(rt,k,a,b)} 分为 key<=k 与 key>k；\texttt{split\_by\_size(rt,k,a,b)} 分出前 k 个。
\item \texttt{rt=merge(a,b)} 要求 a 所有元素应在 b 前面；按键维护时要求 max(a)<=min(b)。
\item \texttt{merge\_it(a,b)} 可合并键域交错的两树。\texttt{push(rt,v)} 给整个非空子树的 key 加 v。
\end{itemize}}
\textbf{最小示例（放入 main）：}
\begin{lstlisting}
FHQ t; t.set(); int rt=0;
for(int v:{2,5,8})rt=t.merge(rt,t.newnode(v));
int a,b; t.split_by_key(rt,5,a,b);
assert(t.siz[a]==2 && t.siz[b]==1);
rt=t.merge(a,b); assert(t.siz[rt]==3);
\end{lstlisting}
\TemplNote{setN(n) 只是创建 n 个孤立节点，不会把它们连成树。按键分裂依赖排序不变量；区间加可能破坏该不变量。没有现成的区间翻转或持久化。全局 rnd 可能与别的模板重名。}

\Needspace{6\baselineskip}\textbf{完整代码}
\begin{lstlisting}[firstnumber=1]
//
// template for FHQ Treap
//
// usage:
//   FHQ t; t.setN(n); split_by_size / split_by_key / merge / merge_it
//   t.push(id, x) for lazy add on key[]
//
mt19937_64 rnd(time(0));
struct FHQ{
int tot=0;
vector<int> ls{0},rs{0},key{0},siz{0},tag{0};
vector<ll> val{0};
inline void push(const int &id,const int &x){
	assert(1<=id&&id<=tot);
	key[id]+=x;tag[id]+=x;
}
inline void pushup(const int &id){
	siz[id]=siz[ls[id]]+siz[rs[id]]+1;
}
inline void pushdown(const int &id){
	if(tag[id]){
		if(ls[id]) push(ls[id],tag[id]);
		if(rs[id]) push(rs[id],tag[id]);
		tag[id]=0;
	}
}
void set(){setN(0);}
void setN(int _n){
	assert(0<=_n&&_n<INT_MAX);tot=_n;
	ls.assign(tot+1,0);rs.assign(tot+1,0);tag.assign(tot+1,0);
	key.assign(tot+1,0);siz.assign(tot+1,1);siz[0]=0;val.resize(tot+1);
	for(int i=1;i<=tot;i++)val[i]=rnd();
}
inline int newnode(int x){
	assert(tot<INT_MAX-1);++tot;
	ls.push_back(0);rs.push_back(0);tag.push_back(0);siz.push_back(1);
	key.push_back(x);val.push_back(rnd());return tot;
}
int leftmost(int x){
	while(ls[x]) pushdown(x),x=ls[x];
	return x;
}
int rightmost(int x){
	while(rs[x]) pushdown(x),x=rs[x];
	return x;
}
int merge(int p,int q){
	if(!p||!q) return p|q;
	if(val[p]>val[q]) {pushdown(p);rs[p]=merge(rs[p],q);pushup(p);return p;}
	else {pushdown(q);ls[q]=merge(p,ls[q]);pushup(q);return q;}
}
void split_by_key(int rt,int k,int &p,int &q){
	if(!rt){p=q=0;return;}
	pushdown(rt);
	if(k<key[rt]) q=rt,split_by_key(ls[rt],k,p,ls[rt]);
	else p=rt,split_by_key(rs[rt],k,rs[rt],q);
	pushup(rt);
}
void split_by_size(int rt,int k,int &p,int &q){
	assert(0<=k&&k<=siz[rt]);
	if(k==0){p=0;q=rt;return;}
	if(k==siz[rt]){p=rt;q=0;return;}
	pushdown(rt);
	if(k<=siz[ls[rt]]) q=rt,split_by_size(ls[rt],k,p,ls[rt]),pushup(rt);
	else p=rt,split_by_size(rs[rt],k-1-siz[ls[rt]],rs[rt],q),pushup(rt);
}
int merge_it(int p,int q){
	// ~O(logn*logV) to merge two set (maybe intersect)
	if(!p||!q) return p|q;
	if(key[rightmost(p)]<=key[leftmost(q)]) return merge(p,q);
	if(key[rightmost(q)]<=key[leftmost(p)]) return merge(q,p);
	if(val[p]<val[q]) swap(p,q);
	int r,s;
	pushdown(p);
	split_by_key(q,key[p],r,s);
	ls[p]=merge_it(ls[p],r);
	rs[p]=merge_it(rs[p],s);
	pushup(p);
	return p;
}
void alldown(int id){
	if(!id) return;
	pushdown(id);
	alldown(ls[id]);
	alldown(rs[id]);
}
};
\end{lstlisting}
\par\vspace{1em}

\Needspace{10\baselineskip}
\subsection{\texttt{ds/pq2stack.cpp}}
\textbf{把优先队列删除转为栈撤销}

\TemplUsage{适合只支持栈式加入/撤销，却需要删除最大权元素的离线或可撤销结构。模板只安排顺序，实际状态由调用方维护。}
\TemplApi{\begin{itemize}
\item \texttt{pq2stack<T> q; q.set()} 初始化；\texttt{push(x,w)} 登记元素及删除优先级。
\item 每次查询前先 \texttt{prep()}，按返回顺序把元素加入外部状态。
\item \texttt{pop()} 删除最大权元素，返回外部栈需要撤销的次数；之后再次 prep 会给出重新加入的元素。
\end{itemize}}
\textbf{最小示例（放入 main）：}
\begin{lstlisting}
pq2stack<int> q; q.set(); vector<int> st;
q.push(10,1); q.push(20,2);
for(int x:q.prep())st.push_back(x);
int k=q.pop(); while(k--)st.pop_back();
for(int x:q.prep())st.push_back(x);
assert(st.size()==1 && st[0]==10);
\end{lstlisting}
\TemplNote{不能把 pop 的返回值当成被删元素。每次重加入都需要重新产生撤销记录；同权元素删除顺序不固定。单次可能重排大量元素，收益依赖整段操作序列。}

\Needspace{6\baselineskip}\textbf{完整代码}
\begin{lstlisting}[firstnumber=1]
//
// template for  pq2stack (a.k.a. fuckGold14526 trick)
// version 1.0 (Last Update Jun 28th, 2026)
//
// usage:
//   pq2stack<T> pq; pq.push(x,w); pq.prep(); pq.pop();
//   pop removes max weight; needs chkmin; call prep() before each query
//
template<typename T>
struct pq2stack{
private:
struct info{
	T e;int w;bool tag;
	inline bool operator < (const info &x)const{return w<x.w;}
};
int n=0,r=0,d=0;
vector<info> I;
public:
void set(){n=r=d=0;I.clear();}
vector<T> prep(){
	vector<T> ans;
	while(d<I.size()) ans.push_back(I[d++].e);
	return ans;
}
void push(T x,int w){
	I.push_back({x,w,0});
}
void imerge(int i,int j,int k){
	if(i==j||j==k) return;
	static vector<info> y;
	y.resize(k-i);
	merge(I.begin()+i,I.begin()+j,I.begin()+j,I.begin()+k,y.begin());
	for(int t=i;t<k;t++) I[t]=y[t-i];
}
int pop(){
	// pq2stack core
	assert(!I.empty());
	int tar=(r?r&-r:inf<int>),mn=I.size()-1,hd=mn;
	for(int i=I.size()-1;i>=0;i--){
		if(I[i].tag){
			I[i].tag=false;
			tar--;
		}
		if(I[i]<I[mn]) mn=i;
		else hd=i;
		if(tar==0) break;
	}
	sort(I.begin()+n,I.end());
	int cut=I.size();
	for(int i=I.size()-1;i>hd;i--) if(I[i]<I[i-1]){
		imerge(i,cut,I.size());
		cut=i;
	}
	imerge(hd,cut,I.size());
	int ans=max(0,d-hd);
	int mx=(int)I.size()-1;
	for(int i=n;i<(int)I.size();i++) if(I[i].w>I[mx].w) mx=i;
	if(mx!=(int)I.size()-1) swap(I[mx],I.back());
	I.pop_back();
	chkmin(d,hd);
	n=I.size();
	if(r==0){
		for(int i=0;i<n;i++) I[i].tag=true;
		r=n;
	}else{
		for(int i=n-(r&-r)+1;i<n;i++) I[i].tag=true;
		r--;
	}
	return ans;
}
};
// !!!!!!!!!!!!!!! before each query, use pq.prep() !!!!!!!!!!!!!!!
\end{lstlisting}
\par\vspace{1em}

\Needspace{10\baselineskip}
\subsection{\texttt{ds/xihe\_tree.cpp}}
\textbf{析合树：返回完整结构}

\TemplUsage{对 1..n 的排列构建析合树。构建时间 O(n log n)，包含四毛子 RMQ、辅助线段树和返回结构在内的空间均为 O(n)。}
\TemplApi{先粘贴 \texttt{basic/ST table (linear).cpp}。\texttt{xihe\_tree tr; auto t=tr.build(a)} 接受不带哨兵的 vector；也可先 setN(n)，填写 tr.a[1..n]，再调用 build()。结果独立持有数据：rt 为根，cnt 为结点数，son[u] 为按位置排序的儿子；typ[u] 为 LEAF/INC/DEC/PRIME；[L[u],R[u]] 为下标闭区间，[mn[u],mx[u]] 为值域闭区间；id[i] 为第 i 项对应的叶子编号。}
\textbf{最小示例（放入 main）：}
\begin{lstlisting}
xihe_tree tr; auto t=tr.build(vector<int>{3,1,2});
int u=t.rt;
assert(t.L[u]==1 && t.R[u]==3);
assert(t.mn[u]==1 && t.mx[u]==3);
assert(t.typ[u]==xihe_tree::DEC && t.son[u].size()==2);
assert(t.mn[t.id[2]]==1);
\end{lstlisting}
\TemplNote{输入必须恰为 1..n 的排列，不接受重复值；必要时在外部离散化，但需确认题目是否只关心排名连续。返回的结点编号、下标区间、id 均从 1 开始，0 不使用。n=0 时 rt=cnt=0。INC/DEC 是儿子值域顺序递增/递减的合点，PRIME 是析点。不要假定父编号大于儿子编号；在返回结果上自行遍历、DP，构建器不含 solve。再次 build 或修改构建器不会改变已返回结果。}

\Needspace{6\baselineskip}\textbf{完整代码}
\begin{lstlisting}[firstnumber=1]
////////////////////////////////////////////////////////////////
//
// template for Divide-Combine Tree
//
// usage:
//   Paste basic/ST table (linear).cpp first.
//   xihe_tree tr; auto t=tr.build(vector<int>{3,1,2});
//   Or setN(n), fill a[1..n], then auto t=tr.build().
//   Input must be a permutation of 1..n. Node / interval indices are 1-based.
//   t.rt, t.son, t.typ, t.L/R, t.mn/mx, t.id; no problem-specific solver.
//   Build O(n log n), space O(n). Children are ordered by position.
//
////////////////////////////////////////////////////////////////
struct xihe_tree{
enum Type{LEAF,INC,DEC,PRIME};
struct result{
	int n=0,rt=0,cnt=0;
	vector<vector<int>> son;
	vector<Type> typ;
	vector<int> L,R,mn,mx,id;
};
int n=0;vector<int> a;
void set(){n=0;a.clear();}
void setN(int _n){assert(0<=_n&&_n<=INT_MAX/4);n=_n;a.assign(n+1,0);}
private:
struct segt{
	int n;vector<int> val,laz;
	segt(int _n):n(_n),val(4*n),laz(4*n){}
	void push(int id,int x){val[id]+=x;laz[id]+=x;}
	void pushdown(int id){if(laz[id]){push(id<<1,laz[id]);push(id<<1|1,laz[id]);laz[id]=0;}}
	void add(int ql,int qr,int x,int id,int l,int r){
		if(ql<=l&&r<=qr){push(id,x);return;}
		int mid=(l+r)>>1;pushdown(id);
		if(ql<=mid)add(ql,qr,x,id<<1,l,mid);
		if(mid<qr)add(ql,qr,x,id<<1|1,mid+1,r);
		val[id]=min(val[id<<1],val[id<<1|1]);
	}
	void add(int l,int r,int x){add(l,r,x,1,1,n);}
	int first(){
		int id=1,l=1,r=n;assert(val[1]<=0);
		while(l<r){
			int mid=(l+r)>>1;pushdown(id);
			if(val[id<<1]<=0){id<<=1;r=mid;}
			else{id=id<<1|1;l=mid+1;}
		}
		return l;
	}
};
public:
result build()const{
	assert(a.size()==(size_t)n+1||!n);
	result res;res.n=n;
	auto &son=res.son;auto &typ=res.typ;auto &L=res.L;auto &R=res.R;auto &id=res.id;
	int &cnt=res.cnt;
	son.resize(2*n+1);typ.resize(2*n+1,LEAF);L.resize(2*n+1);R.resize(2*n+1);id.resize(n+1);
	if(!n){res.mn.resize(1);res.mx.resize(1);return res;}
	vector<char> vis(n+1);
	for(int i=1;i<=n;i++){assert(1<=a[i]&&a[i]<=n&&!vis[a[i]]);vis[a[i]]=true;}
	vector<int> v(a.begin()+1,a.end());
	Linear_RMQ<int> mn;Linear_RMQ<int,true> mx;mn.build(v);mx.build(v);
	auto judge=[&](int l,int r){return mx.ask(l-1,r-1)-mn.ask(l-1,r-1)==r-l;};
	vector<int> M(2*n+1),st1(n+1),st2(n+1),st(n+1);
	int tp1=0,tp2=0,tp=0;segt T(n);
	for(int i=1;i<=n;i++){
		while(tp1&&a[i]<=a[st1[tp1]]) T.add(st1[tp1-1]+1,st1[tp1],a[st1[tp1]]),tp1--;
		while(tp2&&a[i]>=a[st2[tp2]]) T.add(st2[tp2-1]+1,st2[tp2],-a[st2[tp2]]),tp2--;
		T.add(st1[tp1]+1,i,-a[i]);
		T.add(st2[tp2]+1,i,a[i]);
		st1[++tp1]=st2[++tp2]=i;
		id[i]=++cnt;L[cnt]=R[cnt]=M[cnt]=i;typ[cnt]=LEAF;
		int le=T.first(),now=cnt;
		while(tp&&L[st[tp]]>=le){
			if((typ[st[tp]]==INC||typ[st[tp]]==DEC)&&judge(M[st[tp]],i)){
				// 当儿子
				R[st[tp]]=i;M[st[tp]]=L[now];son[st[tp]].push_back(now);now=st[tp--];
			}else if(judge(L[st[tp]],i)){
				// 新合点
				++cnt;
				if(a[L[st[tp]]]<a[i]) typ[cnt]=INC;
				else typ[cnt]=DEC;
				L[cnt]=L[st[tp]];R[cnt]=i;M[cnt]=L[now];
				son[cnt].push_back(st[tp]);
				son[cnt].push_back(now);
				now=cnt;tp--;
			}else{
				// 新析点
				++cnt;typ[cnt]=PRIME;
				son[cnt].push_back(now);
				do son[cnt].push_back(st[tp--]);
				while(tp&&!judge(L[st[tp]],i));
				L[cnt]=(tp?L[st[tp]]:L[now]);R[cnt]=i;
				if(tp) son[cnt].push_back(st[tp--]);
				reverse(son[cnt].begin(),son[cnt].end());
				now=cnt;
			}
		}
		st[++tp]=now;
		T.add(1,i,-1);
	}
	assert(tp==1);res.rt=st[1];
	son.resize(cnt+1);typ.resize(cnt+1);L.resize(cnt+1);R.resize(cnt+1);
	res.mn.resize(cnt+1);res.mx.resize(cnt+1);
	for(int i=1;i<=cnt;i++){res.mn[i]=mn.ask(L[i]-1,R[i]-1);res.mx[i]=mx.ask(L[i]-1,R[i]-1);}
	return res;
}
result build(const vector<int> &v){
	assert(v.size()<=(size_t)INT_MAX/4);setN(v.size());copy(v.begin(),v.end(),a.begin()+1);return build();
}
};
// end for ds/xihe_tree.cpp
/////////////////////////
\end{lstlisting}

\par\vspace{1em}

\clearpage
\section{图论}
\Needspace{10\baselineskip}
\subsection{\texttt{graph/flow.cpp}}
\textbf{Dinic 最大流}

\TemplUsage{有向网络最大流，容量是非负 ll。一般图最坏 O(V\textasciicircum{}2 E)，具体模型常有更好表现。}
\TemplApi{\begin{itemize}
\item \texttt{Flow\_Graph g; g.setN(n); g.setST(s,t)} 初始化，点编号 1..n，源汇必须不同。
\item \texttt{id=g.add\_edge(u,v,cap)} 只添加一条有向边，自动生成反向残量边。
\item \texttt{max\_flow()} 返回累计最大流值；\texttt{edge\_flow(id)} 查看返回的正向边上的流量。
\end{itemize}}
\textbf{最小示例（放入 main）：}
\begin{lstlisting}
Flow_Graph g; g.setN(3); g.setST(1,3);
int id=g.add_edge(1,2,5); g.add_edge(2,3,3);
assert(g.max_flow()==3 && g.edge_flow(id)==3);
assert(g.max_flow()==3);
\end{lstlisting}
\TemplNote{重复求流返回总流量，不是新增流量；添加边后继续增广。setN 清空图，new\_node 保留旧图添加点。无向容量边需要根据建模显式加边。流量总和应在 ll 范围内。}

\Needspace{6\baselineskip}\textbf{完整代码}
\begin{lstlisting}[firstnumber=1]
//
// template for max flow (dinic)
// last update Jun 23rd, 2026
//
// usage:
//   Flow_Graph g; g.setN(n); g.setST(S,T); g.add_edge(u,v,cap);
//   ll f = g.max_flow(); g.edge_flow(id);
//   setN clears the graph; new_node keeps existing edges
//
struct Flow_Graph{
private:
static const ll inf=2.02e18;
struct edge{int v,nxt;ll w;};
vector<edge> e=vector<edge>(2);bool is_flowed=false;
int n=0,m=1,S=0,T=0;ll ans=0;
vector<int> lst{0},level{0},cur{0};
void add(int u,int v,ll w){
	assert(m<INT_MAX-1);e.push_back({v,lst[u],w});lst[u]=++m;
}
bool BFS(){
	for(int i=0;i<=n;i++) level[i]=-1;
	queue<int> q;q.push(S);level[S]=0;
	while(!q.empty()){
		int u=q.front();q.pop();
		if(u==T) return true;
		for(int id=lst[u];id;id=e[id].nxt) if(e[id].w&&level[e[id].v]==-1){
			level[e[id].v]=level[u]+1;
			q.push(e[id].v);
		}
	}
	return false;
}
ll DFS(int u,ll flow){
	if(u==S) return flow;
	ll now=0;
	for(int &i=cur[u];i;i=e[i].nxt) if(level[e[i].v]==level[u]-1&&e[i^1].w){
		ll fo=DFS(e[i].v,min(e[i^1].w,flow-now));
		e[i^1].w-=fo;e[i].w+=fo;now+=fo;
		if(now==flow) break;
	}
	return now;
}
public:
void set(){
	n=S=T=0;m=1;ans=0;is_flowed=false;e.assign(2,{0,0,0});
	lst.assign(1,0);level.assign(1,0);cur.assign(1,0);
}
void setN(int _n){
	assert(0<=_n&&_n<INT_MAX);set();n=_n;
	lst.assign(n+1,0);level.assign(n+1,0);cur.assign(n+1,0);
}
void setST(int _S,int _T){
	assert(1<=min(_S,_T)&&max(_S,_T)<INT_MAX&&_S!=_T);
	if(S!=_S||T!=_T){
		for(int i=2;i<=m;i+=2)e[i].w+=e[i^1].w,e[i^1].w=0;
		ans=0;
	}
	S=_S;T=_T;n=max({n,S,T});
	lst.resize(n+1);level.resize(n+1);cur.resize(n+1);is_flowed=false;
}
int new_node(){
	assert(n<INT_MAX-1);++n;lst.push_back(0);level.push_back(0);cur.push_back(0);
	is_flowed=false;return n;
}
ll edge_flow(int id){
	assert(2<=id&&id<=m&&id%2==0);
	return e[id^1].w;
}
int add_edge(int u,int v,ll w){
	// return forward edge id
	assert(1<=min(u,v)&&max(u,v)<=n&&0<=w);
	assert(m&1);
	is_flowed=false;add(u,v,w);add(v,u,0);return m-1;
}
ll max_flow(){
	assert(1<=min(S,T)&&max(S,T)<=n&&S!=T);
	if(is_flowed) return ans;
	is_flowed=true;
	while(BFS()){
		for(int i=1;i<=n;i++) cur[i]=lst[i];
		ll f=DFS(T,inf);assert(ans<=LLONG_MAX-f);ans+=f;
	}
	return ans;
}
};
\end{lstlisting}
\par\vspace{1em}

\Needspace{10\baselineskip}
\subsection{\texttt{graph/mcmf.cpp}}
\textbf{最小费用最大流}

\TemplUsage{先求最大可行流量，再在最大流中最小化费用。有负费用边时先用 SPFA 求势能；之后 Dijkstra 求最短路，在零约化费用子图上用 Dinic 批量增广，适合大量路径费用相同的网络。DFS 使用显式数组，不依赖递归栈。}
\TemplApi{\begin{itemize}
\item \texttt{Cost\_Flow\_Graph g; g.setN(n); g.setST(s,t)} 初始化。
\item \texttt{add\_edge(u,v,cap,cost)} 添加有向边并返回正向边编号。
\item \texttt{min\_cost\_flow()} 返回 \{总流量,总费用\}；\texttt{edge\_flow(id)} 查询流量。
\end{itemize}}
\textbf{最小示例（放入 main）：}
\begin{lstlisting}
Cost_Flow_Graph g; g.setN(3); g.setST(1,3);
int id=g.add_edge(1,2,2,-1); g.add_edge(2,3,2,3);
auto [f,c]=g.min_cost_flow();
assert(f==2 && c==4 && g.edge_flow(id)==2);
\end{lstlisting}
\TemplNote{允许负费用边，但整个输入图不能有负费用环。没有指定流量参数；要限制流量，增加一个容量为上限的新源边。更改图后重新从原始容量计算；不改图则返回缓存。容量、总流量及计算过程中的总费用须能用 ll 表示，距离与势能使用 \texttt{\_\_int128}。相同最短路费用的增广会合并，但费用层数很多时仍可能较慢。}

\Needspace{6\baselineskip}\textbf{完整代码}
\begin{lstlisting}[firstnumber=1]
//
// template for min cost max flow (potential + dijkstra + blocking flow)
//
// usage:
//   Cost_Flow_Graph g; g.setN(n); g.setST(S,T); g.add_edge(u,v,cap,cost);
//   auto [flow,cost]=g.min_cost_flow(); g.edge_flow(id);
//   negative costs allowed, input graph must have no negative cost cycle
//   setN clears the graph; editing the graph makes the next solve start over
//   flow and cost must fit ll; distances use __int128
//
struct Cost_Flow_Graph{
private:
static constexpr __int128 inf=(__int128)1<<120;
struct edge{int v,nxt;ll w,c;};
vector<edge> e=vector<edge>(2);bool is_flowed=false;
int n=0,m=1,S=0,T=0;ll ans=0,cost=0;
vector<int> lst{0},pre{0},dep,cur,q;
vector<ll> lim,rem;
vector<__int128> h{0},dis{0};
void add(int u,int v,ll w,ll c){
	assert(m<INT_MAX-1);e.push_back({v,lst[u],w,c});lst[u]=++m;
}
void SPFA(){
	fill(h.begin(),h.end(),0);
	bool neg=false;
	for(int i=2;i<=m;i++)if(e[i].w&&e[i].c<0){neg=true;break;}
	if(!neg)return;
	vector<int> cnt(n+1);vector<char> inq(n+1,true);
	queue<int> q;for(int i=1;i<=n;i++)q.push(i);
	while(!q.empty()){
		int u=q.front();q.pop();inq[u]=false;
		for(int id=lst[u];id;id=e[id].nxt)if(e[id].w){
			int v=e[id].v;
			if(h[v]<=h[u]+e[id].c)continue;
			h[v]=h[u]+e[id].c;cnt[v]=cnt[u]+1;
			assert(cnt[v]<n);
			if(!inq[v])inq[v]=true,q.push(v);
		}
	}
}
bool Dijkstra(){
	fill(dis.begin(),dis.end(),inf);dis[S]=0;
	priority_queue<pair<__int128,int>,vector<pair<__int128,int>>,greater<pair<__int128,int>>> q;
	q.push({0,S});
	while(!q.empty()){
		auto [d,u]=q.top();q.pop();
		if(d!=dis[u])continue;
		if(u==T)break;
		for(int id=lst[u];id;id=e[id].nxt)if(e[id].w){
			int v=e[id].v;__int128 w=e[id].c+h[u]-h[v];
			assert(w>=0);
			if(dis[v]<=d+w)continue;
			dis[v]=d+w;q.push({dis[v],v});
		}
	}
	if(dis[T]==inf)return false;
	for(int i=1;i<=n;i++)h[i]+=min(dis[i],dis[T]);
	return true;
}
bool BFS(){
	fill(dep.begin(),dep.end(),-1);dep[S]=0;
	int l=0,r=0;q[r++]=S;
	while(l<r){
		int u=q[l++];
		for(int id=lst[u];id;id=e[id].nxt){
			int v=e[id].v;
			if(!e[id].w||dep[v]!=-1||h[v]!=h[u]+e[id].c)continue;
			dep[v]=dep[u]+1;q[r++]=v;
		}
	}
	return dep[T]!=-1;
}
ll DFS(){
	// iterative DFS: rem[u] is the part of lim[u] not sent yet
	int u=S;lim[u]=rem[u]=LLONG_MAX;
	while(true){
		if(u==T)rem[u]=0;
		int &id=cur[u];
		while(rem[u]&&id){
			int v=e[id].v;
			if(e[id].w&&dep[v]==dep[u]+1&&h[v]==h[u]+e[id].c)break;
			id=e[id].nxt;
		}
		if(!rem[u]||!id){
			ll f=lim[u]-rem[u];
			if(u==S)return f;
			if(rem[u])dep[u]=-1;
			int x=pre[u];u=e[x^1].v;
			e[x].w-=f;e[x^1].w+=f;rem[u]-=f;
			if(rem[u])cur[u]=e[x].nxt;
		}else{
			int v=e[id].v;pre[v]=id;
			lim[v]=rem[v]=min(rem[u],e[id].w);u=v;
		}
	}
}
public:
void set(){
	n=S=T=0;m=1;ans=cost=0;is_flowed=false;e.assign(2,{0,0,0,0});
	dep.clear();cur.clear();q.clear();lim.clear();rem.clear();
	lst.assign(1,0);pre.assign(1,0);h.assign(1,0);dis.assign(1,0);
}
void setN(int _n){
	assert(0<=_n&&_n<INT_MAX);set();n=_n;
	lst.assign(n+1,0);pre.assign(n+1,0);h.assign(n+1,0);dis.assign(n+1,0);
}
void setST(int _S,int _T){
	assert(1<=min(_S,_T)&&max(_S,_T)<INT_MAX&&_S!=_T);
	S=_S;T=_T;n=max({n,S,T});
	lst.resize(n+1);pre.resize(n+1);h.resize(n+1);dis.resize(n+1);is_flowed=false;
}
int new_node(){
	assert(n<INT_MAX-1);++n;lst.push_back(0);pre.push_back(0);h.push_back(0);dis.push_back(0);
	is_flowed=false;return n;
}
ll edge_flow(int id){
	assert(2<=id&&id<=m&&id%2==0);
	return e[id^1].w;
}
int add_edge(int u,int v,ll w,ll c){
	// return forward edge id
	assert(1<=min(u,v)&&max(u,v)<=n&&w>=0&&c!=LLONG_MIN);
	assert(m&1);
	is_flowed=false;add(u,v,w,c);add(v,u,0,-c);return m-1;
}
pair<ll,ll> min_cost_flow(){
	assert(1<=min(S,T)&&max(S,T)<=n&&S!=T);
	if(is_flowed)return {ans,cost};
	for(int i=2;i<=m;i+=2)e[i].w+=e[i^1].w,e[i^1].w=0;
	ans=cost=0;SPFA();
	dep.resize(n+1);cur.resize(n+1);q.resize(n+1);
	lim.resize(n+1);rem.resize(n+1);
	while(Dijkstra())while(BFS()){
		cur=lst;ll f=DFS();assert(f>0);
		__int128 c=h[T]-h[S];
		assert(ans<=LLONG_MAX-f);
		assert(c>=((__int128)LLONG_MIN-cost)/f&&c<=((__int128)LLONG_MAX-cost)/f);
		ans+=f;cost=(ll)(cost+c*f);
	}
	is_flowed=true;return {ans,cost};
}
};
\end{lstlisting}
\par\vspace{1em}

\Needspace{10\baselineskip}
\subsection{\texttt{graph/sssp.cpp}}
\textbf{非负权单源最短路}

\TemplUsage{堆优化 Dijkstra，有向图，复杂度 O((n+m) log n)。}
\TemplApi{\begin{itemize}
\item \texttt{graph\_SSSP g; g.setN(n); g.setS(s)} 设置图和源点。
\item \texttt{add\_edge(u,v,w)} 要求 w>=0；无向图要加两个方向。
\item \texttt{SSSP()} 返回长度 n+1 的距离数组，0 号不用，不可达为 10\textasciicircum{}18。
\end{itemize}}
\textbf{最小示例（放入 main）：}
\begin{lstlisting}
graph_SSSP g; g.setN(3); g.setS(1);
g.add_edge(1,2,4); g.add_edge(2,3,6);
auto d=g.SSSP(); assert(d[3]==10);
\end{lstlisting}
\TemplNote{有限最短路必须小于不可达哨兵 10\textasciicircum{}18，边权相加不能溢出。不支持负权；不会返回具体路径。}

\Needspace{6\baselineskip}\textbf{完整代码}
\begin{lstlisting}[firstnumber=1]
//
// template for Single-source shortest path
// last update Jun 23rd, 2026
//
// usage:
//   graph_SSSP G; G.setN(n); G.setS(s); G.add_edge(u,v,w);
//   auto d = G.SSSP();  // unreachable 1e18, non-negative weights
//
struct graph_SSSP{
private:
int n=0,s=0;
vector<vector<pair<int,ll>>> edg;
public:
void setN(int N){assert(0<=N&&N<INT_MAX);n=N;s=0;edg.assign(n+1,{});}
void setS(int S){s=S;assert(1<=s&&s<=n);}
void add_edge(int u,int v,ll w){
	assert(1<=min(u,v)&&max(u,v)<=n&&0<=w);
	edg[u].push_back(make_pair(v,w));
}
vector<ll> SSSP(){
	// return a vector of length n+1 --- ans[i] shortest path between s and i
	// if no path, ans[i] will be 1e18
	assert(1<=s&&s<=n);
	priority_queue<pair<ll,int>,vector<pair<ll,int>>,greater<pair<ll,int>>> q;
	vector<ll> d(n+1,1e18);
	d[s]=0;q.push(make_pair(0,s));
	while(!q.empty()){
		auto [dd,u]=q.top();q.pop();
		if(dd!=d[u]) continue;
		for(auto [v,w]:edg[u]) if(w<d[v]-d[u]){
			d[v]=d[u]+w;
			q.push(make_pair(d[v],v));
		}
	}
	return d;
}
};
\end{lstlisting}
\par\vspace{1em}

\Needspace{10\baselineskip}
\subsection{\texttt{graph/2-sat.cpp}}
\textbf{布尔约束求解}

\TemplUsage{求满足所有二元或子句的一组赋值，时间与空间 O(n+m)。}
\TemplApi{\begin{itemize}
\item \texttt{two\_sat s; s.setN(n)} 初始化 n 个布尔变量，下标 1..n。
\item \texttt{add(x,a,y,b)} 表示 (变量 x 等于 a) 或 (变量 y 等于 b)，a,b 为 0/1。
\item \texttt{solve()} 无解返回空 vector，否则 ans[1..n] 为可行赋值。
\end{itemize}}
\textbf{最小示例（放入 main）：}
\begin{lstlisting}
two_sat s; s.setN(2);
s.add(1,1,1,1); // force x1=true
s.add(1,0,2,1); // x1 implies x2
auto ans=s.solve();
assert(!ans.empty() && ans[1] && ans[2]);
\end{lstlisting}
\TemplNote{只有“或”子句接口；例如强制 x=a 可重复同一文字。setN 清空约束。不要把 add 当作直接添加两条蕴含边。}

\Needspace{6\baselineskip}\textbf{完整代码}
\begin{lstlisting}[firstnumber=1]
//
// template for 2-SAT
//
// usage:
//   two_sat sat; sat.set(); sat.setN(n); sat.add(x1,o1,x2,o2);
//   auto ans = sat.solve();  // empty if unsat; var i: 2i-1=T, 2i=F
//   setN clears old clauses
//
struct two_sat{
int tot=0,scnt=0,top=0,n=0;
vector<int> low,dfn,stk,bel;
vector<vector<int>> edg;vector<char> inq;
void set(){setN(0);}
void setN(int _n){
	// 1-index, 2i-1 true 2i false
	assert(0<=_n&&_n<(INT_MAX-2)/2);n=_n;
	low.assign(2*n+2,0);dfn.assign(2*n+2,0);stk.assign(2*n+2,0);
	bel.assign(2*n+2,0);inq.assign(2*n+2,0);edg.assign(2*n+2,{});
	tot=scnt=top=0;
}
void tarjan(int u){
	vector<pair<int,int>> q{{u,0}};
	low[u]=dfn[u]=++tot;stk[++top]=u;inq[u]=true;
	while(!q.empty()){
		auto &[x,i]=q.back();
		if(i<(int)edg[x].size()){
			int v=edg[x][i++];
			if(!dfn[v]){
				low[v]=dfn[v]=++tot;stk[++top]=v;inq[v]=true;q.push_back({v,0});
			}else if(inq[v])low[x]=min(low[x],dfn[v]);
		}else{
			int v=x;q.pop_back();
			if(!q.empty()){int p=q.back().first;low[p]=min(low[p],low[v]);}
			if(low[v]==dfn[v]){
				scnt++;int y;
				do{y=stk[top--];inq[y]=false;bel[y]=scnt;}while(y!=v);
			}
		}
	}
}
vector<int> solve(){
	// if impossible return empty
	// else return a vector with length n+1, the a[i] - any valid answer
	tot=scnt=top=0;
	for(int i=1;i<=2*n;i++) low[i]=dfn[i]=bel[i]=inq[i]=0;
	for(int i=1;i<=2*n;i++) if(!dfn[i]) tarjan(i); 
	for(int i=1;i<=n;i++) if(bel[2*i-1]==bel[2*i]) return {};
	vector<int> ans(n+1);
	for(int i=1;i<=n;i++) ans[i]=bel[2*i-1]<bel[2*i];
	return ans;
}
void add(int x1,int o1,int x2,int o2){
	// x1 = o1 or x2 = o2
	assert(1<=min(x1,x2)&&max(x1,x2)<=n&&0<=min(o1,o2)&&max(o1,o2)<=1);
	edg[2*x1-(!o1)].push_back(2*x2-o2);
	edg[2*x2-(!o2)].push_back(2*x1-o1);
}
};
\end{lstlisting}
\par\vspace{1em}

\Needspace{10\baselineskip}
\subsection{\texttt{graph/block-cut-tree.cpp}}
\textbf{点双连通分量与圆方树}

\TemplUsage{无向图的点双分解，支持非连通图和重边，不支持自环。构建时间 O(n+m)。}
\TemplApi{\begin{itemize}
\item \texttt{block\_cut g; g.setN(n)} 后 add\_edge，再 build。
\item 原点仍为 1..n，方点为 n+1..scc，tr 是圆方森林邻接表。
\item 第 k 个方点 n+k 的原图边存在 ec[k] 中。ff、dep 是森林上的父亲和深度。
\end{itemize}}
\textbf{最小示例（放入 main）：}
\begin{lstlisting}
block_cut g; g.setN(3);
g.add_edge(1,2); g.add_edge(2,3); g.build();
assert(g.scc==5 && g.tr[2].size()==2);
\end{lstlisting}
\TemplNote{scc 是最终最大节点编号，不是分量个数，方点数应为 scc-n。割点可以属于多个点双，不能像 SCC 一样给每个原点分配唯一块号。递归深度可能达到 O(n)。}

\Needspace{6\baselineskip}\textbf{完整代码}
\begin{lstlisting}[firstnumber=1]
//
// template for Block Cut Tree
// use add_edge to add (u,v), tr[n+i] for vertex in i-th component
// ec[i] for edge in i-th component
// version 1.0 (Last Update Jun 17th, 2026)
//
// usage:
//   block_cut bc; bc.setN(n); bc.add_edge(u,v); bc.build();
//   vertex u -> tree u; BCC i -> tree n+i; bc.ec[i] edges
//
struct block_cut{
int n=0,tot=0,scc=0,top=0;
vector<char> inq;
vector<int> low,dfn,stk,dep,ff;
vector<vector<int>> edg,tr;
vector<vector<array<int,2>>> ec;
void setN(int _n){
	assert(0<=_n&&_n<(INT_MAX-2)/2);n=_n;
	low.assign(n+2,0);dfn.assign(n+2,0);stk.assign(n+2,0);inq.assign(n+2,0);
	edg.assign(n+2,{});tr.assign(2*n+2,{});ec.assign(n+2,{});
	dep.assign(2*n+2,0);ff.assign(2*n+2,0);tot=scc=top=0;
}
void add_edge(int u,int v){
	assert(1<=min(u,v)&&max(u,v)<=n&&u!=v);
	edg[u].push_back(v);
	edg[v].push_back(u);
}
void tarjan(int u,int fa=0){
	// frame: vertex, parent, next edge, skipped parent edge
	vector<array<int,4>> q{{u,fa,0,0}};
	low[u]=dfn[u]=++tot;stk[++top]=u;inq[u]=true;
	while(!q.empty()){
		auto &[x,p,i,skip]=q.back();
		if(i<(int)edg[x].size()){
			int v=edg[x][i++];
			if(v==p&&!skip){skip=1;continue;}
			if(!dfn[v]){
				low[v]=dfn[v]=++tot;stk[++top]=v;inq[v]=true;q.push_back({v,x,0,0});
			}else if(inq[v])low[x]=min(low[x],dfn[v]);
		}else{
			int v=x,parent=p;q.pop_back();
			if(q.empty())continue;
			low[parent]=min(low[parent],low[v]);
			if(low[v]>=dfn[parent]){
				++scc;tr[scc].push_back(parent);tr[parent].push_back(scc);int y;
				do{y=stk[top--];inq[y]=false;tr[y].push_back(scc);tr[scc].push_back(y);}while(y!=v);
			}
		}
	}
}
void dfs(int u,int fa){
	vector<int> q{u};ff[u]=fa;dep[u]=dep[fa]+1;
	for(int i=0;i<(int)q.size();i++){
		int x=q[i];
		for(int v:tr[x])if(v!=ff[x]){ff[v]=x;dep[v]=dep[x]+1;q.push_back(v);}
	}
}
void build(){
	fill(dfn.begin(),dfn.end(),0);fill(inq.begin(),inq.end(),0);
	fill(dep.begin(),dep.end(),0);fill(ff.begin(),ff.end(),0);
	for(auto &v:tr)v.clear();for(auto &v:ec)v.clear();
	tot=top=0;scc=n;
	for(int i=1;i<=n;i++)if(!dfn[i]){
		tarjan(i);dfs(i,0);inq[i]=false;top=0;
	}
	for(int i=1;i<=n;i++)for(int j:edg[i])if(i<j){
		int x=ff[i]==ff[j]?ff[i]:(dep[i]>dep[j]?ff[i]:ff[j]);
		assert(n<x&&x<=scc);ec[x-n].push_back({i,j});
	}
}
};
\end{lstlisting}
\par\vspace{1em}

\Needspace{10\baselineskip}
\subsection{\texttt{graph/Edge BCC.cpp}}
\textbf{边双连通分量及桥森林}

\TemplUsage{无向图删除桥后缩点，支持重边和非连通图。时间 O(n+m)。}
\TemplApi{\begin{itemize}
\item \texttt{bi\_edge g; g.setN(n)} 后加边。
\item \texttt{auto e=g.build()} 返回连接不同分量的边表；每端是分量编号。
\item \texttt{bel[u]} 表示点所属分量，分量编号 1..scnt。
\end{itemize}}
\textbf{最小示例（放入 main）：}
\begin{lstlisting}
bi_edge g; g.setN(3);
g.add_edge(1,2); g.add_edge(1,2); g.add_edge(2,3);
auto e=g.build();
assert(g.bel[1]==g.bel[2] && g.bel[2]!=g.bel[3]);
assert(e.size()==1);
\end{lstlisting}
\TemplNote{桥森林的顶点是分量，不是原图点。与圆方树分别处理边双、点双，不能互换。}

\Needspace{6\baselineskip}\textbf{完整代码}
\begin{lstlisting}[firstnumber=1]
//
// template for Edge BCC (edge-biconnected components)
//
// usage:
//   bi_edge b; b.setN(n); b.add_edge(u,v); auto E = b.build();
//   bel[u] is EBCC id; E lists edges between different blocks
//
struct bi_edge{
int n=0,top=0,tot=0,scnt=0;
vector<vector<int>> edg;
vector<int> low,dfn,stk,bel;
void setN(int _n){
	assert(0<=_n&&_n<INT_MAX);n=_n;edg.assign(n+1,{});
	low.assign(n+1,0);dfn.assign(n+1,0);stk.assign(n+1,0);bel.assign(n+1,0);
	top=tot=scnt=0;
}
void add_edge(int u,int v){
	assert(1<=min(u,v)&&max(u,v)<=n);
	edg[u].push_back(v);
	edg[v].push_back(u);
}
vector<array<int,2>> build(){
	fill(dfn.begin(),dfn.end(),0);top=tot=scnt=0;
	for(int u=1;u<=n;u++)if(!dfn[u]){
		vector<array<int,4>> q{{u,0,0,0}};
		low[u]=dfn[u]=++tot;stk[++top]=u;
		while(!q.empty()){
			auto &[x,p,i,skip]=q.back();
			if(i<(int)edg[x].size()){
				int v=edg[x][i++];
				if(v==p&&!skip){skip=1;continue;}
				if(!dfn[v]){low[v]=dfn[v]=++tot;stk[++top]=v;q.push_back({v,x,0,0});}
				else low[x]=min(low[x],dfn[v]);
			}else{
				int v=x,parent=p;q.pop_back();
				if(parent)low[parent]=min(low[parent],low[v]);
				if(low[v]==dfn[v]){
					scnt++;int y;do{y=stk[top--];bel[y]=scnt;}while(y!=v);
				}
			}
		}
	}
	vector<array<int,2>> E;
	for(int i=1;i<=n;i++) for(int j:edg[i]) if(bel[i]<bel[j])
		E.push_back({bel[i],bel[j]});
	return E;
}
};
\end{lstlisting}
\par\vspace{1em}

\Needspace{10\baselineskip}
\subsection{\texttt{graph/blossom.cpp}}
\textbf{无权一般图最大匹配}

\TemplUsage{带花树，允许奇环，返回匹配边数，最坏 O(n\textasciicircum{}3)。}
\TemplApi{\begin{itemize}
\item \texttt{blossom g; g.setN(n)} 后 add\_edge 添加无向边。
\item \texttt{solve()} 返回最大匹配的边数；\texttt{match[u]} 为匹配点，未匹配为 0。
\item \texttt{solve()*2==n} 判定存在完全匹配。
\end{itemize}}
\textbf{最小示例（放入 main）：}
\begin{lstlisting}
blossom g; g.setN(4);
g.add_edge(1,2); g.add_edge(2,3);
g.add_edge(3,1); g.add_edge(3,4);
assert(g.solve()==2);
assert(g.match[g.match[1]]==1);
\end{lstlisting}
\TemplNote{这是无权匹配，不会最小化或最大化边权。点数为奇数时不可能完全匹配。重复 solve 会重新求解。}

\Needspace{6\baselineskip}\textbf{完整代码}
\begin{lstlisting}[firstnumber=1]
//
// template for General Graph Matching (blossom), O(n^3)
//
// usage:
//   blossom g; g.setN(n); g.add_edge(u,v); int cnt=g.solve();
//   cnt*2==n if perfect; match[u]=0 if unmatched; 1-index
//
struct blossom{
int n=0;
vector<vector<int>> edg;
vector<int> match,pre,bel,q;
vector<char> inq,vis,flower;
void setN(int _n){
	assert(0<=_n&&_n<INT_MAX);n=_n;
	edg.assign(n+1,{});match.assign(n+1,0);
	pre.resize(n+1);bel.resize(n+1);q.resize(n+1);
	inq.resize(n+1);vis.resize(n+1);flower.resize(n+1);
}
void add_edge(int u,int v){
	assert(1<=min(u,v)&&max(u,v)<=n);
	if(u==v)return;
	edg[u].push_back(v);edg[v].push_back(u);
}
int lca(int u,int v){
	fill(vis.begin(),vis.end(),false);
	while(1){
		u=bel[u];vis[u]=true;
		if(!match[u])break;
		u=pre[match[u]];
	}
	while(!vis[bel[v]])v=pre[match[bel[v]]];
	return bel[v];
}
void shrink(int u,int v,int w){
	while(bel[u]!=w){
		flower[bel[u]]=flower[bel[match[u]]]=true;
		pre[u]=v;v=match[u];u=pre[v];
	}
}
bool BFS(int s){
	fill(inq.begin(),inq.end(),false);
	fill(pre.begin(),pre.end(),0);
	iota(bel.begin(),bel.end(),0);
	int h=0,t=0;q[t++]=s;inq[s]=true;
	while(h<t){
		int u=q[h++];
		for(int v:edg[u]){
			if(bel[u]==bel[v]||match[u]==v)continue;
			if(v==s||(match[v]&&pre[match[v]])){
				int w=lca(u,v);
				fill(flower.begin(),flower.end(),false);
				shrink(u,v,w);shrink(v,u,w);
				for(int i=1;i<=n;i++)if(flower[bel[i]]){
					bel[i]=w;
					if(!inq[i])inq[i]=true,q[t++]=i;
				}
			}else if(!pre[v]){
				pre[v]=u;
				if(!match[v]){
					while(v){
						int x=pre[v],y=match[x];
						match[v]=x;match[x]=v;v=y;
					}
					return true;
				}
				int x=match[v];inq[x]=true;q[t++]=x;
			}
		}
	}
	return false;
}
int solve(){
	fill(match.begin(),match.end(),0);
	int ans=0;
	for(int i=1;i<=n;i++)if(!match[i])ans+=BFS(i);
	return ans;
}
};
\end{lstlisting}
\par\vspace{1em}

\Needspace{10\baselineskip}
\subsection{\texttt{graph/bipartite-matching.cpp}}
\textbf{二分图最大匹配}

\TemplUsage{给定左右两部，多次加边后求最大基数匹配，返回匹配的点对。先贪心匹配，再用多源 BFS 记录起点与前驱，遇到未匹配右点就立即增广；整个过程不递归。}
\TemplApi{\begin{itemize}
\item \texttt{Bipartite\_Matching g; g.setN(n,m)}，左部下标 1..n，右部下标 1..m；两部编号互不影响。
\item \texttt{g.add\_edge(u,v)} 加入左点 u 到右点 v 的边；允许重边。
\item \texttt{auto a=g.solve()} 返回 \texttt{vector<pair<int,int>>}，每个二元组 (u,v) 是一对匹配的左点和右点，按左点递增；\texttt{a.size()} 就是最大匹配数。
\item 求解后的 \texttt{matchL[u]}、\texttt{matchR[v]} 是对应另一部的点，未匹配为 0。\texttt{set()} 清空，setN 清空并设置两部大小。
\end{itemize}}
\textbf{最小示例（放入 main）：}
\begin{lstlisting}
Bipartite_Matching g; g.setN(3,3);
g.add_edge(1,1); g.add_edge(2,1);
g.add_edge(2,2); g.add_edge(3,2);
auto a=g.solve(); assert(a.size()==2);
g.add_edge(3,3); a=g.solve(); assert(a.size()==3);
for(auto [u,v]:a){
    assert(g.matchL[u]==v && g.matchR[v]==u);
}
\end{lstlisting}
\TemplNote{邻接表在 solve 内按左点压成连续存储，保持各点的加边顺序。rt[u] 记录搜索树的未匹配起点，pre[u] 记录增广前驱；某起点完成增广后，跳过该搜索树剩余队列项。一轮允许处理不同长度的增广路。设总点数 V=n+m、边数 E，空间 $O(V+E)$，每轮 $O(n+E)$；保守最坏时间界为 $O((V+E)(\min(n,m)+1))$，不沿用分层 HK 的复杂度保证。空部、空边集允许，返回空匹配。每次 solve 都重新计算；可继续加边再调用。}

\Needspace{6\baselineskip}\textbf{完整代码}
\begin{lstlisting}[firstnumber=1]
////////////////////////////////////////////////////////////////
//
// template for Bipartite Matching
// usage:
//   Bipartite_Matching g; g.setN(n,m); g.add_edge(u,v);
//   auto a=g.solve(); // pairs {u,v}, ordered by the left endpoint
//   Left 1..n, right 1..m; a.size() is the maximum matching size.
//   matchL / matchR are partners, 0 if unmatched; parallel edges are allowed.
//   Greedy initialization + multi-source BFS augmentation; solve rebuilds.
//   O(n+m+E) space; O(n+E) per BFS round.
//
////////////////////////////////////////////////////////////////
struct Bipartite_Matching{
struct edge{int u,v;};
int n=0,m=0;
vector<edge> e;
vector<int> matchL,matchR;
void set(){n=m=0;e.clear();matchL.clear();matchR.clear();}
void setN(int _n,int _m){
	assert(0<=_n&&_n<INT_MAX-1&&0<=_m&&_m<INT_MAX);set();n=_n;m=_m;
	matchL.assign(n+1,0);matchR.assign(m+1,0);
}
void add_edge(int u,int v){
	assert(1<=u&&u<=n&&1<=v&&v<=m&&e.size()<INT_MAX);
	e.push_back({u,v});
}
vector<pair<int,int>> solve(){
	matchL.assign(n+1,0);matchR.assign(m+1,0);
	vector<int> head(n+2),to(e.size()),rt(n+1),pre(n+1),q;
	for(auto [u,v]:e)head[u+1]++;
	partial_sum(head.begin(),head.end(),head.begin());
	copy(head.begin(),head.begin()+n+1,rt.begin());
	for(auto [u,v]:e)to[rt[u]++]=v;
	q.reserve(n);
	for(int u=1;u<=n;u++)for(int i=head[u];i<head[u+1];i++)if(!matchR[to[i]]){
		matchL[u]=to[i];matchR[to[i]]=u;break;
	}
	while(true){
		fill(rt.begin(),rt.end(),0);q.clear();bool found=false;
		for(int u=1;u<=n;u++)if(!matchL[u]&&head[u]!=head[u+1]){
			rt[u]=u;pre[u]=0;q.push_back(u);
		}
		for(int h=0;h<(int)q.size();h++){
			int u=q[h];if(matchL[rt[u]])continue;
			for(int i=head[u];i<head[u+1];i++){
				int v=to[i],w=matchR[v];
				if(!w){
					found=true;
					for(int x=u;v;x=pre[x]){
						int t=matchL[x];matchL[x]=v;matchR[v]=x;v=t;
					}
					break;
				}
				if(!rt[w])rt[w]=rt[u],pre[w]=u,q.push_back(w);
			}
		}
		if(!found)break;
	}
	vector<pair<int,int>> ans;ans.reserve(min(n,m));
	for(int u=1;u<=n;u++)if(matchL[u])ans.emplace_back(u,matchL[u]);
	return ans;
}
};
// end for graph/bipartite-matching.cpp
/////////////////////////
\end{lstlisting}

\Needspace{10\baselineskip}
\subsection{\texttt{graph/regular-bipartite.cpp}}
\textbf{正则二分图的完美匹配分解}

\TemplUsage{两部各 n 个点，每个点度数均为 d，将全部 nd 条边恰好分成 d 组完美匹配。支持多重边，同一对点之间的边按重数计入分解。}
\TemplApi{\begin{itemize}
\item \texttt{Regular\_Bipartite g; g.setN(n)}；\texttt{g.add\_edge(u,v)} 加入左点 u 到右点 v 的边，两部下标均为 1..n。
\item \texttt{auto a=g.solve()} 返回 \texttt{vector<vector<int>>}，共 d 组。每组长度 n+1，\texttt{a[c][u]} 直接是第 c 组中左点 u 匹配的右点；\texttt{a[c][0]=0} 占位。
\item 每组的 a[c][1..n] 都是右部点的一个排列。所有组中 (u,v) 出现的总次数等于该边的输入重数。组编号为 0..d-1，组间没有固定的字典序要求。
\item d 从输入度数推得，solve 用 assert 检查正则条件；d=0 返回空分解。n=0 也返回空分解。
\end{itemize}}
\textbf{最小示例（放入 main）：}
\begin{lstlisting}
Regular_Bipartite g; g.setN(2);
g.add_edge(1,1); g.add_edge(1,2);
g.add_edge(2,1); g.add_edge(2,2);
auto a=g.solve(); assert(a.size()==2);
for(auto &match:a){
    assert(match.size()==3 && match[0]==0);
    assert(1<=match[1] && match[1]<=2);
    assert(1<=match[2] && match[2]<=2);
    assert(match[1]!=match[2]);
}
\end{lstlisting}
\TemplNote{确定性算法，时间 $O(n+nd\log(nd+1))$，包括结果在内空间 $O(n+nd)$。偶数度用欧拉分边；奇数度通过整数重数缩减找完美匹配，再借用已分解的匹配，把另一分支凑成二的幂次度数。无随机失败概率。非正则图不属于此接口，不要关闭断言后直接输入；重复 solve 不修改原边集。}

\Needspace{6\baselineskip}\textbf{完整代码}
\begin{lstlisting}[firstnumber=1]
////////////////////////////////////////////////////////////////
//
// template for Regular Bipartite Matching
// usage:
//   Regular_Bipartite g; g.setN(n); g.add_edge(u,v);
//   auto a=g.solve(); // d perfect matchings, a[c][u] is the right partner
//   Each row has n+1 entries, a[c][0]=0; c is 0-indexed, u is 1-indexed.
//   Both sides are 1..n; every vertex must have the same degree d.
//   Parallel edges are allowed and counted with multiplicity.
//   Deterministic O(n+nd*log(nd+1)) time, O(n+nd) space.
//
////////////////////////////////////////////////////////////////
struct Regular_Bipartite{
struct edge{int u,v;};
int n=0;
vector<edge> e=vector<edge>(1);
void set(){n=0;e.assign(1,{});}
void setN(int _n){assert(0<=_n&&_n<INT_MAX/4);set();n=_n;}
void add_edge(int u,int v){
	assert(1<=min(u,v)&&max(u,v)<=n&&e.size()+n<INT_MAX/2);
	e.push_back({u,v});
}
private:
// Split an even-degree bipartite graph into two equal-degree subgraphs.
// Negative id -u denotes a temporary edge (u,u).
vector<char> split(const vector<int> &a)const{
	assert(a.size()<INT_MAX/2);int m=a.size();
	vector<int> head(2*n+1,-1),nxt(2*m);
	vector<char> vis(m),col(m);
	for(int i=0;i<m;i++){
		int id=a[i],u=id>0?e[id].u:-id,v=(id>0?e[id].v:-id)+n;
		nxt[i*2]=head[u];head[u]=i*2;
		nxt[i*2+1]=head[v];head[v]=i*2+1;
	}
	for(int s=1;s<=2*n;s++){
		int u=s,c=0;
		while(head[u]!=-1){
			int x=head[u];head[u]=nxt[x];int i=x/2;
			if(vis[i])continue;
			vis[i]=true;col[i]=c;c^=1;
			int id=a[i];u=x&1?(id>0?e[id].u:-id):(id>0?e[id].v:-id)+n;
		}
		assert(u==s&&!c);
	}
	return col;
}
vector<int> perfect(const vector<int> &ids,int d)const{
	ll q=1;while(q<(ll)ids.size())q*=2;
	vector<pair<int,ll>> a; a.reserve(ids.size()+n);
	for(int id:ids)a.emplace_back(id,q/d);
	if(q%d)for(int u=1;u<=n;u++)a.emplace_back(-u,q%d);
	while(q>1){
		vector<int> odd;
		for(auto [id,w]:a)if(w&1)odd.push_back(id);
		auto col=split(odd);int c0=0,c1=0;
		for(int i=0;i<(int)odd.size();i++)if(odd[i]<0)(col[i]?c1:c0)++;
		int keep=c1<c0,j=0,k=0;
		for(auto [id,w]:a){
			ll v=w/2;if(w&1)v+=col[k++]==keep;
			if(v)a[j++]={id,v};
		}
		a.resize(j);q/=2;
	}
	vector<int> ans(n);
	for(auto [id,w]:a){assert(id>0&&w==1);ans[e[id].u-1]=id;}
	assert((int)a.size()==n);return ans;
}
void divide(vector<int> a,int d,vector<vector<int>> &ans)const{
	if(!d)return;
	if(d==1){
		vector<int> b(n);for(int id:a)b[e[id].u-1]=id;
		ans.push_back(move(b));return;
	}
	if(d&1){
		ans.push_back(perfect(a,d));const auto &b=ans.back();
		a.erase(remove_if(a.begin(),a.end(),[&](int id){return b[e[id].u-1]==id;}),a.end());
		d--;
	}
	auto col=split(a);vector<int> l,r;l.reserve(a.size()/2);r.reserve(a.size()/2);
	for(int i=0;i<(int)a.size();i++)(col[i]?r:l).push_back(a[i]);
	divide(move(l),d/2,ans);
	// Borrow some completed matchings, so the other degree is a power of two.
	int k=1;while(k<=((d-1)/2))k*=2;
	for(int i=0;i<k-d/2;i++){
		r.insert(r.end(),ans.back().begin(),ans.back().end());ans.pop_back();
	}
	divide(move(r),k,ans);
}
public:
vector<vector<int>> solve()const{
	if(!n)return {};
	vector<int> l(n+1),r(n+1),ids(e.size()-1);
	for(int id=1;id<(int)e.size();id++)l[e[id].u]++,r[e[id].v]++,ids[id-1]=id;
	int d=l[1];for(int u=1;u<=n;u++)assert(l[u]==d&&r[u]==d);
	vector<vector<int>> ans;ans.reserve(d);divide(move(ids),d,ans);
	for(auto &a:ans){
		vector<int> b(n+1);for(int u=1;u<=n;u++)b[u]=e[a[u-1]].v;
		a=move(b);
	}
	return ans;
}
};
// end for graph/regular-bipartite.cpp
/////////////////////////
\end{lstlisting}

\Needspace{10\baselineskip}
\subsection{\texttt{graph/directed-mst.cpp}}
\textbf{指定根的最小树形图}

\TemplUsage{在带权有向图中，以 S 为根，选取 n-1 条边，使 S 能沿选边到达每个点，且总权最小。每个非根点恰有一条入边。}
\TemplApi{\begin{itemize}
\item \texttt{Directed\_MST g; g.setN(n)}；\texttt{id=g.add\_edge(u,v,w)} 加入 u 指向 v 的边，w 为 ll，返回从 1 开始的边编号。
\item \texttt{auto a=g.solve(S)} 返回 \texttt{result}。\texttt{a.ok=false} 表示不存在以 S 为根的生成树形图，此时 id 为空。
\item 有解时 \texttt{a.cost} 为总权，\texttt{a.id[v]} 为 v 的入边编号，\texttt{a.id[S]=0}；遍历所有非根点即可得到实际边集，父亲是 \texttt{g.e[a.id[v]].u}。
\item 点和 id 数组均为 1 下标。n=1 时有解、总权为 0；要求 $1\le S\le n$。允许负权、重边、自环；同端点只保留最小权，平权保留较早加入的边。
\end{itemize}}
\textbf{最小示例（放入 main）：}
\begin{lstlisting}
Directed_MST g; g.setN(4);
g.add_edge(1,2,5); g.add_edge(1,3,8);
g.add_edge(2,3,-2); g.add_edge(3,2,-1);
g.add_edge(3,4,1);
auto a=g.solve(1);
assert(a.ok && a.cost==4 && a.id[1]==0);
assert(a.id[2]==1 && a.id[3]==3 && a.id[4]==5);
assert(!g.solve(4).ok);
\end{lstlisting}
\TemplNote{设 n 为点数、m 为输入边数，时间 $O((n+m)\log(n+1))$，空间 $O(n+m)$；任意显式图至少需要 $O(m)$ 读边，只有稀疏图才能写成 n 乘多对数。使用左偏堆、可回退并查集缩环并恢复选边。中间权值和 cost 都是 \texttt{\_\_int128}；输出总权时先抄 \texttt{basic/int128\_IO.cpp}，不要直接强转 ll。solve 不修改原图，可更换根重新求解。}

\Needspace{6\baselineskip}\textbf{完整代码}
\begin{lstlisting}[firstnumber=1]
////////////////////////////////////////////////////////////////
//
// template for Directed MST (Chu-Liu / Edmonds)
// usage:
//   Directed_MST g; g.setN(n); g.add_edge(u,v,w); auto a=g.solve(S);
//   a.ok; a.cost; a.id[v] is the incoming edge of v, a.id[S]=0.
//   Vertices / edge ids are 1-indexed; edges point away from the root S.
//   Negative weights / parallel edges / loops are allowed. No solution: ok=false.
//   O((n+m)*log(n+1)) time, O(n+m) space; cost uses __int128.
//
////////////////////////////////////////////////////////////////
struct Directed_MST{
struct edge{int u,v;ll w;};
struct result{bool ok=false;__int128 cost=0;vector<int> id;};
int n=0;
vector<edge> e=vector<edge>(1);
void set(){n=0;e.assign(1,{});}
void setN(int _n){assert(0<=_n&&_n<INT_MAX);set();n=_n;}
int add_edge(int u,int v,ll w){
	assert(1<=min(u,v)&&max(u,v)<=n&&e.size()<INT_MAX);
	int id=e.size();e.push_back({u,v,w});return id;
}
private:
struct heap{
	struct node{int ls=0,rs=0,dis=0,id=0;__int128 w=0,lazy=0;};
	vector<node> a=vector<node>(1);
	void add(int u,__int128 w){if(u)a[u].w+=w,a[u].lazy+=w;}
	void down(int u){add(a[u].ls,a[u].lazy);add(a[u].rs,a[u].lazy);a[u].lazy=0;}
	int merge(int u,int v){
		if(!u||!v)return u|v;
		if(a[u].w>a[v].w||(a[u].w==a[v].w&&a[u].id>a[v].id))swap(u,v);
		down(u);a[u].rs=merge(a[u].rs,v);
		if(a[a[u].ls].dis<a[a[u].rs].dis)swap(a[u].ls,a[u].rs);
		a[u].dis=a[a[u].rs].dis+1;return u;
	}
	int push(int u,int id,ll w){
		a.push_back({0,0,1,id,w,0});return merge(u,(int)a.size()-1);
	}
	int pop(int u){down(u);return merge(a[u].ls,a[u].rs);}
};
struct DSU{
	vector<int> fa,siz;
	vector<pair<int,int>> stk;
	DSU(int n):fa(n+1),siz(n+1,1){iota(fa.begin(),fa.end(),0);}
	int find(int u)const{while(fa[u]!=u)u=fa[u];return u;}
	bool join(int u,int v){
		u=find(u);v=find(v);if(u==v)return false;
		if(siz[u]<siz[v])swap(u,v);
		stk.emplace_back(u,v);fa[v]=u;siz[u]+=siz[v];return true;
	}
	void rollback(int t){
		while((int)stk.size()>t){auto [u,v]=stk.back();stk.pop_back();fa[v]=v;siz[u]-=siz[v];}
	}
};
public:
result solve(int S)const{
	assert(1<=S&&S<=n);DSU dsu(n);heap h;
	vector<int> rt(n+1),seen(n+1),path(n),q(n),in(n+1);
	// Keep one cheapest edge per ordered pair, also bounding heap size by n*n.
	{
		vector<vector<int>> edg(n+1);vector<int> vis(n+1),best(n+1),a;
		for(int id=1;id<(int)e.size();id++)if(e[id].u!=e[id].v&&e[id].v!=S)edg[e[id].v].push_back(id);
		for(int v=1;v<=n;v++){
			a.clear();
			for(int id:edg[v]){
				int u=e[id].u;
				if(vis[u]!=v)vis[u]=v,best[u]=id,a.push_back(u);
				else if(e[id].w<e[best[u]].w)best[u]=id;
			}
			for(int u:a)rt[v]=h.push(rt[v],best[u],e[best[u]].w);
		}
	}
	struct cycle{int u,t;vector<int> e;};
	vector<cycle> cyc;__int128 cost=0;seen[S]=-1;
	for(int s=1;s<=n;s++){
		int u=dsu.find(s),top=0;
		while(!seen[u]){
			while(rt[u]&&dsu.find(e[h.a[rt[u]].id].u)==u)rt[u]=h.pop(rt[u]);
			if(!rt[u])return {};
			int id=h.a[rt[u]].id;__int128 w=h.a[rt[u]].w;
			cost+=w;h.add(rt[u],-w);rt[u]=h.pop(rt[u]);
			q[top]=id;path[top++]=u;seen[u]=s;u=dsu.find(e[id].u);
			if(seen[u]==s){
				int end=top,t=dsu.stk.size(),r=0,v;
				do{v=path[--top];r=h.merge(r,rt[v]);}while(dsu.join(u,v));
				u=dsu.find(u);rt[u]=r;seen[u]=0;
				cyc.push_back({u,t,vector<int>(q.begin()+top,q.begin()+end)});
			}
		}
		for(int i=0;i<top;i++)in[dsu.find(e[q[i]].v)]=q[i];
	}
	for(int i=(int)cyc.size()-1;i>=0;i--){
		auto &c=cyc[i];int id=in[c.u];assert(id);
		dsu.rollback(c.t);
		for(int x:c.e)in[dsu.find(e[x].v)]=x;
		in[dsu.find(e[id].v)]=id;
	}
	return {true,cost,move(in)};
}
};
// end for graph/directed-mst.cpp
/////////////////////////
\end{lstlisting}

\Needspace{10\baselineskip}
\subsection{\texttt{graph/chordal.cpp}}
\textbf{弦图判定与完美消除序列}

\TemplUsage{判断无向图中是否不存在长度至少为 4 的无弦诱导环。如果是弦图，返回一个完美消除序列。}
\TemplApi{\begin{itemize}
\item \texttt{Chordal g; g.setN(n)}，点下标 1..n；\texttt{g.add\_edge(u,v)} 加无向边。自环忽略，重边在求解时去重。
\item \texttt{bool ok=g.solve()}。成功时 \texttt{peo[0..n-1]} 按从前到后给出消除顺序，\texttt{rk[u]} 是 u 在序列中的 1 下标位置。
\item 依次消除 peo 中的点时，每个点尚未消除的邻居构成团。不要把 MCS 的选点顺序直接作为消除顺序，两者相反。
\item 失败时清空 peo、rk。空图和非连通图均支持；可以继续加边再求解。
\end{itemize}}
\textbf{最小示例（放入 main）：}
\begin{lstlisting}
Chordal g; g.setN(4);
g.add_edge(1,2); g.add_edge(2,3);
g.add_edge(3,4); g.add_edge(4,1);
assert(!g.solve() && g.peo.empty());
g.add_edge(1,3); assert(g.solve());
assert(g.peo.size()==4);
for(int i=0;i<4;i++)assert(g.rk[g.peo[i]]==i+1);
\end{lstlisting}
\TemplNote{时间、空间均为 $O(n+m)$，m 为输入边数。MCS 用桶维护已选邻居数；验证时把结点按最靠前的后继邻居分组，统一检查后继邻居的包含关系，避免对每个邻域两两枚举。返回 false 时只给出判定，不构造无弦环证据。}

\Needspace{6\baselineskip}\textbf{完整代码}
\begin{lstlisting}[firstnumber=1]
////////////////////////////////////////////////////////////////
//
// template for Chordal Graph Recognition (MCS)
// usage:
//   Chordal g; g.setN(n); g.add_edge(u,v); bool ok=g.solve();
//   If ok, peo[0..n-1] is the elimination order, rk[u] is 1-indexed.
//   Undirected, vertices 1..n; loops are ignored, parallel edges are merged.
//   O(n+m) time / space; a failed solve clears peo and rk.
//
////////////////////////////////////////////////////////////////
struct Chordal{
int n=0;
vector<vector<int>> edg;
vector<int> peo,rk;
void set(){n=0;edg.clear();peo.clear();rk.clear();}
void setN(int _n){assert(0<=_n&&_n<INT_MAX);set();n=_n;edg.resize(n+1);}
void add_edge(int u,int v){
	assert(1<=min(u,v)&&max(u,v)<=n);
	if(u!=v)edg[u].push_back(v),edg[v].push_back(u);
}
bool solve(){
	vector<int> vis(n+1),cnt(n+1);vector<vector<int>> bucket(n+1),son(n+1);
	for(int u=1;u<=n;u++){
		int k=0;
		for(int v:edg[u])if(vis[v]!=u)vis[v]=u,edg[u][k++]=v;
		edg[u].resize(k);bucket[0].push_back(u);
	}
	peo.resize(n);rk.assign(n+1,0);int mx=0;
	for(int i=n;i>=1;i--){
		while(true){
			while(bucket[mx].empty())mx--;
			int u=bucket[mx].back();bucket[mx].pop_back();
			if(rk[u]||cnt[u]!=mx)continue;
			peo[i-1]=u;rk[u]=i;
			for(int v:edg[u])if(!rk[v]){
				bucket[++cnt[v]].push_back(v);mx=max(mx,cnt[v]);
			}
			break;
		}
	}
	for(int u=1;u<=n;u++){
		int p=0;
		for(int v:edg[u])if(rk[v]>rk[u]&&(!p||rk[v]<rk[p]))p=v;
		if(p)son[p].push_back(u);
	}
	fill(vis.begin(),vis.end(),0);
	for(int p=1;p<=n;p++){
		for(int v:edg[p])vis[v]=p;
		for(int u:son[p])for(int v:edg[u])if(v!=p&&rk[v]>rk[u]&&vis[v]!=p){
			peo.clear();rk.clear();return false;
		}
	}
	return true;
}
};
// end for graph/chordal.cpp
/////////////////////////
\end{lstlisting}
\clearpage
\section{树上算法}
\Needspace{10\baselineskip}
\subsection{\texttt{trees/HLD.cpp}}
\textbf{LCA、祖先和路径分段}

\TemplUsage{对一棵连通无向树预处理，支持 O(1) LCA 和 O(log n) 路径分段。}
\TemplApi{\begin{itemize}
\item \texttt{Tree t; t.setN(n)} 后加入 n-1 条边；可设置 t.rt，再 build。
\item \texttt{lca(u,v)}、\texttt{dist(u,v)} 返回最近公共祖先及边数距离；\texttt{kth\_anc(u,k)} 的 k=0 返回自身。
\item \texttt{in[u]} 是 DFS 编号，\texttt{dfn[i]} 反查原点；子树为 [in[u],out[u]]。
\item \texttt{index\_of\_path(u,v)} 返回覆盖点路径的一组 DFS 序闭区间。
\end{itemize}}
\textbf{最小示例（放入 main）：}
\begin{lstlisting}
Tree t; t.setN(3);
t.add_edge(1,2); t.add_edge(1,3); t.build();
assert(t.lca(2,3)==1 && t.dist(2,3)==2);
int cnt=0;
for(auto [l,r]:t.index_of_path(2,3))cnt+=r-l+1;
assert(cnt==3);
\end{lstlisting}
\TemplNote{配线段树时，把点 u 的值放在位置 in[u]，不要直接用 u。路径段没有保证从 u 到 v 的方向顺序，只能直接用于可交换聚合。按边统计时要另行排除 LCA。与 DC.cpp 都定义 Tree。}

\Needspace{6\baselineskip}\textbf{完整代码}
\begin{lstlisting}[firstnumber=1]
//
// template for Heavy-Light Decomposition, prepared by Otomachi Una
// version 1.0 (Last Update Jun 13th, 2026)
//
// usage:
//   Tree tr; tr.setN(n); tr.add_edge(u,v); tr.build();
//   tr.lca(u,v); tr.index_of_path(u,v);
//
struct Tree{
int n=0,rt=1,tot=0;
vector<vector<int>> edg,st;
vector<int> in,out,siz,ff,dep,dfn,hson,top,ed;
void setN(int _n){
	assert(1<=_n&&_n<INT_MAX-1);n=_n;tot=0;rt=1;
	edg.assign(n+1,{});st.resize(__lg(n)+1);
	for(int i=0;i<(int)st.size();i++)st[i].assign(n-(1<<i)+2,0);
	in.assign(n+1,0);
	out.assign(n+1,0);
	siz.assign(n+1,0);
	ff.assign(n+1,0);
	dep.assign(n+1,0);
	dfn.assign(n+1,0);
	hson.assign(n+1,0);
	top.assign(n+1,0);
	ed.assign(n+1,0);
}
void setRt(int _rt){assert(1<=_rt&&_rt<=n);rt=_rt;}
void add_edge(int u,int v){
	assert(1<=min(u,v)&&max(u,v)<=n&&u!=v);
	edg[u].push_back(v);
	edg[v].push_back(u);
}
void HLD(int u,int fa){
	vector<int> q{u};assert(!siz[u]);siz[u]=1;ff[u]=fa;
	for(int i=0;i<(int)q.size();i++){
		int x=q[i];
		for(int v:edg[x])if(v!=ff[x]){
			assert(!siz[v]);siz[v]=1;ff[v]=x;dep[v]=dep[x]+1;q.push_back(v);
		}
	}
	for(int i=(int)q.size()-1;i>0;i--){
		int x=q[i],p=ff[x];siz[p]+=siz[x];
		if(siz[x]>siz[hson[p]])hson[p]=x;
	}
} 
void HLD1(int u,int fa){
	int start=tot;vector<int> q{u};ff[u]=fa;
	while(!q.empty()){
		int x=q.back();q.pop_back();in[x]=++tot;dfn[tot]=x;st[0][tot]=ff[x];
		for(int i=(int)edg[x].size()-1;i>=0;i--){
			int v=edg[x][i];if(v==ff[x]||v==hson[x])continue;
			top[v]=v;q.push_back(v);
		}
		if(hson[x]){top[hson[x]]=top[x];q.push_back(hson[x]);}
		out[x]=in[x]+siz[x]-1;
	}
	for(int i=tot;i>start;i--){int x=dfn[i];ed[x]=hson[x]?ed[hson[x]]:x;}
}
inline int higher(int x,int y){return dep[x]<dep[y]?x:y;}
void build(){
	// build all basic information in O(nlogn) (without lca O(n))
	assert(1<=rt&&rt<=n);tot=0;
	fill(hson.begin(),hson.end(),0);fill(siz.begin(),siz.end(),0);
	dep[rt]=1;top[rt]=rt;HLD(rt,0);HLD1(rt,0);assert(tot==n);
	for(int i=1;i<(int)st.size();i++) for(int j=1;j<=n-(1<<i)+1;j++)
		st[i][j]=higher(st[i-1][j],st[i-1][j+(1<<(i-1))]);
}
inline int lca(int x,int y){
	// find lca ancestor in O(1)
	assert(1<=min(x,y)&&max(x,y)<=n);
	if(x==y) return x;
	x=in[x],y=in[y];
	if(x>y) swap(x,y);
	x++;int k=__lg(y-x+1);
	return higher(st[k][x],st[k][y-(1<<k)+1]);
}
inline int dist(int u,int v){
	return dep[u]+dep[v]-2*dep[lca(u,v)];
}
inline int kth_anc(int u,int k){
	// find k-th ancestor in O(logn)
	assert(1<=u&&u<=n&&0<=k&&k<dep[u]);
	while(1){
		if(dep[u]-dep[top[u]]>=k) return dfn[in[u]-k];
		k-=dep[u]-dep[top[u]]+1;
		u=ff[top[u]];
		if(!u) return rt;
	}
}
inline vector<array<int,2>> index_of_path(int x,int y){
	// return all in ids in path (x,y)
	assert(1<=min(x,y)&&max(x,y)<=n);
	vector<array<int,2>> id;
	while(top[x]!=top[y]){
		if(dep[top[x]]<dep[top[y]]) swap(x,y);
		id.push_back({in[top[x]],in[x]});
		x=ff[top[x]];
	}
	if(dep[x]>dep[y]) swap(x,y);
	id.push_back({in[x],in[y]});
	return id;
}
void output(){
	cerr<<"***********************************\n";
	cerr<<"structure:"<<endl;
	cerr<<"root = "<<rt<<endl;
	for(int i=1;i<=n;i++) for(int j:edg[i]) if(i<j)
		cerr<<i<<' '<<j<<'\n';
	cerr<<"***********************************\n";
	cerr<<"dfn-order: ";
	for(int i=1;i<=n;i++) cerr<<dfn[i]<<" \n"[i==n];
	cerr<<"top: ";
	for(int i=1;i<=n;i++) cerr<<top[i]<<" \n"[i==n];
}
};
\end{lstlisting}
\par\vspace{1em}

\Needspace{10\baselineskip}
\subsection{\texttt{trees/DC.cpp}}
\textbf{点分治框架}

\TemplUsage{提供寻找重心、递归处理剩余连通块的框架。与 HLD 查询共用 Tree，题目需要的统计尚未实现。}
\TemplApi{\begin{itemize}
\item \texttt{Tree t; t.setN(n)} 后加入树边，设置 rt。
\item 需要 LCA 时先 HLD\_build；DC\_build 从整棵树开始点分治。
\item 在 get 中收集子树信息，在 DC 中实现合并、统计和清空逻辑。
\end{itemize}}
\textbf{粘贴本文件之前：}
\begin{lstlisting}
\end{lstlisting}
\textbf{最小示例（放入 main）：}
\begin{lstlisting}
Tree t; t.setN(3);
t.add_edge(1,2); t.add_edge(2,3);
t.HLD_build(); assert(t.dist(1,3)==2);
t.DC_build(); // traversal only, no problem answer
\end{lstlisting}
\TemplNote{本例只演示初始化和遍历，不会输出路径统计答案。不能原样当作“距离为 k 的点对计数”使用。}

\Needspace{6\baselineskip}\textbf{完整代码}
\begin{lstlisting}[firstnumber=1]
//
// template for Centroid Decomposition, prepared by Otomachi Una
// version 1.0 (Last Update Jun 13th, 2026)
//
// usage:
//   Tree tr; tr.setN(n); tr.add_edge(u,v); tr.HLD_build(); tr.DC_build();
//   fill logic inside DC();
//
struct Tree{
int n=0,rt=1,tot=0;
vector<vector<int>> edg,st;
vector<int> in,out,siz,ff,dep,dfn,hson,top,ed,msiz,csiz;
vector<char> vis;int tsiz=0,cen=0;
void setN(int _n){
	assert(1<=_n&&_n<INT_MAX-1);n=_n;tot=0;rt=1;
	edg.assign(n+1,{});st.resize(__lg(n)+1);
	for(int i=0;i<(int)st.size();i++)st[i].assign(n-(1<<i)+2,0);
	in.assign(n+1,0);
	out.assign(n+1,0);
	siz.assign(n+1,0);
	ff.assign(n+1,0);
	dep.assign(n+1,0);
	dfn.assign(n+1,0);
	hson.assign(n+1,0);
	top.assign(n+1,0);
	ed.assign(n+1,0);
	msiz.assign(n+1,0);
	csiz.assign(n+1,0);
	vis.assign(n+1,0);tsiz=cen=0;
}
void add_edge(int u,int v){
	assert(1<=min(u,v)&&max(u,v)<=n&&u!=v);
	edg[u].push_back(v);
	edg[v].push_back(u);
}
void HLD(int u,int fa){
	vector<int> q{u};assert(!siz[u]);siz[u]=1;ff[u]=fa;
	for(int i=0;i<(int)q.size();i++){
		int x=q[i];
		for(int v:edg[x])if(v!=ff[x]){
			assert(!siz[v]);siz[v]=1;ff[v]=x;dep[v]=dep[x]+1;q.push_back(v);
		}
	}
	for(int i=(int)q.size()-1;i>0;i--){
		int x=q[i],p=ff[x];siz[p]+=siz[x];
		if(siz[x]>siz[hson[p]])hson[p]=x;
	}
} 
void HLD1(int u,int fa){
	int start=tot;vector<int> q{u};ff[u]=fa;
	while(!q.empty()){
		int x=q.back();q.pop_back();in[x]=++tot;dfn[tot]=x;st[0][tot]=ff[x];
		for(int i=(int)edg[x].size()-1;i>=0;i--){
			int v=edg[x][i];if(v==ff[x]||v==hson[x])continue;
			top[v]=v;q.push_back(v);
		}
		if(hson[x]){top[hson[x]]=top[x];q.push_back(hson[x]);}
		out[x]=in[x]+siz[x]-1;
	}
	for(int i=tot;i>start;i--){int x=dfn[i];ed[x]=hson[x]?ed[hson[x]]:x;}
}
inline int higher(int x,int y){return dep[x]<dep[y]?x:y;}
void HLD_build(){
	// build all basic information in O(nlogn) (without lca O(n))
	assert(1<=rt&&rt<=n);tot=0;
	fill(hson.begin(),hson.end(),0);fill(siz.begin(),siz.end(),0);
	dep[rt]=1;top[rt]=rt;HLD(rt,0);HLD1(rt,0);assert(tot==n);
	for(int i=1;i<(int)st.size();i++) for(int j=1;j<=n-(1<<i)+1;j++)
		st[i][j]=higher(st[i-1][j],st[i-1][j+(1<<(i-1))]);
}
inline int lca(int x,int y){
	// find lca ancestor in O(1)
	assert(1<=min(x,y)&&max(x,y)<=n);
	if(x==y) return x;
	x=in[x],y=in[y];
	if(x>y) swap(x,y);
	x++;int k=__lg(y-x+1);
	return higher(st[k][x],st[k][y-(1<<k)+1]);
}
inline int dist(int u,int v){
	return dep[u]+dep[v]-2*dep[lca(u,v)];
}
inline int kth_anc(int u,int k){
	// find k-th ancestor in O(logn)
	assert(1<=u&&u<=n&&0<=k&&k<dep[u]);
	while(1){
		if(dep[u]-dep[top[u]]>=k) return dfn[in[u]-k];
		k-=dep[u]-dep[top[u]]+1;
		u=ff[top[u]];
		if(!u) return rt;
	}
}
void get_size(int u,int fa){
	vector<array<int,2>> q{{u,fa}};
	for(int i=0;i<(int)q.size();i++){
		auto [x,p]=q[i];csiz[x]=1;msiz[x]=0;
		for(int v:edg[x])if(v!=p&&!vis[v])q.push_back({v,x});
	}
	for(int i=(int)q.size()-1;i>=0;i--){
		auto [x,p]=q[i];msiz[x]=max(msiz[x],tsiz-csiz[x]);
		if(msiz[x]<msiz[cen])cen=x;
		if(i){csiz[p]+=csiz[x];msiz[p]=max(msiz[p],csiz[x]);}
	}
}
void get(int u,int fa){
	vector<array<int,2>> q{{u,fa}};
	for(int i=0;i<(int)q.size();i++){
		auto [x,p]=q[i];
		// collect information at x here
		for(int v:edg[x])if(v!=p&&!vis[v])q.push_back({v,x});
	}
}
void DC(int u){
	vis[u]=true;
	for(int v:edg[u]) if(!vis[v]){
		// get some info?
		get(v,0);
	}
	// solve whole
	for(int v:edg[u]) if(!vis[v]){
		tsiz=csiz[v];cen=0;
		get_size(v,0);get_size(cen,0);DC(cen);
	}
	// cerr<<"finish"<<endl;
}
void DC_build(){
	assert(1<=rt&&rt<=n);fill(vis.begin(),vis.end(),0);
	tsiz=n;msiz[cen=0]=1e9;get_size(rt,0);get_size(cen,0);DC(cen);
}
void output(){
	cerr<<"***********************************\n";
	cerr<<"structure:"<<endl;
	cerr<<"root = "<<rt<<endl;
	for(int i=1;i<=n;i++) for(int j:edg[i]) if(i<j)
		cerr<<i<<' '<<j<<'\n';
	cerr<<"***********************************\n";
	cerr<<"dfn-order: ";
	for(int i=1;i<=n;i++) cerr<<dfn[i]<<" \n"[i==n];
	cerr<<"top: ";
	for(int i=1;i<=n;i++) cerr<<top[i]<<" \n"[i==n];
}
// modify from here
void solve(){
	// do not forget setN
	// use add_edge(u,v) to add a edge
	// do not forget HLD_build()
	int N;cin>>N;setN(N);
	for(int i=1,f;i<=n;i++){
		cin>>f;
		if(f) add_edge(i,f);
		else rt=i;
	}
	HLD_build();
	DC_build();
}
};
\end{lstlisting}\Needspace{10\baselineskip}
\subsection{\texttt{trees/cartesian.cpp}}
\textbf{小根堆笛卡尔树}

\TemplUsage{O(n) 构建，同时满足下标的中序顺序与值的小根堆序，空间 O(n)。}
\TemplApi{\texttt{auto [ls,rs]=cartesian(a)} 接受从 0 开始的 vector，返回两个同长度 vector，分别为左右儿子下标。无儿子用 -1。相等值保留较早下标在上方。根是原数组中第一个最小值的位置。}
\textbf{最小示例（放入 main）：}
\begin{lstlisting}
vector<int> a{3,1,2,1};
auto [ls,rs]=cartesian(a);
int rt=min_element(a.begin(),a.end())-a.begin();
assert(rt==1 && ls[rt]==0 && rs[rt]==3 && ls[3]==2);
\end{lstlisting}
\TemplNote{空输入返回两个空 vector，此时没有根。构造无递归，单调数组形成长链也可直接构建；后续遍历要自行避免递归爆栈。类型只需支持小于比较。返回值不含原数组或额外哨兵。}

\Needspace{6\baselineskip}\textbf{完整代码}
\begin{lstlisting}[firstnumber=1]
////////////////////////////////////////////////////////////////
//
// template for Cartesian Tree (min heap)
//
// usage:
//   auto [ls,rs]=cartesian(a); // 0-indexed, absent child = -1
//   O(n); equal values keep the earlier index above the later one.
//   The root is the first minimum; empty input returns two empty vectors.
//
////////////////////////////////////////////////////////////////
template<typename T> pair<vector<int>,vector<int>> cartesian(const vector<T> &a){
	assert(a.size()<(size_t)INT_MAX);int n=a.size();
	vector<int> ls(n,-1),rs(n,-1),st;st.reserve(n);
	for(int i=0;i<n;i++){
		int last=-1;
		while(!st.empty()&&a[i]<a[st.back()]){last=st.back();st.pop_back();}
		ls[i]=last;
		if(!st.empty())rs[st.back()]=i;
		st.push_back(i);
	}
	return {move(ls),move(rs)};
}
// end for trees/cartesian.cpp
/////////////////////////
\end{lstlisting}

\par\vspace{1em}

\clearpage
\section{数论}
\Needspace{10\baselineskip}
\subsection{\texttt{number theory/basic.cpp}}
\textbf{模逆元和扩展 CRT}

\TemplUsage{提供 gcd、lcm、逆元与不要求模数互素的同余方程合并。}
\TemplApi{\begin{itemize}
\item \texttt{inv(a,m)} 要求 m>0 且 gcd(a,m)=1。
\item \texttt{COE(r,m)} 表示 x=r (mod m)，a+b 或 a+=b 合并两个方程。
\item 结果 empty() 为 true 表示无解；否则 r 是标准余数、M 是合并模数，accept(x) 判断满足性。
\end{itemize}}
\textbf{最小示例（放入 main）：}
\begin{lstlisting}
COE a(2,3),b(3,5); auto c=a+b;
assert(!c.empty() && c.r==8 && c.M==15);
assert(c.accept(23) && inv(3,11)==4);
assert((COE(0,2)+COE(1,2)).empty());
\end{lstlisting}
\TemplNote{合并后的模数须能放入 ll，内部乘法虽提升到 \_\_int128，但结果不是任意精度。inv 函数会与部分模数模板的全局 inv 数组重名，需要按实际组合改名。}

\Needspace{6\baselineskip}\textbf{完整代码}
\begin{lstlisting}[firstnumber=1]
//
// template for inv, CRT, etc.
// version 1.0 (Last Update Jun 23rd, 2026)
//
// usage:
//    inv(x,M); COE a(r,M); COE c=a+b;
//   c.empty() if no solution
//
using ull=unsigned long long;
inline ull gcd(ull x,ull y){
	// Stein's Algorithm
	if(!x||!y) return x|y;
	int c=__builtin_ctzll(x|y);
	x>>=__builtin_ctzll(x),y>>=__builtin_ctzll(y);
	while(x!=y){
		if(x<y)swap(x,y);
		x-=y;
		x>>=__builtin_ctzll(x);
	}
	return y<<c;
}
inline ull lcm(ull x,ull y){
	if(!x||!y)return 0;
	x/=gcd(x,y);assert(x<=ULLONG_MAX/y);return x*y;
}
inline ll inv(ll x,ll M){
	// return x^{-1} mod M
	assert(M>0);x%=M;if(x<0)x+=M;
	ull a=M,b=x;__int128 y=0,z=1;
	while(b){
		const ull q=a/b,c=a-q*b;
		a=b,b=c;
		const __int128 w=y-static_cast<__int128>(q)*z;
		y=z,z=w;
	}
	assert(a==1U);y%=M;
	return (y>=0?y:M+y);
}
// Congruence equation, make, merge, etc.
struct COE{
ll r,M;
COE():r(0),M(1){}
COE(ll _r,ll _M):r(_r),M(_M){assert(M>0);r%=M;if(r<0)r+=M;}
inline bool empty()const{return M==-1;}
inline bool accept(ll x)const{if(empty())return false;x%=M;return (x<0?x+M:x)==r;}
inline COE& operator +=(const COE &x){
	if(empty()||x.empty()){M=-1;return *this;}
	ll g=gcd(M,x.M);
	if(r%g!=x.r%g){M=-1;return *this;}
	assert(M/g<=LLONG_MAX/x.M);
	ll i=inv(M/g,x.M/g),new_M=M/g*x.M;
	r=((__int128)r+static_cast<__int128>(x.r-r)/g*i%(x.M/g)*M)%new_M;
	M=new_M;
	if(r<0) r+=M;
	return *this;
}
inline COE operator +(const COE &x)const{
	return COE(*this)+=x;
}
};
\end{lstlisting}
\par\vspace{1em}

\Needspace{10\baselineskip}
\subsection{\texttt{number theory/linear sieve.cpp}}
\textbf{线性筛素数}

\TemplUsage{O(n) 时间和空间，适合需要枚举范围内全部素数。}
\TemplApi{\begin{itemize}
\item \texttt{init(n)} 重建范围 0..n，n>=1。
\item \texttt{is\_prime[x]} 判断 x 是否为素数；\texttt{prime[1..m]} 为升序素数，prime[0] 是占位。
\end{itemize}}
\textbf{最小示例（放入 main）：}
\begin{lstlisting}
init(20);
assert(m==8 && prime[1]==2 && prime[m]==19);
assert(is_prime[17] && !is_prime[1]);
\end{lstlisting}
\TemplNote{不是任意大整数判素接口，不能访问 n 以外。全局 init、m、is\_prime 名字较通用，与其他模板组合时应检查冲突。}

\Needspace{6\baselineskip}\textbf{完整代码}
\begin{lstlisting}[firstnumber=1]
//
// template for Linear Sieve
//
// usage:
//   init(n); is_prime[x]; prime[i]  // i-th prime, 1-index
//
vector<char> is_prime;
vector<int> prime;
int m=0;
void init(int n){
	assert(1<=n&&n<INT_MAX);
	is_prime.assign(n+1,true);prime.assign(1,0);m=0;
	is_prime[0]=is_prime[1]=false;
	for(int i=2;i<=n;i++){
		if(is_prime[i])prime.push_back(i),++m;
		for(int j=1;j<=m&&prime[j]<=n/i;j++){
			int x=i*prime[j];is_prime[x]=false;
			if(i%prime[j]==0)break;
		}
	}
}
\end{lstlisting}
\par\vspace{1em}

\Needspace{10\baselineskip}
\subsection{\texttt{number theory/pollard-rho.cpp}}
\textbf{大整数判素与质因数分解}

\TemplUsage{Miller-Rabin 配合 Pollard-Rho，适用于范围内不适合筛到上界的整数。随机算法的运行时间依赖输入。}
\TemplApi{\begin{itemize}
\item \texttt{is\_prime(n)} 判断素数。
\item \texttt{factor(n)} 返回升序质因子，含重数；n=1 返回空。
\item \texttt{pollard\_rho(n)} 提供一个因子，素数时返回自身。
\end{itemize}}
\textbf{最小示例（放入 main）：}
\begin{lstlisting}
auto f=factor(360);
assert((f==vector<ull>{2,2,2,3,3,5}));
assert(is_prime(1000000007ULL));
\end{lstlisting}
\TemplNote{factor 的输入要求 1<=n<2\textasciicircum{}62，不能传 0。不是完整 uint64 范围版本。判素函数与线性筛的 is\_prime 数组同名，不能原样同贴。}

\Needspace{6\baselineskip}\textbf{完整代码}
\begin{lstlisting}[firstnumber=1]
//
// template for Pollard-Rho, prepared by Otomachi Una
// version 1.0 (Last Update Jun 17th, 2026)
//
// usage:
//   is_prime(n); auto f = factor(n);
//
using ull=unsigned long long;
using u128=__uint128_t;
inline ull rho_gcd(ull x,ull y){
	// Stein's Algorithm
	if(!x||!y) return x|y;
	int c=__builtin_ctzll(x|y);
	x>>=__builtin_ctzll(x),y>>=__builtin_ctzll(y);
	while(x!=y){
		if(x<y)swap(x,y);
		x-=y;
		x>>=__builtin_ctzll(x);
	}
	return y<<c;
}
bool is_prime(ull n){
	if(n<2)return false;
	for(auto y:{2,3,5}){
		if(n==y)return true;
		if(n%y==0)return false;
	}
	assert(n<(1ull<<62));// use Montgomery
	ull n2=n*2,r=n&3;
	for(int _=0;_<5;_++)r*=2-n*r;
	r=-r;
	ull t=-n%n,e=-u128(n)%n;
	auto redc=[&](u128 x)->ull{return (x+u128((ull)(x)*r)*n)>>64;};
	auto mul=[&](ull x,ull y)->ull{return redc(u128(x)*y);};
	auto de=[&](ull x)->ull{x=redc(x);return x<n?x:x-n;};
	auto en=[&](ull x)->ull{return mul(x,e);};
	auto pow=[&](ull a,ull b){
		ull res=t,base=en(a);
		while(b){
			if(b&1)res=mul(res,base);
			base=mul(base,base);
			b>>=1;
		}
		return res;
	};
	ull d=n-1;int z=__builtin_ctzll(d);d>>=z;
	auto miller_rabin=[&](ull b){
		if(b==0)return true;
		ull y=pow(b,d);
		if(de(y)==1)return true;
		for(int i=0;i<z;i++){
			if(de(y)==n-1)return true;
			y=mul(y,y);
		}
		return false;
	};
	if(n<4759123141ull){
		for(auto b:{2,7,61})if(!miller_rabin(b%n))return false;
	}else{
		for(auto b:{2,325,9375,28178,450775,9780504,1795265022})if(!miller_rabin(b%n))return false;
	}
	return true;
}
ull pollard_rho(ull p){
	assert(p>=2);
	if(p%2==0)return 2;
	if(is_prime(p))return p;
	assert(p<(1ull<<62)); // use Montgomery
	ull n=p,n2=n*2,r=n&3;
	for(int _=0;_<5;_++)r*=2-n*r;
	r=-r;
	ull t=-n%n;
	auto redc=[&](u128 x)->ull{return (x+u128((ull)(x)*r)*n)>>64;};
	auto mul=[&](ull x,ull y)->ull{return redc(u128(x)*y);};
	auto dif=[&](ull x,ull y)->ull{x+=n2-y;return x<n2?x:x-n2;};
	auto de=[&](ull x)->ull{x=redc(x);return x<n?x:x-n;};
	mt19937_64 rnd;
	while(1){
		ull c=rnd()%(p-1)+1,y=rnd()%(p-1)+1;
		auto f=[&](ull x)->ull{return redc(u128(x)*x+c);};
		for(ull s=1;;s<<=1){
			ull x=y;
			const ull m=min(1ull<<max(int(__lg(p))/3-8,0),s);
			for(ull i=0;i<s/m;i++){
				ull w=t,z=y;
				for(ull j=0;j<m;j++){
					y=f(y);
					w=mul(w,dif(y,x));
				}
				ull g=rho_gcd(de(w),n);
				if(g>1){
					if(g<n)return g;
					for(ull j=0;j<m;j++){
						z=f(z);
						if((g=rho_gcd(de(dif(z,x)),n))!=1){
							if(g<n)return g;
							else goto fail;
						}
					}
				}
			}
		}
		fail:;
	}
}
vector<ull> factor(ull x){
	vector<ull> ans;
	auto dfs=[&](auto self,ull x){
		if(x==1)return;
		ull y=pollard_rho(x);
		if(x==y){ans.push_back(x);return;}
		self(self,y);self(self,x/y);
	};dfs(dfs,x);
	sort(ans.begin(),ans.end());
	return ans;
}
\end{lstlisting}
\par\vspace{1em}

\Needspace{10\baselineskip}
\subsection{\texttt{number theory/floor sum.cpp}}
\textbf{类欧几里得整除求和}

\TemplUsage{计算闭区间 i=l..r 上 floor((a*i+b)/c) 的和，时间 O(log c)，返回 \_\_int128。}
\TemplApi{\begin{itemize}
\item \texttt{floor\_sum(l,r,a,b,c)} 参数都是 ll。
\item l,r,a,b 均非负，c>0；l>r 返回 0。
\end{itemize}}
\textbf{最小示例（放入 main）：}
\begin{lstlisting}
auto ans=floor_sum(1,4,3,1,2);
assert(ans==16); // 2+3+5+6
\end{lstlisting}
\TemplNote{不能直接处理负系数或负下标。输出可配 int128\_IO；不要先在 ll 中计算 a*l。答案和内部累计值须在有符号 128 位范围内。}

\Needspace{6\baselineskip}\textbf{完整代码}
\begin{lstlisting}[firstnumber=1]
//
// template for Floor Sum, O(log c)
//
// usage:
//   floor_sum(l,r,a,b,c); // sum floor((a*i+b)/c), l<=i<=r
//   0<=l,r,a,b; c>0; return __int128, answer must fit
//
__int128 floor_sum(ll l,ll r,ll a,ll b,ll c){
	assert(l>=0&&r>=0&&a>=0&&b>=0&&c>0);
	if(l>r)return 0;
	__int128 n=(__int128)r-l+1,A=a,B=(__int128)a*l+b,C=c,ans=0;
	while(1){
		ans+=n*(n-1)/2*(A/C)+n*(B/C);
		A%=C;B%=C;
		__int128 y=A*n+B;
		if(y<C)break;
		n=y/C;B=y%C;swap(A,C);
	}
	return ans;
}
\end{lstlisting}
\par\vspace{1em}

\Needspace{10\baselineskip}
\subsection{\texttt{number theory/du jiao sieve.cpp}}
\textbf{\ensuremath{\mu}、\ensuremath{\varphi} 前缀和}

\TemplUsage{预筛小范围并缓存整除分块递归，分别查询 sum \ensuremath{\mu}(i)、sum \ensuremath{\varphi}(i)，i=1..n。}
\TemplApi{\begin{itemize}
\item \texttt{du\_jiao d; d.init\_mu(B)} 后调用 sum\_mu(n)，返回 ll。
\item \texttt{init\_phi(B)} 后调用 sum\_phi(n)，返回 \_\_int128。
\item B>=1，n>=0；两个初始化和缓存相互独立，可只使用一个。
\end{itemize}}
\textbf{最小示例（放入 main）：}
\begin{lstlisting}
du_jiao d; d.init_mu(3);
assert(d.sum_mu(10)==-1);
d.init_phi(3); assert(d.sum_phi(10)==32);
assert(d.sum_mu(0)==0 && d.sum_phi(0)==0);
\end{lstlisting}
\TemplNote{必须先初始化对应函数。B 是预处理上界而非查询上限，常从 n\textasciicircum{}(2/3) 附近结合内存选择；预筛 O(B)，递归性能随 B 和查询集改变。重复初始化会清空该项缓存。}

\Needspace{6\baselineskip}\textbf{完整代码}
\begin{lstlisting}[firstnumber=1]
//
// template for Du Jiao Sieve
//
// usage:
//   du_jiao d; d.init_mu(B); ll mu=d.sum_mu(n);
//   d.init_phi(B); __int128 phi=d.sum_phi(n);
//   sum over [1,n], n>=0; init only what is needed
//   B ~ n^(2/3); O(B+n/sqrt(B)) time, O(B+n/B) space
//
struct du_jiao{
vector<int> mu;
vector<ll> phi;
unordered_map<ll,ll> mem_mu;
unordered_map<ll,__int128> mem_phi;
void init_mu(int m){
	assert(m>=1&&m<INT_MAX);
	vector<int> prime;vector<bool> vis(m+1,false);
	mu.assign(m+1,0);mem_mu.clear();mu[1]=1;
	for(int i=2;i<=m;i++){
		if(!vis[i])prime.push_back(i),mu[i]=-1;
		for(int p:prime){
			if(p>m/i)break;
			int x=i*p;vis[x]=true;
			if(i%p==0)break;
			mu[x]=-mu[i];
		}
	}
	for(int i=1;i<=m;i++)mu[i]+=mu[i-1];
}
void init_phi(int m){
	assert(m>=1&&m<INT_MAX);
	vector<int> prime;vector<bool> vis(m+1,false);
	phi.assign(m+1,0);mem_phi.clear();phi[1]=1;
	for(int i=2;i<=m;i++){
		if(!vis[i])prime.push_back(i),phi[i]=i-1;
		for(int p:prime){
			if(p>m/i)break;
			int x=i*p;vis[x]=true;
			if(i%p==0){phi[x]=phi[i]*p;break;}
			phi[x]=phi[i]*(p-1);
		}
	}
	for(int i=1;i<=m;i++)phi[i]+=phi[i-1];
}
ll sum_mu(ll n){
	assert(n>=0&&!mu.empty());
	if(n<(ll)mu.size())return mu[n];
	auto it=mem_mu.find(n);
	if(it!=mem_mu.end())return it->second;
	__int128 ans=1;
	for(ll l=2,r;l<=n;){
		r=n/(n/l);ans-=(__int128)(r-l+1)*sum_mu(n/l);
		if(r==n)break;
		l=r+1;
	}
	return mem_mu[n]=(ll)ans;
}
__int128 sum_phi(ll n){
	assert(n>=0&&!phi.empty());
	if(n<(ll)phi.size())return phi[n];
	auto it=mem_phi.find(n);
	if(it!=mem_phi.end())return it->second;
	__int128 ans=(__int128)n*(n+(__int128)1)/2;
	for(ll l=2,r;l<=n;){
		r=n/(n/l);ans-=(__int128)(r-l+1)*sum_phi(n/l);
		if(r==n)break;
		l=r+1;
	}
	return mem_phi[n]=ans;
}
};
\end{lstlisting}
\par\vspace{1em}

\clearpage
\section{多项式}
\Needspace{10\baselineskip}
\subsection{\texttt{polynomial/poly-fast.cpp}}
\textbf{998244353 下的常用多项式}

\TemplUsage{适用于固定模数下的大规模卷积与形式幂级数。卷积 O(n log n)，BM 和普通插值为二次复杂度。}
\TemplApi{\begin{itemize}
\item \texttt{poly f=\{a0,a1,...\}} 表示 a0+a1*x+...，下标就是次数。
\item \texttt{f*g} 是卷积；加减和数乘也支持。\texttt{Inv(f)} 返回模 x\textasciicircum{}n 的倒数，n=f.size()，f[0] 必须非零。
\item \texttt{Ln(f)} 要求 f[0]=1；\texttt{Exp(f)} 要求 f[0]=0；返回长度与输入相同。
\item \texttt{diff/integ} 求导/积分，积分常数取零；\texttt{value(f,x)} 是单点求值。
\item \texttt{BM(seq)} 求线性递推的分母多项式，\texttt{FSPE(F,G,k)} 求 F/G 的第 k 项，\texttt{RSPE(seq,k)} 从已有递推样本推第 k 项，k 从 0 开始。
\end{itemize}}
\textbf{最小示例（放入 main）：}
\begin{lstlisting}
poly f{1,1,0,0};
auto g=Inv(f); auto h=f*g; h.resize(4);
assert(h[0].x==1 && !h[1] && !h[2] && !h[3]);
auto e=Exp(Ln(f));
for(int i=0;i<4;i++)assert(e[i].x==f[i].x);
assert(FSPE({0,1},{1,-1,-1},10).x==55);
\end{lstlisting}
\TemplNote{精度由输入 vector 长度决定，想求前 n 项就先 resize(n)，不要删掉需要保留的高位零。\texttt{operator/} 是相关运算 (f/g)[i]=sum\_j f[i+j]*g[j]，不是多项式除法。形式幂级数相除要 f*Inv(g) 后截断。各多项式文件均内置模数类型和全局表，只选一份，不要再粘贴 ModInt。 NTT 长度不超过 2\textasciicircum{}23，Inv/Ln/Exp 的中间卷积也计入这个限制。lagrange(x,y) 要求取值点互异，允许零点。}
\par\textbf{考场版：}本节已换成普通迭代 NTT；fft 输出按自然顺序排列，invFft 接受同样顺序。不能与原版变换内核混搭。

\Needspace{6\baselineskip}\textbf{完整代码}
\begin{lstlisting}[firstnumber=1]
//
// template for ModInt
// version 3.0 (Last Update Jun 15th, 2026)
//
// usage:
//   init(n); mint a,b; C(n,k); binom(n,k);
//
template <unsigned _M> struct ModInt{
static constexpr unsigned MOD=_M;
static_assert(1<MOD&&MOD<=INT_MAX);
unsigned x;
constexpr ModInt():x(0){}
constexpr ModInt(unsigned y):x(y%MOD){}
constexpr ModInt(int y):x((y%=static_cast<int>(MOD))<0?y+MOD:y){}
constexpr ModInt(unsigned long long y):x(y%MOD){}
constexpr ModInt(long long y):x((y%=static_cast<long long>(MOD))<0?y+MOD:y){}
ModInt& operator +=(const ModInt &a){x=(((x+=a.x)>=MOD)?x-MOD:x);return *this;}
ModInt& operator -=(const ModInt &a){x=(((x-=a.x)>=MOD)?x+MOD:x);return *this;}
ModInt& operator *=(const ModInt &a){x=static_cast<unsigned long long>(x)*a.x%MOD;return *this;}
ModInt& operator /=(const ModInt &a){return (*this)*=a.inv();}
bool operator ==(const ModInt &a)const{return x==a.x;}
bool operator !=(const ModInt &a)const{return x!=a.x;}
explicit operator bool() const { return x != 0; }
bool operator!() const { return x == 0; }
ModInt inv() const {
	unsigned a=MOD,b=x;ll y=0,z=1;
	while(b){
		const unsigned q=a/b,c=a-q*b;
		a=b,b=c;
		const ll w=y-(ll)q*z;
		y=z,z=w;
	}
	assert(a==1);return ModInt(y);
}
ModInt pow(long long k) const {
	unsigned long long e=k;
	ModInt a=*this,b=1;
	if(k<0){a=a.inv();e=0-e;}
	for(;e;e>>=1){if(e&1)b*=a;a*=a;}
	return b;
}
ModInt& operator++() {*this+=1;return *this;}
ModInt& operator--() {*this-=1;return *this;}
ModInt operator +()const{return *this;}
ModInt operator -()const{ModInt a;a.x=(x?MOD-x:0u);return a;}
ModInt operator +(const ModInt &a)const{return ModInt(*this)+=a;}
ModInt operator -(const ModInt &a)const{return ModInt(*this)-=a;}
ModInt operator *(const ModInt &a)const{return ModInt(*this)*=a;}
ModInt operator /(const ModInt &a)const{return ModInt(*this)/=a;}

template<typename T> ModInt friend operator +(T a,const ModInt &b){return ModInt(a)+=b;}
template<typename T> ModInt friend operator -(T a,const ModInt &b){return ModInt(a)-=b;}
template<typename T> ModInt friend operator *(T a,const ModInt &b){return ModInt(a)*=b;}
template<typename T> ModInt friend operator /(T a,const ModInt &b){return ModInt(a)/=b;}
friend istream& operator >> (istream &i,ModInt &x){ll y;if(i>>y)x=ModInt(y);return i;}
friend ostream& operator << (ostream &o,const ModInt &x){o<<x.x;return o;}
};

const int MOD=998244353;
using mint = ModInt<MOD>;
vector<mint> fac{1},ifac{1},inv{0,1};
void init(int n=0){
	assert(0<=n&&n<MOD);
	int m=fac.size();if(n<m)return;
	n=max(n,(int)min(2ll*m,(ll)MOD-1));
	fac.resize(n+1);ifac.resize(n+1);inv.resize(max(n+1,2));
	for(int i=m;i<=n;i++){
		fac[i]=fac[i-1]*i;
		if(i>1)inv[i]=-(MOD/i)*inv[MOD%i];
		ifac[i]=ifac[i-1]*inv[i];
	}
}
inline mint C(int x,int y){
	// choose x from y
	if(x<0||y<x)return 0;
	init(y);return fac[y]*ifac[x]*ifac[y-x];
}
inline mint binom(int y,int x){return C(x,y);}


//
// template for polynomial
// version 7.0 (Last Update Jun 15th, 2026)
// NTT core replaced with the faster version
//
// usage:
//    init(n);
//   poly h=f*g; Inv, Ln, Exp, BM, FSPE, value;
//

using ll=long long;
using poly = vector<mint>;

constexpr unsigned MO = 998244353U;
constexpr unsigned MO2 = 2U * MO;
constexpr int FFT_MAX=23;
void ntt(mint *a,int n,bool inverse){
	assert(1<=n&&n<=(1<<FFT_MAX)&&!(n&(n-1)));
	for(int i=1,j=0;i<n;i++){
		int k=n>>1;
		for(;j&k;k>>=1)j^=k;
		j^=k;if(i<j)swap(a[i],a[j]);
	}
	for(int len=2;len<=n;len<<=1){
		mint step=mint(3).pow((MOD-1)/len);
		if(inverse)step=step.inv();
		for(int i=0;i<n;i+=len){
			mint cur=1;
			for(int j=0;j<len/2;j++,cur*=step){
				mint x=a[i+j],y=a[i+j+len/2]*cur;
				a[i+j]=x+y;a[i+j+len/2]=x-y;
			}
		}
	}
	if(inverse){mint v=mint(n).inv();for(int i=0;i<n;i++)a[i]*=v;}
}
void fft(mint *a,int n){ntt(a,n,false);}
void invFft(mint *a,int n){ntt(a,n,true);}

void fft(poly &as) {
	fft(as.data(), as.size());
}
void invFft(poly &as) {
	invFft(as.data(), as.size());
}

// ----------------------------------------------------------------

poly operator *(poly f, poly g) {
	if(f.empty()||g.empty())return {};
	assert(f.size()+g.size()-1<=(1<<FFT_MAX));
	int n = f.size() + g.size();
	int l = 0; while ((1 << l) < n-1) ++l;
	f.resize(1 << l); g.resize(1 << l);
	fft(f); fft(g);
	for (int i = 0; i < (1 << l); i++) f[i] *= g[i];
	invFft(f);
	f.resize(n - 1);
	return f;
}

poly operator /(poly f, poly g) {
	if(f.empty()||g.empty())return poly(f.size());
	int m = g.size();
	reverse(g.begin(), g.end());
	g = f * g;
	for (int i = 0; i < (int)f.size(); i++) f[i] = g[i + m - 1];
	return f;
}

void operator +=(poly &f, const poly &g) {
	if (f.size() < g.size()) f.resize(g.size());
	for (int i = 0; i < (int)g.size(); i++) f[i] += g[i];
}

poly operator +(poly f, const poly &g) { f += g; return f; }

void operator -=(poly &f, const poly &g) {
	if (f.size() < g.size()) f.resize(g.size());
	for (int i = 0; i < (int)g.size(); i++) f[i] -= g[i];
}

poly operator -(poly f, const poly &g) { f -= g; return f; }

poly operator *(poly f, mint x) {
	for (mint &i : f) i *= x;
	return f;
}

poly operator *(mint x, poly f) {
	for (mint &i : f) i *= x;
	return f;
}

mint value(const poly &f,mint x){
	mint ans=0;
	for(int i=f.size()-1;i>=0;i--) ans=ans*x+f[i];
	return ans;
}

poly Inv(poly f) {
	assert(!f.empty()&&f[0]);
	int n = f.size();
	int l = 0; while ((1 << l) < n) ++l;
	f.resize(1 << l);
	poly g{f[0].inv()}, _f;
	for (int i = 1; i <= l; i++) {
		_f = poly(begin(f), begin(f) + (1 << i));
		g.resize(1 << (i+1)); _f.resize(1 << (i+1));
		fft(g); fft(_f);
		for (int j = 0; j < (1 << (i+1)); j++)
			g[j] = mint(2) * g[j] - g[j] * g[j] * _f[j];
		invFft(g);
		fill(begin(g) + (1 << i), end(g), 0);
	}
	g.resize(n);
	return g;
}

poly integ(poly f) {
	int n = f.size(); ::init(n);f.resize(n + 1);
	for (int i = n; i >= 1; i--) f[i] = f[i - 1] * ::inv[i];
	f[0] = 0;
	return f;
}

poly diff(poly f) {
	if(f.empty())return {};
	int n = f.size();
	for (int i = 0; i < n - 1; i++)
		f[i] = f[i + 1] * (i + 1);
	f.pop_back();
	return f;
}

poly Ln(poly f) {
	assert(!f.empty()&&f[0].x==1);
	poly f_ = diff(f), _f = Inv(f);
	int n = f_.size(), m = _f.size();
	int l = 0; while ((1 << l) < n + m) ++l;
	f_.resize(1 << l); _f.resize(1 << l);
	fft(f_); fft(_f);
	for (int i = 0; i < (1 << l); i++) f_[i] *= _f[i];
	invFft(f_);
	f_ = integ(f_);
	f_.resize(f.size());
	return f_;
}

poly Exp(poly f) {
	assert(!f.empty()&&!f[0]);
	poly g{1}, _f, _g;
	int n = f.size();
	int l = 0; while ((1 << l) < n) ++l;
	f.resize(1 << l);
	for (int i = 1; i <= l; i++) {
		_f = poly(begin(f), begin(f) + (1 << i));
		_g = Ln(g);
		g.resize(1 << (i+1)); _f.resize(1 << (i+1)); _g.resize(1 << (i+1));
		fft(g); fft(_f); fft(_g);
		for (int j = 0; j < (1 << (i+1)); j++)
			g[j] *= mint(1) - _g[j] + _f[j];
		invFft(g);
		fill(begin(g) + (1 << i), end(g), 0);
	}
	g.resize(n);
	return g;
}

poly BM(poly a) {
	poly C{1}, B{1};
	int L = 0, m = 1;
	mint b = 1;
	for (int n = 0; n < (int)a.size(); n++) {
		mint d = 0;
		for (int i = 0; i <= L; i++) d += C[i] * a[n - i];
		if (!d) {
			m++;
		} else {
			poly T = C;
			mint coef = d * b.inv();
			poly xmB(m, 0);
			for (mint val : B) xmB.push_back(val);
			if (C.size() < xmB.size()) C.resize(xmB.size(), 0);
			for (int i = 0; i < (int)xmB.size(); i++) C[i] -= coef * xmB[i];
			if (2 * L <= n) {
				L = n + 1 - L; B = T; b = d; m = 1;
			} else {
				m++;
			}
		}
	}
	return C;
}

mint FSPE(poly F, poly G, ll t) {
	// find [x^t] F/G
	assert(t>=0&&!G.empty()&&G[0]);
	if(F.empty())return 0;
	while (t) {
		poly G1 = G;
		for (int i = 1; i < (int)G1.size(); i += 2) G1[i] = -G1[i];
		F = F * G1; G = G * G1;
		poly f, g;
		for (int i = t & 1; i < (int)F.size(); i += 2) f.push_back(F[i]);
		for (int i = 0; i < (int)G.size(); i += 2) g.push_back(G[i]);
		t >>= 1;
		F = f; G = g;
		if(F.empty())return 0;
	}
	return F[0] / G[0];
}

mint RSPE(poly f, ll x) {
	// find f[x] in O(n^2), note that f must be recursion
	poly g = BM(f);
	poly s(g.size()-1);
	for (int i = 0; i < (int)s.size(); i++)
		for (int j = 0; j <= i && j < (int)f.size(); j++)
			s[i] += f[j] * g[i - j];
	return FSPE(s, g, x);
}

poly lagrange(vector<mint> x,vector<mint> y){
	int n=x.size();assert(y.size()==x.size());
	if(!n)return {};
	poly g(n+1),q(n),f(n);g[0]=1;
	for(int i=0;i<n;i++){
		for(int j=i+1;j>=1;j--)g[j]=g[j-1]-x[i]*g[j];
		g[0]*=-x[i];
	}
	for(int i=0;i<n;i++){
		q[n-1]=g[n];
		for(int j=n-2;j>=0;j--)q[j]=g[j+1]+x[i]*q[j+1];
		mint d=value(q,x[i]);assert(d);
		mint c=y[i]/d;for(int j=0;j<n;j++)f[j]+=q[j]*c;
	}
	return f;
}
// use init(n) before accessing fac / ifac / inv directly
\end{lstlisting}
\par\vspace{1em}

\Needspace{10\baselineskip}
\subsection{\texttt{polynomial/poly-new.cpp}}
\textbf{带长除、多点求值的多项式}

\TemplUsage{同样固定模数 998244353，提供常用幂级数和 Sqrt、Div、Eval 等接口。}
\TemplApi{\begin{itemize}
\item \texttt{poly f=\{a0,a1,...\}} 表示 a0+a1*x+...，下标就是次数。
\item \texttt{f*g} 是卷积；加减和数乘也支持。\texttt{Inv(f)} 返回模 x\textasciicircum{}n 的倒数，n=f.size()，f[0] 必须非零。
\item \texttt{Ln(f)} 要求 f[0]=1；\texttt{Exp(f)} 要求 f[0]=0；返回长度与输入相同。
\item \texttt{diff/integ} 求导/积分，积分常数取零；\texttt{value(f,x)} 是单点求值。
\item \texttt{BM(seq)} 求线性递推的分母多项式，\texttt{FSPE(F,G,k)} 求 F/G 的第 k 项，\texttt{RSPE(seq,k)} 从已有递推样本推第 k 项，k 从 0 开始。
\item \texttt{Sqrt(f)} 当前仅接受 f[0]=1，取常数项为 1 的根。
\item \texttt{Div(f,g,q,r)} 得到多项式长除的商与余式，g 最高次系数非零。\texttt{Eval(f,xs)} 返回各取值点的函数值。
\end{itemize}}
\textbf{最小示例（放入 main）：}
\begin{lstlisting}
poly f{1,1,0,0};
auto g=Inv(f); auto h=f*g; h.resize(4);
assert(h[0].x==1 && !h[1] && !h[2] && !h[3]);
auto e=Exp(Ln(f));
for(int i=0;i<4;i++)assert(e[i].x==f[i].x);
assert(FSPE({0,1},{1,-1,-1},10).x==55);
poly q,r; Div({1,2,1},{1,1},q,r);
assert(q.size()==2 && q[0].x==1 && q[1].x==1);
auto v=Eval({1,2,1},{0,1,2});
assert(v[0].x==1 && v[1].x==4 && v[2].x==9);
\end{lstlisting}
\TemplNote{精度由输入 vector 长度决定，想求前 n 项就先 resize(n)，不要删掉需要保留的高位零。\texttt{operator/} 是相关运算 (f/g)[i]=sum\_j f[i+j]*g[j]，不是多项式除法。形式幂级数相除要 f*Inv(g) 后截断。各多项式文件均内置模数类型和全局表，只选一份，不要再粘贴 ModInt。 不支持任意常数项的模平方根。长除前清除除数尾部多余的零。lagrange 要求互异取值点，允许零点。NTT 上限为 2\textasciicircum{}23。}
\par\textbf{考场版：}本节已换成普通迭代 NTT；fft 输出按自然顺序排列，invFft 接受同样顺序。不能与原版变换内核混搭。

\Needspace{6\baselineskip}\textbf{完整代码}
\begin{lstlisting}[firstnumber=1]
//
// template for ModInt, prepared by Otomachi Una
// version 3.0 (Last Update Jun 15th, 2026)
//
// usage:
//   init(n); mint a,b; poly f={...};
//
template <unsigned _M> struct ModInt{
static constexpr unsigned MOD=_M;
static_assert(1<MOD&&MOD<=INT_MAX);
unsigned x;
constexpr ModInt():x(0){}
constexpr ModInt(unsigned y):x(y%MOD){}
constexpr ModInt(int y):x((y%=static_cast<int>(MOD))<0?y+MOD:y){}
constexpr ModInt(unsigned long long y):x(y%MOD){}
constexpr ModInt(long long y):x((y%=static_cast<long long>(MOD))<0?y+MOD:y){}
ModInt& operator +=(const ModInt &a){x=(((x+=a.x)>=MOD)?x-MOD:x);return *this;}
ModInt& operator -=(const ModInt &a){x=(((x-=a.x)>=MOD)?x+MOD:x);return *this;}
ModInt& operator *=(const ModInt &a){x=static_cast<unsigned long long>(x)*a.x%MOD;return *this;}
ModInt& operator /=(const ModInt &a){return (*this)*=a.inv();}
bool operator ==(const ModInt &a)const{return x==a.x;}
bool operator !=(const ModInt &a)const{return x!=a.x;}
explicit operator bool() const { return x != 0; }
bool operator!() const { return x == 0; }
ModInt inv() const {
	unsigned a=MOD,b=x;ll y=0,z=1;
	while(b){
		const unsigned q=a/b,c=a-q*b;
		a=b,b=c;
		const ll w=y-(ll)q*z;
		y=z,z=w;
	}
	assert(a==1);return ModInt(y);
}
ModInt pow(long long k) const {
	unsigned long long e=k;
	ModInt a=*this,b=1;
	if(k<0){a=a.inv();e=0-e;}
	for(;e;e>>=1){if(e&1)b*=a;a*=a;}
	return b;
}
ModInt& operator++() {*this+=1;return *this;}
ModInt& operator--() {*this-=1;return *this;}
ModInt operator +()const{return *this;}
ModInt operator -()const{ModInt a;a.x=(x?MOD-x:0u);return a;}
ModInt operator +(const ModInt &a)const{return ModInt(*this)+=a;}
ModInt operator -(const ModInt &a)const{return ModInt(*this)-=a;}
ModInt operator *(const ModInt &a)const{return ModInt(*this)*=a;}
ModInt operator /(const ModInt &a)const{return ModInt(*this)/=a;}

template<typename T> ModInt friend operator +(T a,const ModInt &b){return ModInt(a)+=b;}
template<typename T> ModInt friend operator -(T a,const ModInt &b){return ModInt(a)-=b;}
template<typename T> ModInt friend operator *(T a,const ModInt &b){return ModInt(a)*=b;}
template<typename T> ModInt friend operator /(T a,const ModInt &b){return ModInt(a)/=b;}
friend istream& operator >> (istream &i,ModInt &x){ll y;if(i>>y)x=ModInt(y);return i;}
friend ostream& operator << (ostream &o,const ModInt &x){o<<x.x;return o;}
};

const int MOD=998244353;
using mint = ModInt<MOD>;
vector<mint> fac{1},ifac{1},inv{0,1};
void init(int n=0){
	assert(0<=n&&n<MOD);
	int m=fac.size();if(n<m)return;
	n=max(n,(int)min(2ll*m,(ll)MOD-1));
	fac.resize(n+1);ifac.resize(n+1);inv.resize(max(n+1,2));
	for(int i=m;i<=n;i++){
		fac[i]=fac[i-1]*i;
		if(i>1)inv[i]=-(MOD/i)*inv[MOD%i];
		ifac[i]=ifac[i-1]*inv[i];
	}
}
inline mint C(int x,int y){
	// choose x from y
	if(x<0||y<x)return 0;
	init(y);return fac[y]*ifac[x]*ifac[y-x];
}
inline mint binom(int y,int x){return C(x,y);}


//
// template for polynomial, prepared by Otomachi Una
// version 6.1 (Last Update Jun 13th, 2026)
// NTT core replaced with the "tech" version
//
// usage:
//    init(n);
//   Inv, Ln, Exp, Div, Sqrt, Eval, BM;
//

using poly = vector<mint>;

// ---------- NTT constants and roots (tech version) ----------
constexpr unsigned MO = 998244353U;
constexpr unsigned MO2 = 2U * MO;
constexpr int FFT_MAX=23;
void ntt(mint *a,int n,bool inverse){
	assert(1<=n&&n<=(1<<FFT_MAX)&&!(n&(n-1)));
	for(int i=1,j=0;i<n;i++){
		int k=n>>1;
		for(;j&k;k>>=1)j^=k;
		j^=k;if(i<j)swap(a[i],a[j]);
	}
	for(int len=2;len<=n;len<<=1){
		mint step=mint(3).pow((MOD-1)/len);
		if(inverse)step=step.inv();
		for(int i=0;i<n;i+=len){
			mint cur=1;
			for(int j=0;j<len/2;j++,cur*=step){
				mint x=a[i+j],y=a[i+j+len/2]*cur;
				a[i+j]=x+y;a[i+j+len/2]=x-y;
			}
		}
	}
	if(inverse){mint v=mint(n).inv();for(int i=0;i<n;i++)a[i]*=v;}
}
void fft(mint *a,int n){ntt(a,n,false);}
void invFft(mint *a,int n){ntt(a,n,true);}

void fft(poly &as) {
	fft(as.data(), as.size());
}
void invFft(poly &as) {
	invFft(as.data(), as.size());
}

// ----------------------------------------------------------------

poly operator *(poly f, poly g) {
	if(f.empty()||g.empty())return {};
	assert(f.size()+g.size()-1<=(1<<FFT_MAX));
	int n = f.size() + g.size();
	int l = 0; while ((1 << l) < n-1) ++l;
	f.resize(1 << l); g.resize(1 << l);
	fft(f); fft(g);
	for (int i = 0; i < (1 << l); i++) f[i] *= g[i];
	invFft(f);
	f.resize(n - 1);
	return f;
}

poly operator /(poly f, poly g) {
	if(f.empty()||g.empty())return poly(f.size());
	int m = g.size();
	reverse(g.begin(), g.end());
	g = f * g;
	for (int i = 0; i < (int)f.size(); i++) f[i] = g[i + m - 1];
	return f;
}

void operator +=(poly &f, const poly &g) {
	if (f.size() < g.size()) f.resize(g.size());
	for (int i = 0; i < (int)g.size(); i++) f[i] += g[i];
}

poly operator +(poly f, const poly &g) { f += g; return f; }

void operator -=(poly &f, const poly &g) {
	if (f.size() < g.size()) f.resize(g.size());
	for (int i = 0; i < (int)g.size(); i++) f[i] -= g[i];
}

poly operator -(poly f, const poly &g) { f -= g; return f; }

poly operator *(poly f, mint x) {
	for (mint &i : f) i *= x;
	return f;
}

poly operator *(mint x, poly f) {
	for (mint &i : f) i *= x;
	return f;
}

mint value(const poly &f,mint x){
	mint ans=0;
	for(int i=f.size()-1;i>=0;i--) ans=ans*x+f[i];
	return ans;
}

poly Inv(poly f) {
	assert(!f.empty()&&f[0]);
	int n = f.size();
	int l = 0; while ((1 << l) < n) ++l;
	f.resize(1 << l);
	poly g{f[0].inv()}, _f;
	for (int i = 1; i <= l; i++) {
		_f = poly(begin(f), begin(f) + (1 << i));
		g.resize(1 << (i+1)); _f.resize(1 << (i+1));
		fft(g); fft(_f);
		for (int j = 0; j < (1 << (i+1)); j++)
			g[j] = mint(2) * g[j] - g[j] * g[j] * _f[j];
		invFft(g);
		fill(begin(g) + (1 << i), end(g), 0);
	}
	g.resize(n);
	return g;
}

poly integ(poly f) {
	int n = f.size(); ::init(n);f.resize(n + 1);
	for (int i = n; i >= 1; i--) f[i] = f[i - 1] * ::inv[i];
	f[0] = 0;
	return f;
}

poly diff(poly f) {
	if(f.empty())return {};
	int n = f.size();
	for (int i = 0; i < n - 1; i++)
		f[i] = f[i + 1] * (i + 1);
	f.pop_back();
	return f;
}

poly Ln(poly f) {
	assert(!f.empty()&&f[0].x==1);
	poly f_ = diff(f), _f = Inv(f);
	int n = f_.size(), m = _f.size();
	int l = 0; while ((1 << l) < n + m) ++l;
	f_.resize(1 << l); _f.resize(1 << l);
	fft(f_); fft(_f);
	for (int i = 0; i < (1 << l); i++) f_[i] *= _f[i];
	invFft(f_);
	f_ = integ(f_);
	f_.resize(f.size());
	return f_;
}

poly Exp(poly f) {
	assert(!f.empty()&&!f[0]);
	poly g{1}, _f, _g;
	int n = f.size();
	int l = 0; while ((1 << l) < n) ++l;
	f.resize(1 << l);
	for (int i = 1; i <= l; i++) {
		_f = poly(begin(f), begin(f) + (1 << i));
		_g = Ln(g);
		g.resize(1 << (i+1)); _f.resize(1 << (i+1)); _g.resize(1 << (i+1));
		fft(g); fft(_f); fft(_g);
		for (int j = 0; j < (1 << (i+1)); j++)
			g[j] *= mint(1) - _g[j] + _f[j];
		invFft(g);
		fill(begin(g) + (1 << i), end(g), 0);
	}
	g.resize(n);
	return g;
}

poly Sqrt(poly f) {
	assert(!f.empty()&&f[0].x==1);
	poly g{1}, _f, _g;
	int n = f.size();
	int l = 0; while ((1 << l) < n) ++l;
	f.resize(1 << l);
	mint inv2 = mint(2).inv();
	for (int i = 1; i <= l; i++) {
		_f = poly(begin(f), begin(f) + (1 << i));
		_g = Inv(g);
		g.resize(1 << (i+1)); _f.resize(1 << (i+1)); _g.resize(1 << (i+1));
		fft(_f); fft(_g); fft(g);
		for (int j = 0; j < (1 << (i+1)); j++)
			g[j] = (_f[j] + g[j] * g[j]) * inv2 * _g[j];
		invFft(g);
		fill(begin(g) + (1 << i), end(g), 0);
	}
	g.resize(n);
	return g;
}

void Div(poly f, poly g, poly &q, poly &r) {
	assert(!g.empty()&&g.back());
	if(f.size()<g.size()){q={0};r=f;return;}
	int n = f.size() - 1, m = g.size() - 1;
	reverse(f.begin(), f.end());
	reverse(g.begin(), g.end());
	g.resize(n + 1);
	q = f * Inv(g);
	q.resize(n - m + 1);
	reverse(q.begin(), q.end());
	g.resize(m + 1);
	reverse(g.begin(), g.end());
	reverse(f.begin(), f.end());
	poly h = q * g;
	r.resize(m);
	for (int i = 0; i < m; i++) r[i] = f[i] - h[i];
}

poly Eval(poly f, poly x) {
	if(x.empty())return {};
	vector<poly> func;
	map<pair<int, int>, int> mp;
	function<poly(int, int)> prod = [&](int l, int r) {
		poly ans;
		if (l == r) ans = poly{-x[l], 1};
		else {
			int mid = (l + r) >> 1;
			ans = prod(l, mid) * prod(mid + 1, r);
		}
		func.push_back(ans);
		mp[{l, r}] = func.size() - 1;
		return ans;
	};
	prod(0, x.size() - 1);

	function<poly(poly, int, int)> solve = [&](poly f, int l, int r) {
		if (l == r) {
			mint val = 0;
			for (int i = f.size() - 1; i >= 0; i--)
				val = val * x[l] + f[i];
			return poly{val};
		}
		int mid = (l + r) >> 1;
		poly p0 = func[mp[{l, mid}]], p1 = func[mp[{mid + 1, r}]];
		poly d, q;
		Div(f, p0, d, q);
		poly ans = solve(q, l, mid);
		Div(f, p1, d, q);
		poly _ans = solve(q, mid + 1, r);
		for (mint i : _ans) ans.push_back(i);
		return ans;
	};
	return solve(f, 0, x.size() - 1);
}

poly BM(poly a) {
	poly C{1}, B{1};
	int L = 0, m = 1;
	mint b = 1;
	for (int n = 0; n < (int)a.size(); n++) {
		mint d = 0;
		for (int i = 0; i <= L; i++) d += C[i] * a[n - i];
		if (!d) {
			m++;
		} else {
			poly T = C;
			mint coef = d * b.inv();
			poly xmB(m, 0);
			for (mint val : B) xmB.push_back(val);
			if (C.size() < xmB.size()) C.resize(xmB.size(), 0);
			for (int i = 0; i < (int)xmB.size(); i++) C[i] -= coef * xmB[i];
			if (2 * L <= n) {
				L = n + 1 - L; B = T; b = d; m = 1;
			} else {
				m++;
			}
		}
	}
	return C;
}

mint FSPE(poly F, poly G, ll t) {
	// find [x^t] F/G
	assert(t>=0&&!G.empty()&&G[0]);
	if(F.empty())return 0;
	while (t) {
		poly G1 = G;
		for (int i = 1; i < (int)G1.size(); i += 2) G1[i] = -G1[i];
		F = F * G1; G = G * G1;
		poly f, g;
		for (int i = t & 1; i < (int)F.size(); i += 2) f.push_back(F[i]);
		for (int i = 0; i < (int)G.size(); i += 2) g.push_back(G[i]);
		t >>= 1;
		F = f; G = g;
		if(F.empty())return 0;
	}
	return F[0] / G[0];
}

mint RSPE(poly f, ll x) {
	// find f[x] in O(n^2), note that f must be recursion
	poly g = BM(f);
	poly s(g.size()-1);
	for (int i = 0; i < (int)s.size(); i++)
		for (int j = 0; j <= i && j < (int)f.size(); j++)
			s[i] += f[j] * g[i - j];
	return FSPE(s, g, x);
}

poly lagrange(vector<mint> x,vector<mint> y){
	int n=x.size();assert(y.size()==x.size());
	if(!n)return {};
	poly g(n+1),q(n),f(n);g[0]=1;
	for(int i=0;i<n;i++){
		for(int j=i+1;j>=1;j--)g[j]=g[j-1]-x[i]*g[j];
		g[0]*=-x[i];
	}
	for(int i=0;i<n;i++){
		q[n-1]=g[n];
		for(int j=n-2;j>=0;j--)q[j]=g[j+1]+x[i]*q[j+1];
		mint d=value(q,x[i]);assert(d);
		mint c=y[i]/d;for(int j=0;j<n;j++)f[j]+=q[j]*c;
	}
	return f;
}

poly to_ex(poly f){ // f(e^x) 
	int n=f.size();if(!n)return {};::init(n);
	function<pair<poly,poly>(int,int)> solve=[&](int l,int r){
		if(l==r) return make_pair(poly{f[l]},poly{1,MOD-l});
		int mid=l+r>>1;
		auto ls=solve(l,mid),rs=solve(mid+1,r);
		return make_pair(ls.first*rs.second+ls.second*rs.first,ls.second*rs.second);
	};
	auto ans=solve(0,n-1);
	poly g=ans.first*Inv(ans.second);
	g.resize(n);
	for(int i=0;i<n;i++) g[i]=g[i]*ifac[i]; 
	return g;
}

poly S2line(int n){
	// return S(n,i)
	assert(n>=0);::init(n);
	poly F(n+1),G(n+1);
	// S(n,i) * binom(j,i) -> F(j)
	for(int i=0;i<=n;i++){
		G[i]=ifac[i];F[i]=mint(i).pow(n)*ifac[i];
		if(i&1) F[i]*=-1;
	}
	F=F*G;
	F.resize(n+1);
	for(int i=0;i<=n;i++){
		if(i&1) F[i]*=-1;
	}
	return F;
}
\end{lstlisting}
\par\vspace{1em}

\Needspace{10\baselineskip}
\subsection{\texttt{polynomial/poly-old.cpp}}
\textbf{ll 系数的旧版多项式}

\TemplUsage{保留旧题适配接口，系数类型为 ll，默认 MOD=998244353。新题优先用 poly-fast 或 MTT。}
\TemplApi{\begin{itemize}
\item \texttt{poly} 实际是 vector<long long>，输入系数应先归一化到 [0,MOD)。
\item \texttt{f*g} 为卷积，\texttt{inv/ln/exp/sqrt} 使用小写命名。
\item \texttt{init(n)} 准备 fac、ifac、Inv 表；注意 Inv 是整数逆元数组，inv 才是幂级数求逆函数。
\end{itemize}}
\textbf{最小示例（放入 main）：}
\begin{lstlisting}
init(10); poly f{1,2},g{3,4}; auto h=f*g;
assert(h[0]==3 && h[1]==10 && h[2]==8);
poly a{1,1,0,0}; auto b=inv(a);
assert(b[0]==1 && b[1]==MOD-1);
\end{lstlisting}
\TemplNote{不要根据其他版本的接口大小写直接替换。旧版保留许多历史实现，本手册的可运行示例验证乘法和求逆；做递推推项优先选已核对接口的 MTT/poly-fast，避免依赖旧示例 solve。文件含 poly、LL、MP 等宏，组合时应留意污染。}

\Needspace{6\baselineskip}\textbf{完整代码}
\begin{lstlisting}[firstnumber=1]
#define LL __int128
#define MP make_pair
#define poly vector<long long>
//
// template for polynomial (ll NTT)
// version 1.0 (last update: May 24th, 2026)
//
// usage:
//   init(n); poly f,g; f*g; ln(f);
//
const ll MOD=998244353,inv2=(MOD+1)/2;
vector<ll> fac{1},ifac{1},Inv{0,1};
vector<int> bfly(1);vector<ll> inver{0,1};
int clogg(int x){assert(x>0);return x==1?0:32-__builtin_clz(x-1);} 
ll ksm(ll a,int b){ll res=1;while(b){if(b&1)res=res*a%MOD;a=a*a%MOD,b>>=1;}return res;}
void butterfly(int l){
	assert(0<=l&&l<=23);
	static int las=-1;
	if(las!=l){
		las=l;bfly.assign(1<<l,0); 
		for(int i=1;i<(1<<l);i++)
			bfly[i]=(bfly[i>>1]>>1)|((i&1)<<l-1);
	} 
}
void NTT(poly &f,int l,int typ){
	butterfly(l);f.resize(1<<l);
	for(int i=0;i<(1<<l);i++)
		if(bfly[i]<i) swap(f[i],f[bfly[i]]);
	for(int i=0;i<l;i++){
		ll step=ksm(3,MOD-1+(MOD-1>>i+1)*typ);
		for(int j=0;j<(1<<l);j+=(1<<i+1)){
			ll cur=1,l=j+(1<<i);
			for(int k=j;k<l;k++){
				ll u=f[k],v=f[k|(1<<i)]*cur%MOD;
				f[k]=(u+v>=MOD?u+v-MOD:u+v);
				f[k|(1<<i)]=(u>=v?u-v:u-v+MOD);
				cur=cur*step%MOD;
			}
		}
	}
	if(typ==-1){
		ll val=ksm(1<<l,MOD-2);
		for(int i=0;i<(1<<l);i++)
			f[i]=val*f[i]%MOD;
	}
	return;
}
poly operator *(poly f,poly g){
	// h[i+j] <- f[i]*g[j]
	if(f.empty()||g.empty())return {};
	assert(f.size()+g.size()-1<=(1<<23));
	int n=f.size()+g.size(),l=clogg(f.size()+g.size()-1);
	if(n<=64){
		poly h(n-1);
		for(int i=0;i<f.size();i++)
			for(int j=0;j<g.size();j++)
				h[i+j]+=f[i]*g[j]%MOD;
		for(ll &i:h) i%=MOD;
		return h;
	}
	NTT(f,l,1);NTT(g,l,1);
	for(int i=0;i<(1<<l);i++)
		f[i]=f[i]*g[i]%MOD;
	NTT(f,l,-1);f.resize(n-1);
	return f;
}
poly operator /(poly f,poly g){
	if(f.empty()||g.empty())return poly(f.size());
	// h[i-j] <- f[i]*g[j]
	int m=g.size();
	reverse(g.begin(),g.end());
	g=f*g;
	for(int i=0;i<f.size();i++) f[i]=g[i+m-1];
	return f;
}
void operator +=(poly &f,const poly &g){
	int n=max(f.size(),g.size());
	f.resize(n);
	for(int i=0;i<g.size();i++)
		f[i]=(f[i]+g[i]>=MOD?f[i]+g[i]-MOD:f[i]+g[i]);
	return;
}
poly operator +(poly f,poly g){
	f+=g;
	return f;
}
void operator -=(poly &f,const poly &g){
	int n=max(f.size(),g.size());
	f.resize(n);
	for(int i=0;i<g.size();i++)
		f[i]=(f[i]-g[i]<0?f[i]-g[i]+MOD:f[i]-g[i]);
	return;
}
poly operator -(poly f,poly g){
	f-=g;
	return f;
}
poly operator *(poly f,ll x){
	for(int i=0;i<f.size();i++) (f[i]*=x)%=MOD;
	return f;
}
poly operator *(ll x,poly f){
	for(int i=0;i<f.size();i++) (f[i]*=x)%=MOD;
	return f;
}
poly inv(poly f){
	assert(!f.empty()&&f[0]);
	int n=f.size(),l=clogg(n);
	f.resize(1<<l);
	poly g{ksm(f[0],MOD-2)},_f;
	for(int i=1;i<=l;i++){
		_f=poly(begin(f),begin(f)+(1<<i));
		NTT(g,i+1,1);NTT(_f,i+1,1);
		for(int j=0;j<(1<<i+1);j++)
			g[j]=(2*g[j]-g[j]*g[j]%MOD*_f[j]%MOD+MOD)%MOD;
		NTT(g,i+1,-1);
		fill(begin(g)+(1<<i),end(g),0);
	}
	g.resize(n);
	return g;
}
poly integ(poly f){
	int n=f.size();f.resize(n+1);
	assert(n<MOD);
	int m=inver.size();
	if(n>=m){
		int k=max(n,(int)min(2ll*m,MOD-1));inver.resize(k+1);
		for(int i=m;i<=k;i++)inver[i]=MOD-(MOD/i)*inver[MOD%i]%MOD;
	}
	for(int i=n;i>=1;i--) f[i]=f[i-1]*inver[i]%MOD;
	f[0]=0;
	return f;
}
poly diff(poly f){
	if(f.empty())return {};
	int n=f.size();
	for(int i=0;i<n-1;i++)
		f[i]=f[i+1]*(i+1)%MOD;
	f.pop_back();
	return f;
} 
poly ln(poly f){
	assert(!f.empty()&&f[0]==1);
	poly f_=diff(f),_f=inv(f);
	int n=f_.size(),m=_f.size();
	int l=clogg(n+m);
	NTT(f_,l,1);NTT(_f,l,1);
	for(int i=0;i<(1<<l);i++)
		f_[i]=f_[i]*_f[i]%MOD;
	NTT(f_,l,-1);
	f_=integ(f_);
	f_.resize(f.size());
	return f_;
}
poly exp(poly f){
	assert(!f.empty()&&f[0]==0);
	poly g{1},_f,_g;
	int n=f.size(),l=clogg(n);
	f.resize(1<<l);
	for(int i=1;i<=l;i++){
		_f=poly(begin(f),begin(f)+(1<<i));
		_g=ln(g);
		NTT(g,i+1,1);NTT(_f,i+1,1);NTT(_g,i+1,1);
		for(int j=0;j<(1<<i+1);j++)
			g[j]=g[j]*(1-_g[j]+_f[j]+MOD)%MOD;
		NTT(g,i+1,-1);
		fill(begin(g)+(1<<i),end(g),0);
	}
	g.resize(n);
	return g;
}
poly sqrt(poly f){
	assert(!f.empty()&&f[0]==1);
	poly g{1},_f,_g;
	int n=f.size(),l=clogg(n);
	f.resize(1<<l);
	for(int i=1;i<=l;i++){
		_f=poly(begin(f),begin(f)+(1<<i));
		_g=inv(g);
		NTT(_f,i+1,1);NTT(_g,i+1,1);NTT(g,i+1,1);
		for(int j=0;j<(1<<i+1);j++)
			g[j]=(_f[j]+g[j]*g[j]%MOD)*inv2%MOD*_g[j]%MOD;
		NTT(g,i+1,-1);
		fill(begin(g)+(1<<i),end(g),0);
	}
	g.resize(n);
	return g;
} 
void div(poly f,poly g,poly &q,poly &r){// ¹¹Ôì p,q ʹµÃ f=g*q+r; 
	assert(!g.empty()&&g.back());
	if(f.size()<g.size()){q={0};r=f;return;}
	int n=f.size()-1,m=g.size()-1;
	reverse(f.begin(),f.end());
	reverse(g.begin(),g.end());
	g.resize(n+1);
	q=(f*inv(g));q.resize(n-m+1);
	reverse(q.begin(),q.end());
	g.resize(m+1);
	reverse(g.begin(),g.end());
	reverse(f.begin(),f.end());
	poly h=q*g;
	r.resize(m);
	for(int i=0;i<m;i++) r[i]=(f[i]-h[i]+MOD)%MOD;
	return;
}
poly eval(poly f,poly x){
	if(x.empty())return {};
	vector<poly> func; map<pair<int,int>,int> mp;
	function<poly(int,int)> prod=[&](int l,int r){
		poly ans;
		if(l==r) ans=(poly){MOD-x[l],1};
		else {
			int mid=l+r>>1;
			ans=prod(l,mid)*prod(mid+1,r);
		} 
		func.push_back(ans);
		mp[MP(l,r)]=func.size()-1;
		return ans;
	};prod(0,x.size()-1);
	int t=-1;
	function<poly(poly,int,int)> solve=[&](poly f,int l,int r){
		if(l==r){
			ll val=0;
			for(int i=f.size()-1;i>=0;i--)
				val=(val*x[l]+f[i])%MOD;
			return (poly){val};
		}
		int mid=l+r>>1;
		poly p0=func[mp[MP(l,mid)]],p1=func[mp[MP(mid+1,r)]];
		poly d,q;
		div(f,p0,d,q);
		poly ans=solve(q,l,mid);
		div(f,p1,d,q);
		poly _ans=solve(q,mid+1,r);
		for(int i:_ans) ans.push_back(i);
		return ans;
	};
	return solve(f,0,x.size()-1);
}
poly BM(poly a){
	poly c{1},b{1};int l=0,m=1;ll last=1;
	for(int i=0;i<(int)a.size();i++){
		ll d=0;for(int j=0;j<=l;j++)d=(d+c[j]*a[i-j])%MOD;
		if(!d){m++;continue;}
		poly t=c;ll v=d*ksm(last,MOD-2)%MOD;
		if(c.size()<b.size()+m)c.resize(b.size()+m);
		for(int j=0;j<(int)b.size();j++)c[j+m]=(c[j+m]-v*b[j]%MOD+MOD)%MOD;
		if(2*l<=i){l=i+1-l;b=move(t);last=d;m=1;}
		else m++;
	}
	return c;
}
ll FSPE(poly F,poly G,ll t){// [x^t]F(x)/G(x)
	assert(t>=0&&!G.empty()&&G[0]);
	if(F.empty())return 0;
	while(t){
		poly G1=G;
		for(int i=1;i<G1.size();i+=2) G1[i]=MOD-G1[i];
		F=F*G1,G=G*G1;
		poly f,g;
		for(int i=t&1;i<F.size();i+=2) f.push_back(F[i]);
		for(int i=0;i<G.size();i+=2) g.push_back(G[i]);
		t>>=1;
		F=f,G=g;
		if(F.empty())return 0;
	}
	return F[0]*ksm(G[0],MOD-2)%MOD;

}
ll RSPE(poly f,ll x){
	poly g=BM(f),s(g.size()-1);
	for(int i=0;i<(int)s.size();i++)
		for(int j=0;j<=i&&j<(int)f.size();j++)s[i]=(s[i]+f[j]*g[i-j])%MOD;
	return FSPE(s,g,x);
}
mt19937 rnd(time(0));
void init(int n=0){
	assert(0<=n&&n<MOD);
	int m=fac.size();if(n<m)return;
	n=max(n,(int)min(2ll*m,MOD-1));
	fac.resize(n+1);ifac.resize(n+1);Inv.resize(max(n+1,2));
	for(int i=m;i<=n;i++){
		fac[i]=fac[i-1]*i%MOD;
		if(i>=2)Inv[i]=MOD-Inv[MOD%i]*(MOD/i)%MOD;
		ifac[i]=ifac[i-1]*Inv[i]%MOD;
	}
}
ll C(int x,int y){if(x<0||x>y)return 0;init(y);return fac[y]*ifac[x]%MOD*ifac[y-x]%MOD;} 
\end{lstlisting}
\par\vspace{1em}

\Needspace{10\baselineskip}
\subsection{\texttt{polynomial/MTT.cpp}}
\textbf{非 NTT 模数下的多项式}

\TemplUsage{默认目标模数 10\textasciicircum{}9+7，三模 NTT 后 CRT 还原；常用 FPS 功能与 poly-fast 对齐，但每次卷积都会回到目标模数。}
\TemplApi{\begin{itemize}
\item \texttt{poly f=\{a0,a1,...\}} 表示 a0+a1*x+...，下标就是次数。
\item \texttt{f*g} 是卷积；加减和数乘也支持。\texttt{Inv(f)} 返回模 x\textasciicircum{}n 的倒数，n=f.size()，f[0] 必须非零。
\item \texttt{Ln(f)} 要求 f[0]=1；\texttt{Exp(f)} 要求 f[0]=0；返回长度与输入相同。
\item \texttt{diff/integ} 求导/积分，积分常数取零；\texttt{value(f,x)} 是单点求值。
\item \texttt{BM(seq)} 求线性递推的分母多项式，\texttt{FSPE(F,G,k)} 求 F/G 的第 k 项，\texttt{RSPE(seq,k)} 从已有递推样本推第 k 项，k 从 0 开始。
\item \texttt{Poly} 与 \texttt{poly} 是同一个类型；\texttt{lagrange(xs,ys)} 支持互异取值点，包括零点，返回次数小于点数的插值多项式。
\end{itemize}}
\textbf{最小示例（放入 main）：}
\begin{lstlisting}
poly f{1,1,0,0};
auto g=Inv(f); auto h=f*g; h.resize(4);
assert(h[0].x==1 && !h[1] && !h[2] && !h[3]);
auto e=Exp(Ln(f));
for(int i=0;i<4;i++)assert(e[i].x==f[i].x);
assert(FSPE({0,1},{1,-1,-1},10).x==55);
auto p=lagrange({0,1,2},{1,4,9});
assert(p[0].x==1 && p[1].x==2 && p[2].x==1);
\end{lstlisting}
\TemplNote{精度由输入 vector 长度决定，想求前 n 项就先 resize(n)，不要删掉需要保留的高位零。\texttt{operator/} 是相关运算 (f/g)[i]=sum\_j f[i+j]*g[j]，不是多项式除法。形式幂级数相除要 f*Inv(g) 后截断。各多项式文件均内置模数类型和全局表，只选一份，不要再粘贴 ModInt。 卷积结果长度最多 2\textasciicircum{}23，FPS 输入长度最多 2\textasciicircum{}22，Ln/Exp 长度还要小于目标质数。没有 poly-new 的 Sqrt/Div/Eval。RSPE 需要输入确实足以确定递推，通常至少给 2 倍阶数的样本，不适用于任意序列外推。}
\par\textbf{考场版：}本节已换成普通迭代 NTT；fft 输出按自然顺序排列，invFft 接受同样顺序。不能与原版变换内核混搭。

\Needspace{6\baselineskip}\textbf{完整代码}
\begin{lstlisting}[firstnumber=1]
//
// template for ModInt, use x.pow(k) / x.inv()
// version 2.0 (Last Update Jun 13th, 2026)
//
// usage:
//   init(n); mint a,b; C(n,k); MOD=1e9+7
//
template <unsigned _M> struct ModInt{
static constexpr unsigned MOD=_M;
static_assert(1<MOD&&MOD<=INT_MAX);
unsigned x;
constexpr ModInt():x(0){}
constexpr ModInt(unsigned y):x(y%MOD){}
constexpr ModInt(int y):x((y%=static_cast<int>(MOD))<0?y+MOD:y){}
constexpr ModInt(unsigned long long y):x(y%MOD){}
constexpr ModInt(long long y):x((y%=static_cast<long long>(MOD))<0?y+MOD:y){}
ModInt& operator +=(const ModInt &a){x=(((x+=a.x)>=MOD)?x-MOD:x);return *this;}
ModInt& operator -=(const ModInt &a){x=(((x-=a.x)>=MOD)?x+MOD:x);return *this;}
ModInt& operator *=(const ModInt &a){x=static_cast<unsigned long long>(x)*a.x%MOD;return *this;}
ModInt& operator /=(const ModInt &a){return (*this)*=a.inv();}
bool operator ==(const ModInt &a)const{return x==a.x;}
bool operator !=(const ModInt &a)const{return x!=a.x;}
explicit operator bool() const { return x != 0; }
bool operator!() const { return x == 0; }
ModInt inv() const {
	unsigned a=MOD,b=x;ll y=0,z=1;
	while(b){
		const unsigned q=a/b,c=a-q*b;
		a=b,b=c;
		const ll w=y-(ll)q*z;
		y=z,z=w;
	}
	assert(a==1);return ModInt(y);
}
ModInt pow(long long k) const {
	unsigned long long e=k;
	ModInt a=*this,b=1;
	if(k<0){a=a.inv();e=0-e;}
	for(;e;e>>=1){if(e&1)b*=a;a*=a;}
	return b;
}
ModInt& operator++() {*this+=1;return *this;}
ModInt& operator--() {*this-=1;return *this;}
ModInt operator +()const{return *this;}
ModInt operator -()const{ModInt a;a.x=(x?MOD-x:0u);return a;}
ModInt operator +(const ModInt &a)const{return ModInt(*this)+=a;}
ModInt operator -(const ModInt &a)const{return ModInt(*this)-=a;}
ModInt operator *(const ModInt &a)const{return ModInt(*this)*=a;}
ModInt operator /(const ModInt &a)const{return ModInt(*this)/=a;}

template<typename T> ModInt friend operator +(T a,const ModInt &b){return ModInt(a)+=b;}
template<typename T> ModInt friend operator -(T a,const ModInt &b){return ModInt(a)-=b;}
template<typename T> ModInt friend operator *(T a,const ModInt &b){return ModInt(a)*=b;}
template<typename T> ModInt friend operator /(T a,const ModInt &b){return ModInt(a)/=b;}
friend istream& operator >> (istream &i,ModInt &x){ll y;if(i>>y)x=ModInt(y);return i;}
friend ostream& operator << (ostream &o,const ModInt &x){o<<x.x;return o;}
};
// Basic function, fac, ifac, binom
// init(n) prepares fac / ifac / inv through n
const int MOD=1e9+7;
using mint=ModInt<MOD>;
using poly=vector<mint>;
using Poly=poly;
vector<mint> fac{1},ifac{1},inv{0,1};
void init(int n=0){
	assert(0<=n&&n<MOD);
	int m=fac.size();if(n<m)return;
	n=max(n,(int)min(2ll*m,(ll)MOD-1));
	fac.resize(n+1);ifac.resize(n+1);inv.resize(max(n+1,2));
	for(int i=m;i<=n;i++){
		fac[i]=fac[i-1]*i;
		if(i>1)inv[i]=-(MOD/i)*inv[MOD%i];
		ifac[i]=ifac[i-1]*inv[i];
	}
}
inline mint C(int x,int y){
	// choose x from y
	if(x<0||y<x)return 0;
	init(y);return fac[y]*ifac[x]*ifac[y-x];
}
inline mint binom(int y,int x){return C(x,y);}

//
// template for MTT (three-prime CRT)
// usage:
//   poly h=f*g; Inv(f); Ln(f); Exp(f);
//   diff, integ, value, BM, FSPE, RSPE, lagrange;
//   MOD is prime; Ln needs f[0]=1, Exp needs f[0]=0
//   convolution length <= 2^23; FPS length <= 2^22
//
const int FFT_MAX=23;
template<unsigned MO> struct MTT_NTT{
using num=ModInt<MO>;
void ntt(num *a,int n,bool inverse)const{
	assert(1<=n&&n<=(1<<FFT_MAX)&&!(n&(n-1)));
	for(int i=1,j=0;i<n;i++){
		int k=n>>1;
		for(;j&k;k>>=1)j^=k;
		j^=k;if(i<j)swap(a[i],a[j]);
	}
	for(int len=2;len<=n;len<<=1){
		num step=num(3).pow((MO-1)/len);
		if(inverse)step=step.inv();
		for(int i=0;i<n;i+=len){
			num cur=1;
			for(int j=0;j<len/2;j++,cur*=step){
				num x=a[i+j],y=a[i+j+len/2]*cur;
				a[i+j]=x+y;a[i+j+len/2]=x-y;
			}
		}
	}
	if(inverse){num v=num(n).inv();for(int i=0;i<n;i++)a[i]*=v;}
}
void fft(num *a,int n)const{ntt(a,n,false);}
void invFft(num *a,int n)const{ntt(a,n,true);}
};

template<unsigned MO> vector<ModInt<MO>> convolution(const poly &f,const poly &g,int n){
	static const MTT_NTT<MO> ntt;
	vector<ModInt<MO>> a(n),b(n);
	for(int i=0;i<(int)f.size();i++)a[i]=f[i].x;
	for(int i=0;i<(int)g.size();i++)b[i]=g[i].x;
	ntt.fft(a.data(),n);ntt.fft(b.data(),n);
	for(int i=0;i<n;i++)a[i]*=b[i];
	ntt.invFft(a.data(),n);a.resize(f.size()+g.size()-1);return a;
}
poly operator *(const poly &f,const poly &g){
	if(f.empty()||g.empty())return {};
	assert(f.size()+g.size()-1<=(1<<FFT_MAX));
	int m=f.size()+g.size()-1;
	poly h(m);
	if(min(f.size(),g.size())<=32){
		for(int i=0;i<(int)f.size();i++)for(int j=0;j<(int)g.size();j++)h[i+j]+=f[i]*g[j];
		return h;
	}
	constexpr unsigned mod1=998244353,mod2=167772161,mod3=469762049;
	constexpr unsigned long long mod12=1ull*mod1*mod2;
	static const unsigned inv1=ModInt<mod2>(mod1).inv().x;
	static const unsigned inv2=ModInt<mod3>(mod12).inv().x;
	assert((__int128)min(f.size(),g.size())*(MOD-1)*(MOD-1)<(__int128)mod12*mod3);
	int n=1;while(n<m)n<<=1;
	auto a=convolution<mod1>(f,g,n);
	auto b=convolution<mod2>(f,g,n);
	auto c=convolution<mod3>(f,g,n);
	for(int i=0;i<m;i++){
		unsigned long long x=(b[i].x+mod2-a[i].x%mod2)*1ull*inv1%mod2;
		x=x*mod1+a[i].x;
		unsigned long long y=(c[i].x+mod3-x%mod3)*inv2%mod3;
		h[i]=x%MOD+(mod12%MOD)*y%MOD;
	}
	return h;
}
// Transposed multiplication: (f/g)[i] = sum_j f[i+j]*g[j].
poly operator /(const poly &f,poly g){
	if(f.empty()||g.empty())return poly(f.size());
	int m=g.size();reverse(g.begin(),g.end());g=f*g;
	return poly(g.begin()+m-1,g.end());
}
void operator +=(poly &f,const poly &g){
	if(f.size()<g.size())f.resize(g.size());
	for(int i=0;i<(int)g.size();i++)f[i]+=g[i];
}
poly operator +(poly f,const poly &g){f+=g;return f;}
void operator -=(poly &f,const poly &g){
	if(f.size()<g.size())f.resize(g.size());
	for(int i=0;i<(int)g.size();i++)f[i]-=g[i];
}
poly operator -(poly f,const poly &g){f-=g;return f;}
poly operator *(poly f,mint x){for(auto &y:f)y*=x;return f;}
poly operator *(mint x,poly f){for(auto &y:f)y*=x;return f;}
mint value(const poly &f,mint x){
	mint ans=0;for(int i=(int)f.size()-1;i>=0;i--)ans=ans*x+f[i];return ans;
}
poly Inv(const poly &f){
	assert(!f.empty()&&f[0]&&f.size()<=(1<<(FFT_MAX-1)));
	int n=f.size();poly g{f[0].inv()};
	while((int)g.size()<n){
		int m=min(n,(int)g.size()*2);
		poly h=poly(f.begin(),f.begin()+m)*g;h.resize(m);
		for(auto &x:h)x=-x;
		h[0]+=2;g=g*h;g.resize(m);
	}
	return g;
}
poly diff(poly f){
	if(f.empty())return {};
	for(int i=1;i<(int)f.size();i++)f[i-1]=f[i]*i;
	f.pop_back();return f;
}
poly integ(poly f){
	int n=f.size();init(n);f.resize(n+1);
	for(int i=n;i;i--)f[i]=f[i-1]*inv[i];
	f[0]=0;return f;
}
poly Ln(const poly &f){
	assert(!f.empty()&&f[0].x==1&&f.size()<(size_t)MOD&&f.size()<=(1<<(FFT_MAX-1)));
	poly g=diff(f)*Inv(f);g.resize(f.size()-1);return integ(move(g));
}
poly Exp(const poly &f){
	assert(!f.empty()&&!f[0]&&f.size()<(size_t)MOD&&f.size()<=(1<<(FFT_MAX-1)));
	int n=f.size();poly g{1};
	while((int)g.size()<n){
		int m=min(n,(int)g.size()*2);poly h=g;h.resize(m);h=Ln(h);
		for(int i=0;i<m;i++)h[i]=f[i]-h[i];
		h[0]+=1;g=g*h;g.resize(m);
	}
	return g;
}
poly BM(poly a) {
	poly C{1}, B{1};
	int L = 0, m = 1;
	mint b = 1;
	for (int n = 0; n < (int)a.size(); n++) {
		mint d = 0;
		for (int i = 0; i <= L; i++) d += C[i] * a[n - i];
		if (!d) {
			m++;
		} else {
			poly T = C;
			mint coef = d * b.inv();
			poly xmB(m, 0);
			for (mint val : B) xmB.push_back(val);
			if (C.size() < xmB.size()) C.resize(xmB.size(), 0);
			for (int i = 0; i < (int)xmB.size(); i++) C[i] -= coef * xmB[i];
			if (2 * L <= n) {
				L = n + 1 - L; B = T; b = d; m = 1;
			} else {
				m++;
			}
		}
	}
	return C;
}

mint FSPE(poly F, poly G, ll t) {
	// find [x^t] F/G
	assert(t>=0&&!G.empty()&&G[0]);
	if(F.empty())return 0;
	while (t) {
		poly G1 = G;
		for (int i = 1; i < (int)G1.size(); i += 2) G1[i] = -G1[i];
		F = F * G1; G = G * G1;
		poly f, g;
		for (int i = t & 1; i < (int)F.size(); i += 2) f.push_back(F[i]);
		for (int i = 0; i < (int)G.size(); i += 2) g.push_back(G[i]);
		t >>= 1;
		F = f; G = g;
		if(F.empty())return 0;
	}
	return F[0] / G[0];
}

mint RSPE(poly f, ll x) {
	// find f[x] in O(n^2), note that f must be recursion
	poly g = BM(f);
	poly s(g.size()-1);
	for (int i = 0; i < (int)s.size(); i++)
		for (int j = 0; j <= i && j < (int)f.size(); j++)
			s[i] += f[j] * g[i - j];
	return FSPE(s, g, x);
}


poly lagrange(const vector<mint> &x,const vector<mint> &y){
	int n=x.size();assert(y.size()==x.size());
	poly g(n+1),f(n),q(n);g[0]=1;
	for(int i=0;i<n;i++){
		for(int j=i+1;j;j--)g[j]=g[j-1]-x[i]*g[j];
		g[0]*=-x[i];
	}
	for(int i=0;i<n;i++){
		q[n-1]=g[n];
		for(int j=n-2;j>=0;j--)q[j]=g[j+1]+x[i]*q[j+1];
		mint d=value(q,x[i]);assert(d);
		mint c=y[i]/d;
		for(int j=0;j<n;j++)f[j]+=q[j]*c;
	}
	return f;
}
\end{lstlisting}
\par\vspace{1em}

\Needspace{10\baselineskip}
\subsection{\texttt{polynomial/FWT.cpp}}
\textbf{按位 AND / OR / XOR 卷积}

\TemplUsage{求 $c_k=\sum_{i\mathbin{\mathrm{op}}j=k}a_i b_j$，其中 op 为按位 AND、OR 或 XOR。长度 $n=2^k$，每次变换 $O(n\log n)$，额外空间 $O(1)$。}
\TemplApi{\begin{itemize}
\item \texttt{convolution\_and(a,b)}、\texttt{convolution\_or(a,b)}、\texttt{convolution\_xor(a,b)} 直接返回对应卷积，不修改传入数组。支持长度不同、非二次幂长度；任一输入为空则返回空数组，否则结果长度为不小于两者最大长度的最小二次幂。
\item \texttt{fwt\_and(a)}、\texttt{fwt\_or(a)}、\texttt{fwt\_xor(a)} 原地正变换；第二个参数传 \texttt{true} 做逆变换。
\item 支持 \texttt{vector<int>}、\texttt{vector<ll>} 和 \texttt{vector<mint>}。mint 需先抄模整数定义，FWT 本身不依赖多项式文件。
\item 两个数组补零到相同的非零二次幂长度 n，使 n 大于最大下标；不必补到两者长度之和。
\item 两边做同一种正变换，逐点相乘，再做对应逆变换，即得卷积。
\end{itemize}}
\textbf{最小示例（放入 main）：}
\begin{lstlisting}
vector<int> a{1,2},b{3,4};
auto c=convolution_and(a,b); assert((c==vector<int>{13,8}));
c=convolution_or(a,b); assert((c==vector<int>{3,18}));
c=convolution_xor(a,b); assert((c==vector<int>{11,10}));
// vector<mint> works with the same interfaces.
\end{lstlisting}
\TemplNote{下标从 0 开始。int/ll 表示普通整数运算，不会自动取模，正变换和逐点乘法的中间值也必须不溢出。整数 XOR 逆变换逐层除以 2，要求每步可整除，不能把任意整数数组当成合法频域数据。XOR 模运算要求 2 可逆，通常使用奇数模数；AND/OR 无此限制。不需要归一化第二次。}
\Needspace{6\baselineskip}\textbf{完整代码}
\begin{lstlisting}[firstnumber=1]
//
// template for FWT and bitwise convolution
//
// usage:
//   auto c=convolution_and(a,b); // convolution_or / convolution_xor
//   fwt_and(a); fwt_and(a,true); // forward / inverse; also or / xor
//   vector<int>, vector<ll> or vector<mint>; 0-indexed
//   convolution pads to a power of 2; empty input returns {}
//   transform size must be a positive power of 2
//   integer intermediates must fit the type; mint XOR needs inverse of 2
//
template<typename T>
void fwt_and(vector<T> &a,bool inverse=false){
	assert(!a.empty()&&a.size()<=INT_MAX&&!(a.size()&(a.size()-1)));
	int n=a.size();
	for(int k=1;k<n;k<<=1)
		for(int i=0;i<n;i+=k<<1)
			for(int j=0;j<k;j++){
				if(inverse)a[i+j]-=a[i+j+k];
				else a[i+j]+=a[i+j+k];
			}
}
template<typename T>
void fwt_or(vector<T> &a,bool inverse=false){
	assert(!a.empty()&&a.size()<=INT_MAX&&!(a.size()&(a.size()-1)));
	int n=a.size();
	for(int k=1;k<n;k<<=1)
		for(int i=0;i<n;i+=k<<1)
			for(int j=0;j<k;j++){
				if(inverse)a[i+j+k]-=a[i+j];
				else a[i+j+k]+=a[i+j];
			}
}
template<typename T>
void fwt_xor(vector<T> &a,bool inverse=false){
	assert(!a.empty()&&a.size()<=INT_MAX&&!(a.size()&(a.size()-1)));
	int n=a.size();T inv2=1;
	if constexpr(!is_integral_v<T>)if(inverse)inv2=T(1)/T(2);
	for(int k=1;k<n;k<<=1)
		for(int i=0;i<n;i+=k<<1)
			for(int j=0;j<k;j++){
				if constexpr(is_integral_v<T>){
					using W=conditional_t<(sizeof(T)<=4),long long,__int128>;
					W x=a[i+j],y=a[i+j+k],u=x+y,v=x-y;
					if(inverse){assert(u%2==0&&v%2==0);u/=2;v/=2;}
					assert(numeric_limits<T>::lowest()<=u&&u<=numeric_limits<T>::max());
					assert(numeric_limits<T>::lowest()<=v&&v<=numeric_limits<T>::max());
					a[i+j]=T(u);a[i+j+k]=T(v);
				}else{
					T x=a[i+j],y=a[i+j+k];
					a[i+j]=x+y;a[i+j+k]=x-y;
					if(inverse){a[i+j]*=inv2;a[i+j+k]*=inv2;}
				}
			}
}

template<typename T>
vector<T> convolution_and(vector<T> a,vector<T> b){
	if(a.empty()||b.empty())return {};
	assert(max(a.size(),b.size())<=(1u<<30));
	int n=1;while(n<(int)max(a.size(),b.size()))n<<=1;
	a.resize(n);b.resize(n);
	fwt_and(a);fwt_and(b);
	for(int i=0;i<n;i++)a[i]*=b[i];
	fwt_and(a,true);return a;
}

template<typename T>
vector<T> convolution_or(vector<T> a,vector<T> b){
	if(a.empty()||b.empty())return {};
	assert(max(a.size(),b.size())<=(1u<<30));
	int n=1;while(n<(int)max(a.size(),b.size()))n<<=1;
	a.resize(n);b.resize(n);
	fwt_or(a);fwt_or(b);
	for(int i=0;i<n;i++)a[i]*=b[i];
	fwt_or(a,true);return a;
}

template<typename T>
vector<T> convolution_xor(vector<T> a,vector<T> b){
	if(a.empty()||b.empty())return {};
	assert(max(a.size(),b.size())<=(1u<<30));
	int n=1;while(n<(int)max(a.size(),b.size()))n<<=1;
	a.resize(n);b.resize(n);
	fwt_xor(a);fwt_xor(b);
	for(int i=0;i<n;i++)a[i]*=b[i];
	fwt_xor(a,true);return a;
}
\end{lstlisting}
\clearpage
\section{字符串}
\Needspace{10\baselineskip}
\subsection{\texttt{string/sa.cpp}}
\textbf{后缀数组和 O(1) LCP}

\TemplUsage{计数排序倍增构建后缀排序；附带 SA\_LCP，用 height 数组和 ST 表查询两个后缀的最长公共前缀。}
\TemplApi{\begin{itemize}
\item \texttt{SA sa; sa.build(" "+s)} 输入前面补一个占位空格，真实字符下标 1..n。
\item \texttt{sa.sa[r]} 是排名第 r 的后缀起点；\texttt{sa.rk[i]} 是后缀 i 的排名。
\item \texttt{SA\_LCP l; l.build(sa)}，然后 \texttt{l.lcp(i,j)} 或 \texttt{l.ask(i,j)} 查询，i,j 是原串位置。
\item 也可 \texttt{l.build(" "+s)} 直接建。height[r] 为排名 r 与 r-1 的后缀 LCP，height[1]=0。
\end{itemize}}
\textbf{最小示例（放入 main）：}
\begin{lstlisting}
SA sa; sa.build(" banana");
assert(sa.sa[1]==6);
SA_LCP l; l.build(sa);
assert(l.lcp(2,4)==3); // anana / ana
assert(l.lcp(3,3)==4);
\end{lstlisting}
\TemplNote{SA 构建 O(n log n)，工作空间 O(n)，LCP 结构构建与空间 O(n log n)，查询 O(1)。不要把后缀排名传给 lcp。默认字符编码范围 1..127，任意非零字节传 M=256；不支持串内零字节。修改原串后须重建两者。}

\Needspace{6\baselineskip}\textbf{完整代码}
\begin{lstlisting}[firstnumber=1]
struct SA{
int n=0,m=0;
vector<int> sa,rk;
vector<unsigned char> s;
void set(){n=m=0;sa.clear();rk.clear();s.clear();}
void build(string S,int M=128){
	assert(!S.empty()&&S[0]==' '&&1<=M&&M<=256);
	assert(S.size()<INT_MAX/2);
	n=S.size()-1;m=M;
	sa.assign(n+1,0);rk.assign(n+1,0);s.assign(n+1,0);
	vector<int> y(n+1),old(n+1),cnt(max(n+1,M),0);
	for(int i=1;i<=n;i++){
		s[i-1]=(unsigned char)S[i];
		assert(0<s[i-1]&&s[i-1]<M);
		rk[i]=s[i-1];cnt[rk[i]]++;
	}
	for(int i=1;i<m;i++)cnt[i]+=cnt[i-1];
	for(int i=n;i;i--)sa[cnt[rk[i]]--]=i;
	for(int k=1;k<n;k<<=1){
		int p=0;
		for(int i=n-k+1;i<=n;i++)y[++p]=i;
		for(int i=1;i<=n;i++)if(sa[i]>k)y[++p]=sa[i]-k;
		fill(cnt.begin(),cnt.end(),0);
		for(int i=1;i<=n;i++)cnt[rk[i]]++;
		for(int i=1;i<m;i++)cnt[i]+=cnt[i-1];
		for(int i=n;i;i--)sa[cnt[rk[y[i]]]--]=y[i];
		old=rk;p=0;
		for(int i=1;i<=n;i++){
			int x=sa[i],z=sa[i-1];
			if(i==1||old[x]!=old[z]||
				(x+k<=n?old[x+k]:0)!=(z+k<=n?old[z+k]:0))p++;
			rk[x]=p;
		}
		m=p+1;if(p==n)break;
	}
	for(int i=1;i<=n;i++)rk[sa[i]]=i;
}
void output(){
	cerr<<"string: ";for(int i=0;i<n;i++)cerr<<s[i];cerr<<'\n';
	cerr<<"sa: ";for(int i=1;i<=n;i++)cerr<<sa[i]<<" \n"[i==n];
	cerr<<"rk: ";for(int i=1;i<=n;i++)cerr<<rk[i]<<" \n"[i==n];
}
void solve(){for(int i=1;i<=n;i++)cout<<sa[i]<<" \n"[i==n];}
};
// SA_LCP L; L.build(sa); L.lcp(i,j); // 1-indexed
// O(n log n) build, O(1) query
struct SA_LCP{
int n=0;
vector<int> rk,height;
vector<vector<int>> st;
void build(const SA &S){
	n=S.n;rk=S.rk;height.assign(n+1,0);st.clear();
	assert((int)S.sa.size()==n+1&&(int)rk.size()==n+1);
	for(int i=1,k=0;i<=n;i++){
		int r=rk[i];
		if(r==1){k=0;continue;}
		int j=S.sa[r-1];
		while(i+k<=n&&j+k<=n&&S.s[i+k-1]==S.s[j+k-1])k++;
		height[r]=k;if(k)k--;
	}
	if(!n)return;
	st.push_back(height);
	for(int k=1;(1<<k)<=n;k++){
		st.emplace_back(n-(1<<k)+2);
		for(int i=1;i+(1<<k)-1<=n;i++)st[k][i]=min(st[k-1][i],st[k-1][i+(1<<(k-1))]);
	}
}
void build(const string &s,int M=128){SA S;S.build(s,M);build(S);}
int lcp(int i,int j)const{
	assert(1<=min(i,j)&&max(i,j)<=n);
	if(i==j)return n-i+1;
	int l=rk[i],r=rk[j];if(l>r)swap(l,r);l++;
	int k=__lg(r-l+1);return min(st[k][l],st[k][r-(1<<k)+1]);
}
int ask(int i,int j)const{return lcp(i,j);}
};
\end{lstlisting}
\par\vspace{1em}

\Needspace{10\baselineskip}
\subsection{\texttt{string/sam.cpp}}
\textbf{单串后缀自动机}

\TemplUsage{构建单个小写字母串的子串状态图，并统计每个状态对应子串的出现次数。时间及空间 O(n)，字母表固定 26。}
\TemplApi{\begin{itemize}
\item \texttt{SAM t; t.build(" "+s)} 每次重新构建。
\item 根为 0；\texttt{ch[u][c]} 为字符 c 的转移，0 表示不存在。
\item 状态 u 对应长度区间 (len[fail[u]],len[u]]，同一状态内的子串出现次数均为 siz[u]。
\end{itemize}}
\textbf{最小示例（放入 main）：}
\begin{lstlisting}
SAM t; t.build(" ababa");
int u=0; for(char c:string("aba"))u=t.ch[u][c-'a'];
assert(u && t.siz[u]==2);
ll distinct=0;
for(int i=1;i<=t.tot;i++)distinct+=t.len[i]-t.len[t.fail[i]];
assert(distinct==9);
\end{lstlisting}
\TemplNote{匹配模式串时遇到不存在转移应立即判定不出现。build 内已累计 siz，不能再次累计。裸调用 extend 不会自动维护整套出现次数；多串应用请用 GSAM。solve 是“重复子串长度乘次数最大值”的特定示例，不是通用查询。}

\Needspace{6\baselineskip}\textbf{完整代码}
\begin{lstlisting}[firstnumber=1]
// string algorithm template by Otomachi Una (Junlin Ye)
// version 1.0 (last update May 24th 2026)
//
// usage:
//   SAM sam; sam.build(" abaab");  // leading space, lowercase
//   sam.len[u], sam.fail[u], sam.siz[u]
//
// your string should index from 1, begin with a space
struct SAM{
// this SAM rooted at 0
string s;
int n=0,tot=0,lst=0;
vector<array<int,26>> ch;
vector<int> fail,len,siz;
vector<vector<int>> son;
SAM(){set();}
void set(int m=0){
	assert(0<=m&&m<(INT_MAX-1)/2);n=tot=lst=0;
	ch.assign(1,{});fail.assign(1,-1);len.assign(1,0);siz.assign(1,0);son.assign(1,{});
	ch.reserve(2*m+1);fail.reserve(2*m+1);len.reserve(2*m+1);
	siz.reserve(2*m+1);son.reserve(2*m+1);
}
int new_node(){
	assert(tot<INT_MAX-1);ch.push_back({});fail.push_back(0);
	len.push_back(0);siz.push_back(0);son.emplace_back();return ++tot;
}
inline int clone(int x){
	new_node();ch[tot]=ch[x];
	fail[tot]=fail[x];
	return tot;
}
void extend(char c){
	assert('a'<=c&&c<='z');c-='a'; // attention the alphabet
	int cur=new_node();len[cur]=len[lst]+1;
	while(~lst&&!ch[lst][c]) ch[lst][c]=tot,lst=fail[lst];
	if(lst==-1) fail[cur]=0;
	else{
		int p=ch[lst][c];
		if(len[p]==len[lst]+1) fail[cur]=p;
		else{
			int q=clone(p);len[q]=len[lst]+1;
			while(~lst&&ch[lst][c]==p) ch[lst][c]=q,lst=fail[lst];
			fail[p]=fail[cur]=q;
		}
	}
	lst=cur;
}
void output(){
	// this allow you to see the structure of SAM
	// output fail, ed, len etc.
	cerr<<"string:"<<s<<endl;
	cerr<<"# node = "<<tot<<endl;
	cerr<<"fail: ";for(int i=0;i<=tot;i++) cerr<<fail[i]<<" \n"[i==tot];
	cerr<<"len: ";for(int i=0;i<=tot;i++) cerr<<len[i]<<" \n"[i==tot];
	// other else?
	cerr<<"siz: ";for(int i=0;i<=tot;i++) cerr<<siz[i]<<" \n"[i==tot];
}
void dfs(int u){
	vector<int> q{u};
	for(int i=0;i<(int)q.size();i++)for(int v:son[q[i]])q.push_back(v);
	for(int i=(int)q.size()-1;i>0;i--)siz[fail[q[i]]]+=siz[q[i]];
}
void build(string S){
	assert(!S.empty()&&S[0]==' '&&S.size()<(INT_MAX-1)/2);
	set((int)S.size()-1);s=move(S);n=(int)s.size()-1;
	for(int i=1;i<=n;i++){
		extend(s[i]);
		// something record here
		siz[lst]=1;
	}
	for(int i=1;i<=tot;i++) son[fail[i]].push_back(i);
	dfs(0);
	// output();
}
void solve(){
	// this is the main work
	ll ans=0;
	for(int i=1;i<=tot;i++) if(siz[i]>1) ans=max(ans,1ll*siz[i]*len[i]);
	cout<<ans<<'\n';
}
};
\end{lstlisting}
\par\vspace{1em}

\Needspace{10\baselineskip}
\subsection{\texttt{string/ac.cpp}}
\textbf{多模式串 AC 自动机}

\TemplUsage{一次加入多个模式串，再统计它们在文本中的出现次数。允许重叠出现；构建 O(26*节点数)，每次 count 为 O(文本长度+节点数)。}
\TemplApi{\begin{itemize}
\item \texttt{AC ac; id=ac.insert(" "+pattern)} 返回该模式串末节点编号。
\item 插入完所有模式串后调用 build，之后不能再插入，除非 set 清空重建。
\item \texttt{auto cnt=ac.count(" "+text)}，模式串答案是 cnt[id]。
\item build 后 ch 是补全后的转移，可直接用于自动机 DP；fail 是失配链接。
\end{itemize}}
\textbf{最小示例（放入 main）：}
\begin{lstlisting}
AC ac; int a=ac.insert(" aba"),b=ac.insert(" ba");
ac.build(); auto cnt=ac.count(" ababa");
assert(cnt[a]==2 && cnt[b]==2);
\end{lstlisting}
\TemplNote{模式和文本都只支持小写字母，前导空格不参加匹配。重复模式串返回同一编号，但它们各自的答案仍是该模式在文本中的出现次数。空模式返回根编号 0，次数为文本长度+1。不同 count 之间不会累计。}

\Needspace{6\baselineskip}\textbf{完整代码}
\begin{lstlisting}[firstnumber=1]
//
// template for Aho-Corasick
// usage:
//   AC ac; int id=ac.insert(" aba"); ac.build();
//   auto cnt=ac.count(" ababa"); cnt[id];
//   leading space, lowercase; insert before build
//
struct AC{
int tot=0;
bool built=false;
vector<array<int,26>> ch;
vector<int> fail,q;
AC(){set();}
void set(int m=0){
	assert(0<=m&&m<INT_MAX-1);tot=0;built=false;
	ch.assign(1,{});fail.assign(1,0);q.clear();
	ch.reserve(m+1);fail.reserve(m+1);
}
int new_node(){
	assert(tot<INT_MAX-1);ch.push_back({});fail.push_back(0);return ++tot;
}
int insert(const string &s){
	assert(!built&&!s.empty()&&s[0]==' '&&s.size()<INT_MAX);
	int u=0;
	for(int i=1;i<(int)s.size();i++){
		assert('a'<=s[i]&&s[i]<='z');int c=s[i]-'a';
		if(!ch[u][c]){int v=new_node();ch[u][c]=v;}
		u=ch[u][c];
	}
	return u;
}
void build(){
	if(built)return;
	q.clear();q.reserve(tot+1);q.push_back(0);
	for(int i=0;i<(int)q.size();i++){
		int u=q[i];
		for(int c=0;c<26;c++){
			int v=ch[u][c];
			if(v){fail[v]=u?ch[fail[u]][c]:0;q.push_back(v);}
			else if(u)ch[u][c]=ch[fail[u]][c];
		}
	}
	built=true;
}
vector<ll> count(const string &s)const{
	assert(built&&!s.empty()&&s[0]==' '&&s.size()<INT_MAX);
	vector<ll> cnt(tot+1);cnt[0]=1;int u=0;
	for(int i=1;i<(int)s.size();i++){
		assert('a'<=s[i]&&s[i]<='z');u=ch[u][s[i]-'a'];cnt[u]++;
	}
	for(int i=(int)q.size()-1;i;i--)cnt[fail[q[i]]]+=cnt[q[i]];
	return cnt;
}
};
\end{lstlisting}
\par\vspace{1em}

\Needspace{10\baselineskip}
\subsection{\texttt{string/gsam.cpp}}
\textbf{广义后缀自动机}

\TemplUsage{维护多个小写串的所有子串；出现次数是这些串中次数之和，重复插入的串也重复计数，不允许跨串边界匹配。}
\TemplApi{\begin{itemize}
\item \texttt{GSAM t; t.insert(" "+s)} 加入一条字符串，之后 \texttt{t.build()} 累计出现次数。
\item \texttt{count(" "+p)} 返回模式的总出现次数；\texttt{find(" "+p)} 返回状态编号，不存在为 -1。
\item \texttt{distinct()} 返回所有插入串合起来的不同非空子串数。
\item 可继续 insert，但下一次 count 前必须重新 build。
\end{itemize}}
\textbf{最小示例（放入 main）：}
\begin{lstlisting}
GSAM t; t.insert(" aba"); t.insert(" bab");
t.build(); assert(t.count(" ab")==2);
assert(t.count(" abab")==0 && t.distinct()==6);
t.insert(" aba"); t.build(); assert(t.count(" aba")==2);
\end{lstlisting}
\TemplNote{count 不是“包含该子串的字符串个数”。根为 0，fail[0]=-1；别把普通 SAM 简单重置 lst 的做法当成这里的插入实现。set(m) 的 m 用于预留总字符量级的空间，build 的计数按当前全部插入内容重新计算。}

\Needspace{6\baselineskip}\textbf{完整代码}
\begin{lstlisting}[firstnumber=1]
//
// template for Generalized SAM
// usage:
//   GSAM sam; sam.insert(" aba"); sam.insert(" bab"); sam.build();
//   sam.count(" ab"); sam.distinct();
//   leading space, lowercase; count includes repeated strings
//
struct GSAM{
int n=0,tot=0;
ll strings=0;
bool built=false;
vector<array<int,26>> ch;
vector<int> fail,len;
vector<ll> cnt,siz;
GSAM(){set();}
void set(int m=0){
	assert(0<=m&&m<(INT_MAX-1)/2);n=tot=0;strings=0;built=false;
	ch.assign(1,{});fail.assign(1,-1);len.assign(1,0);cnt.assign(1,0);siz.clear();
	ch.reserve(2*m+1);fail.reserve(2*m+1);len.reserve(2*m+1);cnt.reserve(2*m+1);
}
int new_node(){
	assert(tot<INT_MAX-1);ch.push_back({});fail.push_back(0);len.push_back(0);cnt.push_back(0);return ++tot;
}
int clone(int x,int l){
	int y=new_node();ch[y]=ch[x];fail[y]=fail[x];len[y]=l;return y;
}
int extend(int p,int c){
	if(ch[p][c]){
		int q=ch[p][c];if(len[q]==len[p]+1)return q;
		int r=clone(q,len[p]+1);
		while(p!=-1&&ch[p][c]==q)ch[p][c]=r,p=fail[p];
		fail[q]=r;return r;
	}
	int cur=new_node();len[cur]=len[p]+1;
	while(p!=-1&&!ch[p][c])ch[p][c]=cur,p=fail[p];
	if(p==-1)fail[cur]=0;
	else{
		int q=ch[p][c];
		if(len[q]==len[p]+1)fail[cur]=q;
		else{
			int r=clone(q,len[p]+1);
			while(p!=-1&&ch[p][c]==q)ch[p][c]=r,p=fail[p];
			fail[q]=fail[cur]=r;
		}
	}
	return cur;
}
int insert(const string &s){
	assert(!s.empty()&&s[0]==' '&&s.size()<INT_MAX);
	int u=0;n=max(n,(int)s.size()-1);strings++;built=false;
	for(int i=1;i<(int)s.size();i++){
		assert('a'<=s[i]&&s[i]<='z');u=extend(u,s[i]-'a');cnt[u]++;
	}
	return u;
}
void build(){
	vector<int> c(n+1),q(tot+1);siz=cnt;
	for(int i=0;i<=tot;i++)c[len[i]]++;
	for(int i=1;i<=n;i++)c[i]+=c[i-1];
	for(int i=0;i<=tot;i++)q[--c[len[i]]]=i;
	for(int i=tot;i;i--)siz[fail[q[i]]]+=siz[q[i]];
	siz[0]+=strings;built=true;
}
int find(const string &s)const{
	assert(!s.empty()&&s[0]==' '&&s.size()<INT_MAX);int u=0;
	for(int i=1;i<(int)s.size();i++){
		assert('a'<=s[i]&&s[i]<='z');u=ch[u][s[i]-'a'];if(!u)return -1;
	}
	return u;
}
ll count(const string &s)const{
	assert(built);int u=find(s);return u==-1?0:siz[u];
}
ll distinct()const{
	ll ans=0;for(int i=1;i<=tot;i++)ans+=len[i]-len[fail[i]];return ans;
}
};
\end{lstlisting}
\par\vspace{1em}

\Needspace{10\baselineskip}
\subsection{\texttt{string/kmp.cpp}}
\textbf{前缀函数和 Z 函数}

\TemplUsage{线性时间计算字符串的 border 与各后缀和全串的公共前缀长度。}
\TemplApi{\begin{itemize}
\item \texttt{border(" "+s)} 返回 pi，pi[i] 为 s[1..i] 的最长真 border 长度。
\item \texttt{z\_function(" "+s)} 返回 z，z[i] 为 s[i..n] 与 s[1..n] 的 LCP。
\item 数组长度 n+1，0 号是占位。
\end{itemize}}
\textbf{最小示例（放入 main）：}
\begin{lstlisting}
auto pi=border(" ababa"); auto z=z_function(" ababa");
assert(pi[5]==3 && z[3]==3);
\end{lstlisting}
\TemplNote{这个实现 z[1]=0，而不是全串长度。没有封装“模式串在文本的所有位置”：可自行 KMP 扫描，或拼接 pattern+分隔符+text 后用 Z；分隔符不能出现在两个原串里。}

\Needspace{6\baselineskip}\textbf{完整代码}
\begin{lstlisting}[firstnumber=1]
// string algorithm template by Otomachi Una (Junlin Ye)
// version 1.0 (last update May 24th 2026)
//
// usage:
//   string s = " ababa"; auto b = border(s); auto z = z_function(s);
//   1-indexed with leading space
//
// your string should index from 1, begin with a space
vector<int> border(string s){
	// s is a string with length n+1,
	// this function return a vector with length n+1
	// indication the border array of s
	assert(!s.empty()&&s.size()<=INT_MAX);
	int n=(int)s.length()-1;
	assert(!s.empty()&&s[0]==' ');
	vector<int> a(n+1);
	for(int i=2,j=0;i<=n;i++){
		while(j&&s[j+1]!=s[i]) j=a[j];
		if(s[j+1]==s[i]) j++;
		a[i]=j;
	}
	return a;
}
vector<int> z_function(string s){
	// s is a string with length n+1,
	// this function return a vector with length n+1
	// indication the a[i]=lcp(s[i:],s[])
	assert(!s.empty()&&s.size()<=INT_MAX);
	int n=(int)s.length()-1;
	assert(!s.empty()&&s[0]==' ');
	vector<int> z(n+1);
	for(int i=2,l=1,r=0;i<=n;i++){
		z[i]=r<i?0:min(r-i+1,z[i-l+1]);
		while(i+z[i]<=n&&s[1+z[i]]==s[i+z[i]]) z[i]++;
		if(i+z[i]>r) l=i,r=i+z[i]-1;
	}
	return z;
}
\end{lstlisting}
\par\vspace{1em}

\Needspace{10\baselineskip}
\subsection{\texttt{string/hash.cpp}}
\textbf{双模前缀哈希}

\TemplUsage{提供子串指纹与拼接，不是无碰撞的字符串判等。构建前缀数组 O(n)，常规子串取值 O(1)。}
\TemplApi{\begin{itemize}
\item \texttt{auto h=Hash\_of(" "+s)} 产生前缀哈希。
\item 子串 [l,r] 为 h[r]-h[l-1]。
\item \texttt{a+b} 拼接两段，\texttt{h.v()} 打包哈希值；\texttt{.l} 是长度。
\item \texttt{init\_hash(n)} 可预留幂次；实际操作也会自动扩容。
\end{itemize}}
\textbf{最小示例（放入 main）：}
\begin{lstlisting}
auto h=Hash_of(" ababa");
auto a=h[3]-h[0],b=h[5]-h[2];
assert(a.l==b.l && a.v()==b.v());
\end{lstlisting}
\TemplNote{比较子串时同时比较长度和哈希。底数每个进程随机，不能把数值用于跨进程永久存储。减法只表示删去前缀，不是任意删除；Hash 对象应使用 \{\} 零初始化。}

\Needspace{6\baselineskip}\textbf{完整代码}
\begin{lstlisting}[firstnumber=1]
//
// template for Hash, prepared by Otomachi_Una
// version 1.2 (Last Update Jun 13th, 2026)
//
// usage:
//   init_hash(); auto H = Hash_of(" abc"); Hash sub = H[r]-H[l-1];
//   sub.v() for compare; random bases each run
//
mt19937 rnd(time(0));
const int MOD1 = 998244353;
const int MOD2 = 1e9 + 7;
const int B1 = 256 + rnd() % (MOD1-257);
const int B2 = 256 + rnd() % (MOD2-257);

vector<ll> pw1{1}, pw2{1};
void init_hash(int n=0);

struct Hash {
	int l=0;
	ll h1=0,h2=0;
	ll v() const { return h1 << 30 | h2; }
};

inline Hash operator+(const Hash &x, const Hash &y) {
	// append y to end of x
	assert(x.l>=0&&y.l>=0&&x.l<=INT_MAX-y.l);init_hash(y.l);
	return {x.l + y.l,
			(x.h1 * pw1[y.l] + y.h1) % MOD1,
			(x.h2 * pw2[y.l] + y.h2) % MOD2};
}

inline Hash operator-(const Hash &x, const Hash &y) {
	// y is prefix of x, delete y from x
	assert(0<=y.l&&y.l<=x.l);init_hash(x.l-y.l);
	return {x.l - y.l,
			(x.h1 - y.h1 * pw1[x.l - y.l] % MOD1 + MOD1) % MOD1,
			(x.h2 - y.h2 * pw2[x.l - y.l] % MOD2 + MOD2) % MOD2};
}

template<typename T>
inline Hash operator+(const Hash &h, const T &x) {
	// append x to end of h
	assert(0<=h.l&&h.l<INT_MAX);
	return {h.l + 1,
			(h.h1 * B1 + x) % MOD1,
			(h.h2 * B2 + x) % MOD2};
}

template<typename T>
inline void operator+=(Hash &h, const T &x) {
	// append x to end of h
	h = h + x;
}

void init_hash(int n) {
	assert(0<=n&&n<INT_MAX);
	int m=pw1.size();if(n<m)return;
	n=max(n,(int)min(2ll*m,(ll)INT_MAX-1));
	pw1.resize(n+1);pw2.resize(n+1);
	for(int i=m;i<=n;i++){
		pw1[i]=pw1[i-1]*B1%MOD1;pw2[i]=pw2[i-1]*B2%MOD2;
	}
}

vector<Hash> Hash_of(const string &s) {
	int n = (int)s.length() - 1;
	assert(!s.empty()&&s[0]==' '&&s.size()<=INT_MAX);init_hash(n);
	vector<Hash> a(n+1);
	for (int i = 1; i <= n; ++i)
		a[i] = a[i-1] + (unsigned char)s[i];
	return a;
}
\end{lstlisting}
\par\vspace{1em}

\Needspace{10\baselineskip}
\subsection{\texttt{string/pam.cpp}}
\textbf{回文自动机}

\TemplUsage{在线维护不同回文子串、最长回文后缀及回文出现次数。字符集为小写字母，建树与空间均为 $O(n)$。}
\TemplApi{\begin{itemize}
\item \texttt{PAM p; p.build(s)} 从整串重建；s 带前导空格。也可 \texttt{p.set(n)} 预留容量，再逐字 \texttt{p.extend(c)}，返回当前最长回文后缀结点。
\item 0 是长度 0 的根，1 是长度 -1 的根；真正的回文结点为 2..tot，不同回文数为 tot-1。\texttt{lst} 是当前最长回文后缀。
\item \texttt{len[u]} 为长度，\texttt{fail[u]} 为最长真回文后缀；\texttt{num[u]} 为该回文的非空回文后缀个数。因此每次加入后的 \texttt{num[lst]} 是当前前缀末尾的回文数。
\item \texttt{pos[u]} 为首次出现的右端点，首次区间是 [pos[u]-len[u]+1,pos[u]]。\texttt{cnt=p.count()} 返回每个回文的总出现次数，可重叠；根的次数置 0。
\end{itemize}}
\textbf{最小示例（放入 main）：}
\begin{lstlisting}
PAM p; p.build(" ababa");
auto cnt=p.count();
assert(p.tot-1==5 && p.len[p.lst]==5);
assert(p.num[p.lst]==3 && cnt[p.lst]==1);
int u=p.ch[1]['a'-'a']; assert(cnt[u]==3);
p.extend('b'); cnt=p.count(); // still safe after count()
assert(p.len[p.lst]==5 && cnt[u]==3);
\end{lstlisting}
\TemplNote{\texttt{siz} 只记录结点作为最长后缀的次数，不能直接作为总出现次数。count 在副本上沿 fail 累加，单次 $O(n)$，可以重复调用，也能继续 extend；不要每加入一个字符都调用 count。所有回文出现次数之和用 ll 保存。空串写作一个空格。}

\Needspace{6\baselineskip}\textbf{完整代码}
\begin{lstlisting}[firstnumber=1]
////////////////////////////////////////////////////////////////
//
// template for Palindromic Automaton
// usage:
//   PAM pam; pam.build(" ababa"); auto cnt=pam.count();
//   Or pam.set(n), then pam.extend(c); returns the longest suffix node.
//   Leading space, lowercase; real nodes are 2..tot, count() is repeatable.
//   len / fail / num / pos: length, suffix link, suffix count, first end.
//   O(n) build / space for a fixed alphabet, O(n) count.
//
////////////////////////////////////////////////////////////////
struct PAM{
int n=0,tot=1,lst=0;
string s;
vector<array<int,26>> ch;
vector<int> len,fail,num,pos,siz;
PAM(){set();}
void set(int m=0){
	assert(0<=m&&m<INT_MAX-2);n=0;tot=1;lst=0;s=" ";
	ch.assign(2,{});len={0,-1};fail={1,1};
	num.assign(2,0);pos.assign(2,0);siz.assign(2,0);
	s.reserve(m+1);ch.reserve(m+2);len.reserve(m+2);fail.reserve(m+2);
	num.reserve(m+2);pos.reserve(m+2);siz.reserve(m+2);
}
int get_fail(int u)const{
	while(s[n-len[u]-1]!=s[n])u=fail[u];
	return u;
}
int extend(char c){
	assert('a'<=c&&c<='z'&&n<INT_MAX-3);s+=c;n++;c-='a';
	int u=get_fail(lst);
	if(!ch[u][c]){
		int v=++tot;
		ch.push_back({});len.push_back(len[u]+2);
		fail.push_back(len[v]==1?0:ch[get_fail(fail[u])][c]);
		num.push_back(num[fail[v]]+1);pos.push_back(n);siz.push_back(0);
		ch[u][c]=v;
	}
	lst=ch[u][c];siz[lst]++;return lst;
}
void build(string S){
	assert(!S.empty()&&S[0]==' '&&S.size()<INT_MAX-2);
	set((int)S.size()-1);
	for(int i=1;i<(int)S.size();i++)extend(S[i]);
}
vector<int> count()const{
	vector<int> cnt=siz;
	for(int u=tot;u>=2;u--)cnt[fail[u]]+=cnt[u];
	cnt[0]=cnt[1]=0;return cnt;
}
};
// end for string/pam.cpp
/////////////////////////
\end{lstlisting}

\Needspace{10\baselineskip}
\subsection{\texttt{string/manacher.cpp}}
\textbf{Manacher 回文半径}

\TemplUsage{线性求出所有奇数、偶数长度回文的最大半径，并以 $O(1)$ 判断某个子串是否回文。}
\TemplApi{\begin{itemize}
\item \texttt{Manacher t; t.build(s)}；s 带前导空格，实际下标 1..n；\texttt{set()} 清空。
\item \texttt{d1[i]} 包含中心字符：最长奇回文为 [i-d1[i]+1,i+d1[i]-1]，长度 2*d1[i]-1。
\item \texttt{d2[i]} 的中心在 i-1 与 i 之间：最长偶回文为 [i-d2[i],i+d2[i]-1]，长度 2*d2[i]，半径可以为 0。
\item \texttt{ask(l,r)} 判断非空闭区间是否回文，要求 $1\le l\le r\le n$。所有回文出现次数为各位置 d1+d2 的总和，不是不同回文数。
\end{itemize}}
\textbf{最小示例（放入 main）：}
\begin{lstlisting}
Manacher t; t.build(" abba");
assert(t.d2[3]==2 && t.d1[2]==1);
assert(t.ask(1,4) && t.ask(2,3));
assert(!t.ask(1,3));
long long cnt=0;
for(int i=1;i<=t.n;i++)cnt+=t.d1[i]+t.d2[i];
assert(cnt==6);
\end{lstlisting}
\TemplNote{时间、空间均为 $O(n)$；不插入分隔字符，可处理任意字节。奇偶半径含义不同，不能直接把 d1 当成回文长度。空串允许 build，但没有可查询的非空区间。求最长回文取所有 2*d1[i]-1 与 2*d2[i] 的最大值，空串答案为 0。}

\Needspace{6\baselineskip}\textbf{完整代码}
\begin{lstlisting}[firstnumber=1]
////////////////////////////////////////////////////////////////
//
// template for Manacher
// usage:
//   Manacher t; t.build(" abba"); t.ask(1,4);
//   Leading space, 1-indexed; ask(l,r) checks a nonempty closed interval.
//   Odd: [i-d1[i]+1,i+d1[i]-1]; even: [i-d2[i],i+d2[i]-1].
//   O(n) build / space, O(1) query.
//
////////////////////////////////////////////////////////////////
struct Manacher{
int n=0;
vector<int> d1,d2;
void set(){n=0;d1.clear();d2.clear();}
void build(const string &s){
	assert(!s.empty()&&s[0]==' '&&s.size()<INT_MAX/2);
	n=(int)s.size()-1;d1.assign(n+1,0);d2.assign(n+1,0);
	for(int i=1,l=1,r=0;i<=n;i++){
		int k=i>r?1:min(d1[l+r-i],r-i+1);
		while(i-k>=1&&i+k<=n&&s[i-k]==s[i+k])k++;
		d1[i]=k--;
		if(i+k>r)l=i-k,r=i+k;
	}
	for(int i=1,l=1,r=0;i<=n;i++){
		int k=i>r?0:min(d2[l+r-i+1],r-i+1);
		while(i-k-1>=1&&i+k<=n&&s[i-k-1]==s[i+k])k++;
		d2[i]=k;
		if(i+k-1>r)l=i-k,r=i+k-1;
	}
}
bool ask(int l,int r)const{
	assert(1<=l&&l<=r&&r<=n);int k=r-l+1;
	return k&1?d1[(l+r)/2]>=k/2+1:d2[(l+r+1)/2]>=k/2;
}
};
// end for string/manacher.cpp
/////////////////////////
\end{lstlisting}

\Needspace{10\baselineskip}
\subsection{\texttt{string/runs.cpp}}
\textbf{Runs 极大周期子串}

\TemplUsage{枚举所有长度至少为最小周期两倍、且不能保持该周期向左或向右扩展的子串。返回若干三元组 $(l,r,p)$，其中 p 必须是最小周期。}
\TemplApi{先抄 \texttt{string/sa.cpp} 和 \texttt{basic/ST table (linear).cpp}，再抄本文件。调用 \texttt{Runs t; auto a=t.build(s)}，返回 \texttt{vector<array<int,3>>}，依次为 l、r、p。输入有前导空格，区间为 1 下标闭区间，结果按三元组字典序排序并去重。空串、单字符及无重复的串返回空向量。}
\textbf{最小示例（放入 main）：}
\begin{lstlisting}
Runs t; auto a=t.build(" ababab");
assert((a==vector<array<int,3>>{{1,6,2}}));
a=t.build(" aabaa");
assert((a==vector<array<int,3>>{{1,2,1},{4,5,1}}));
assert(t.build(" abc").empty());
\end{lstlisting}
\TemplNote{实际字符限字节 1..255，0 留给 SA 哨兵。采用两种字母序下的最长 Lyndon 前缀，加正反串 LCP 向两侧扩展；LCP 使用四毛子，不使用哈希。源码版时间 $O(n\log n)$（最后排序），空间 $O(n)$；考场倍增 SA 也满足此界。Runs 不包括所有周期区间，只给出极大区间；长度不必是 p 的倍数，例如 ababa 返回 (1,5,2)。}

\Needspace{6\baselineskip}\textbf{完整代码}
\begin{lstlisting}[firstnumber=1]
////////////////////////////////////////////////////////////////
//
// template for Runs
// usage:
//   Paste string/sa.cpp and basic/ST table (linear).cpp first.
//   Runs t; auto a=t.build(" ababab"); // {1,6,2}
//   Leading space, bytes 1..255; returns vector<array<int,3>> of {l,r,p}.
//   [l,r] is maximal with least period p, r-l+1>=2*p; 1-indexed.
//   Sorted and unique; O(n log n) time, O(n) space.
//
////////////////////////////////////////////////////////////////
struct Runs{
private:
struct LCP{
	int n=0;
	vector<int> rk;
	Linear_RMQ<int> st;
	void build(const string &s){
		SA t;t.build(s,256);n=t.n;rk=move(t.rk);
		vector<int> h(n+1);
		for(int i=1,k=0;i<=n;i++){
			int r=rk[i];if(r==1){k=0;continue;}
			int j=t.sa[r-1];
			while(i+k<=n&&j+k<=n&&s[i+k]==s[j+k])k++;
			h[r]=k;if(k)k--;
		}
		st.build(h);
	}
	int ask(int i,int j)const{
		if(i>n||j>n)return 0;
		if(i==j)return n-i+1;
		int l=rk[i],r=rk[j];if(l>r)swap(l,r);
		return st.ask(l+1,r);
	}
};
public:
vector<array<int,3>> build(const string &s)const{
	assert(!s.empty()&&s[0]==' '&&s.size()<(INT_MAX-1024)/2);
	int n=(int)s.size()-1;
	for(int i=1;i<=n;i++)assert((unsigned char)s[i]>0);
	vector<array<int,3>> ans;if(n<2)return ans;
	LCP a,b;a.build(s);
	string t=s;reverse(t.begin()+1,t.end());b.build(t);
	vector<int> st;st.reserve(n);
	for(int op=0;op<2;op++){
		st.clear();
		for(int i=n;i>=1;i--){
			while(!st.empty()){
				int j=st.back(),k=a.ask(i,j);
				// Reversing the alphabet does not reverse the prefix rule.
				if(j+k>n)break;
				bool less=(unsigned char)s[i+k]<(unsigned char)s[j+k];
				if(less==bool(op))break;
				st.pop_back();
			}
			if(!st.empty()){
				int j=st.back(),p=j-i;
				int l=b.ask(n-i+2,n-j+2),r=a.ask(i,j);
				if(l+r>=p)ans.push_back({i-l,j+r-1,p});
			}
			st.push_back(i);
		}
	}
	sort(ans.begin(),ans.end());ans.erase(unique(ans.begin(),ans.end()),ans.end());
	return ans;
}
};
// end for string/runs.cpp
/////////////////////////
\end{lstlisting}

\Needspace{10\baselineskip}
\subsection{\texttt{string/lyndon.cpp}}
\textbf{Lyndon 分解}

\TemplUsage{用 Duval 将串唯一分解为字典序不增的 Lyndon 串。Lyndon 串严格小于自己的每个非空真后缀，相同因子可以连续出现。}
\TemplApi{\texttt{auto a=lyndon(s)}；s 带前导空格。返回 \texttt{vector<pair<int,int>>}，各项为一个因子的 1 下标闭区间 [l,r]，按原串从左到右排列且恰好覆盖全串。需要内容时用 \texttt{s.substr(l,r-l+1)}。}
\textbf{最小示例（放入 main）：}
\begin{lstlisting}
string s=" ababab"; auto a=lyndon(s);
assert((a==vector<pair<int,int>>{{1,2},{3,4},{5,6}}));
assert((lyndon(" cba")==vector<pair<int,int>>{{1,1},{2,2},{3,3}}));
assert(lyndon(" ").empty());
\end{lstlisting}
\TemplNote{时间 $O(n)$，除返回区间外额外空间 $O(1)$。比较按 unsigned char 的字节顺序，可处理任意字节；不要求小写字母。同样的因子不会合并，ababab 分为三个 ab，而 aaaa 分为四个 a。空串返回空向量。本接口是整串分解，不是每个位置的最长 Lyndon 前缀数组。}

\Needspace{6\baselineskip}\textbf{完整代码}
\begin{lstlisting}[firstnumber=1]
////////////////////////////////////////////////////////////////
//
// template for Lyndon Factorization
// usage:
//   auto a=lyndon(" ababab"); // [1,2], [3,4], [5,6]
//   Leading space, 1-indexed closed intervals, ordered from left to right.
//   The factors are Lyndon words in non-increasing lexicographic order.
//   O(n) time, O(1) extra space excluding the returned intervals.
//
////////////////////////////////////////////////////////////////
vector<pair<int,int>> lyndon(const string &s){
	assert(!s.empty()&&s[0]==' '&&s.size()<INT_MAX);
	int n=(int)s.size()-1;vector<pair<int,int>> ans;
	for(int i=1;i<=n;){
		int j=i+1,k=i;
		while(j<=n&&(unsigned char)s[k]<=(unsigned char)s[j]){
			k=s[k]==s[j]?k+1:i;j++;
		}
		int p=j-k;
		while(i<=k){ans.emplace_back(i,i+p-1);i+=p;}
	}
	return ans;
}
// end for string/lyndon.cpp
/////////////////////////
\end{lstlisting}
\clearpage
\section{线性代数}
\Needspace{10\baselineskip}
\subsection{\texttt{math/gauss.cpp}}
\textbf{线性方程组与行列式}

\TemplUsage{用消元求 Ax=b 的一组解，或计算方阵行列式。支持质数模数下的 mint 以及浮点类型；不能用 int/ll 实例化。}
\TemplApi{\begin{itemize}
\item \texttt{Gauss<T> g; res=g.solve(a,m)}，a 是 n 行 m+1 列的增广矩阵，最后一列为 b。
\item 返回 -1 无解（x 被清空），0 多解，1 唯一解；解存 g.x，长度 m，下标 0..m-1。
\item 多解时自由变量取 0；g.rank 是系数矩阵秩，where[j]=-1 表示自由变量。
\item \texttt{g.det(a)} 接受 n*n 方阵，返回 T，不改变 x/rank。空方阵行列式为 1。
\end{itemize}}
\textbf{最小示例（放入 main）：}
\begin{lstlisting}
Gauss<long double> g;
// x+y=3, 2x-y=0
int res=g.solve({{1,1,3},{2,-1,0}},2);
assert(res==1 && abs(g.x[0]-1)<1e-9L);
assert(abs(g.x[1]-2)<1e-9L);
assert(abs(g.det({{1,2},{3,4}})+2)<1e-9L);
\end{lstlisting}
\TemplNote{Gauss<mint> 需要先粘贴 ModInt 或一个内置 mint 的多项式文件。浮点使用最大绝对值选主元和绝对阈值 eps（默认 1e-12），尺度差异大或病态方程需要自行缩放并检查残差。方阵时间 O(n\textasciicircum{}3)，复制矩阵后原输入保持不变；多个 RHS 要自行扩展。}

\Needspace{6\baselineskip}\textbf{完整代码}
\begin{lstlisting}[firstnumber=1]
//
// template for Gaussian elimination and determinant
// usage:
//   Gauss<mint> G; // prime modulus, or Gauss<long double>
//   int res=G.solve(a,m); // augmented matrix, m variables
//   -1: no solution, 0: multiple, 1: unique; answer in G.x
//   free variables are set to zero; G.rank, G.where
//   G.det(a); // square matrix; floating point uses G.eps
//
template<typename T> struct Gauss{
static_assert(!is_integral_v<T>,"use a field type");
long double eps=1e-12L;
int rank=0;
vector<int> where;
vector<T> x;
bool zero(const T &v)const{
	if constexpr(is_floating_point_v<T>)return abs(v)<=eps;
	else return !v;
}
int pivot(const vector<vector<T>> &a,int r,int c)const{
	int p=-1;
	for(int i=r;i<(int)a.size();i++)if(!zero(a[i][c])){
		if constexpr(is_floating_point_v<T>){if(p==-1||abs(a[i][c])>abs(a[p][c]))p=i;}
		else return i;
	}
	return p;
}
int solve(vector<vector<T>> a,int m){
	assert(m>=0&&m<INT_MAX&&a.size()<INT_MAX&&eps>=0);
	int n=a.size();for(auto &v:a)assert(v.size()==(size_t)m+1);
	rank=0;where.assign(m,-1);x.assign(m,T(0));
	for(int c=0;c<m&&rank<n;c++){
		int p=pivot(a,rank,c);if(p==-1)continue;
		swap(a[p],a[rank]);where[c]=rank;
		T z=T(1)/a[rank][c];a[rank][c]=T(1);
		for(int j=c+1;j<=m;j++)a[rank][j]*=z;
		for(int i=rank+1;i<n;i++){
			T v=a[i][c];a[i][c]=T(0);if(zero(v))continue;
			for(int j=c+1;j<=m;j++)a[i][j]-=v*a[rank][j];
		}
		rank++;
	}
	for(int i=rank;i<n;i++)if(!zero(a[i][m])){x.clear();return -1;}
	for(int c=m-1;c>=0;c--)if(where[c]!=-1){
		int r=where[c];x[c]=a[r][m];
		for(int j=c+1;j<m;j++)x[c]-=a[r][j]*x[j];
	}
	return rank==m?1:0;
}
T det(vector<vector<T>> a)const{
	assert(a.size()<INT_MAX&&eps>=0);int n=a.size();
	for(auto &v:a)assert(v.size()==(size_t)n);
	T ans=1;
	for(int c=0;c<n;c++){
		int p=pivot(a,c,c);if(p==-1)return T(0);
		if(p!=c){swap(a[p],a[c]);ans=-ans;}
		ans*=a[c][c];T z=T(1)/a[c][c];
		for(int i=c+1;i<n;i++){
			T v=a[i][c];a[i][c]=T(0);if(zero(v))continue;v*=z;
			for(int j=c+1;j<n;j++)a[i][j]-=v*a[c][j];
		}
	}
	return ans;
}
};
\end{lstlisting}
\par\vspace{1em}

\Needspace{10\baselineskip}
\subsection{\texttt{math/linear basis/xor.cpp}}
\textbf{异或线性基}
\TemplApi{维护 64 位无符号整数的异或张成空间。\texttt{Xor\_Basis B} 默认空基；\texttt{set()} 清空，\texttt{insert(x)} 返回是否增秩，\texttt{rank} 为当前秩。\texttt{contains(x)} 判断能否异或得到 x；\texttt{ask(x=0)} 返回 x 与基中任意子集异或的最大值，允许空集。}
\textbf{最小示例（放入 main）：}
\begin{lstlisting}
Xor_Basis B;
assert(B.insert(1) && B.insert(2));
assert(!B.insert(3) && B.rank==2);
assert(B.contains(3) && !B.contains(4));
assert(B.ask()==3 && B.ask(4)==7);
\end{lstlisting}
\TemplNote{插入、查询均 O(64)，空间 O(64)。输入用 ull，不按有符号大小比较；第 63 位也可用。零向量不增秩。a 中存的是消元后的向量，不是原输入元素。}
\Needspace{6\baselineskip}\textbf{完整代码}
\begin{lstlisting}[firstnumber=1]
// Xor_Basis B; B.insert(x); B.contains(x); B.ask(x);
struct Xor_Basis{
int rank=0;
vector<ull> a=vector<ull>(64);
void set(){rank=0;fill(a.begin(),a.end(),0);}
bool insert(ull x){
	for(int i=63;i>=0;i--)if(x>>i&1){
		if(a[i])x^=a[i];
		else{a[i]=x;rank++;return true;}
	}
	return false;
}
bool contains(ull x)const{
	for(int i=63;i>=0;i--)if(x>>i&1)x^=a[i];
	return x==0;
}
ull ask(ull x=0)const{
	for(int i=63;i>=0;i--)x=max(x,x^a[i]);
	return x;
}
};
\end{lstlisting}
\par\vspace{1em}
\Needspace{10\baselineskip}
\subsection{\texttt{math/linear basis/xor-prefix.cpp}}
\textbf{前缀异或线性基}
\TemplApi{\texttt{Prefix\_Xor\_Basis B}；按加入顺序调用 \texttt{insert(x,id)}，id 必须为严格递增的正整数，可不连续。增秩返回 0，否则返回本次相关关系中最早且可删除的元素编号；零向量返回自身编号。\texttt{set()} 清空并重置编号。}
\textbf{最小示例（放入 main）：}
\begin{lstlisting}
Prefix_Xor_Basis B;
assert(B.insert(1,1)==0 && B.insert(2,2)==0);
assert(B.insert(2,3)==2); // removing id=1 would lower rank
assert(B.insert(3,4)==1);
assert(B.get_rank(3)==2 && B.contains(1,3));
assert(!B.contains(1,4) && B.ask(0,4)==3);
assert(B.insert(0,5)==5);
\end{lstlisting}
\TemplNote{处理到右端点 r 后，\texttt{get\_rank(l)}、\texttt{contains(x,l)}、\texttt{ask(x,l)} 分别查询编号 [l,r] 的秩、可表示性和最大异或值。l 默认 1；l 大于 r 时视为空基。\texttt{rank} 是全部已插入元素的秩。pos 中的正编号是保留的原元素编号，a 是消元后的基。单次 O(64)，空间 O(64)；查询历史右端点须离线按 r 扫描或自己保存副本，没有持久化。}
\Needspace{6\baselineskip}\textbf{完整代码}
\begin{lstlisting}[firstnumber=1]
// Prefix_Xor_Basis B; int removed=B.insert(x,id);
// id must increase; 0: rank increased, otherwise removed id
struct Prefix_Xor_Basis{
int rank=0,last=0;
vector<ull> a=vector<ull>(64);
vector<int> pos=vector<int>(64);
void set(){rank=last=0;fill(a.begin(),a.end(),0);fill(pos.begin(),pos.end(),0);}
int insert(ull x,int id){
	assert(id>last);last=id;
	for(int i=63;i>=0;i--)if(x>>i&1){
		if(!a[i]){a[i]=x;pos[i]=id;rank++;return 0;}
		if(pos[i]<id){swap(a[i],x);swap(pos[i],id);}
		x^=a[i];
	}
	return id;
}
int get_rank(int l=1)const{
	assert(l>=1);int ans=0;
	for(int i=0;i<64;i++)ans+=pos[i]>=l;
	return ans;
}
bool contains(ull x,int l=1)const{
	assert(l>=1);
	for(int i=63;i>=0;i--)if((x>>i&1)&&pos[i]>=l)x^=a[i];
	return x==0;
}
ull ask(ull x=0,int l=1)const{
	assert(l>=1);
	for(int i=63;i>=0;i--)if(pos[i]>=l)x=max(x,x^a[i]);
	return x;
}
};
\end{lstlisting}
\par\vspace{1em}
\Needspace{10\baselineskip}
\subsection{\texttt{math/linear basis/linear.cpp}}
\textbf{模意义线性基}
\TemplApi{先抄一份 ModInt；\texttt{Linear\_Basis<mint> B; B.setN(n)} 设置向量维数。\texttt{insert(x)} 返回是否增秩；\texttt{contains(x)} 判断 x 是否在张成空间中；\texttt{rank} 是当前秩。输入 vector 长度必须为 n，坐标下标 0..n-1。}
\textbf{最小示例（放入 main）：}
\begin{lstlisting}
Linear_Basis<mint> B; B.setN(2);
assert(B.insert({1,0}) && B.insert({0,1}));
assert(!B.insert({2,3}) && B.rank==2);
assert(B.contains({7,9}));
B.setN(2); assert(!B.contains({1,0}));
\end{lstlisting}
\TemplNote{T 要求精确域上的运算，通常使用质数模数 mint；不支持普通 int/ll、合数模数或浮点数。插入和查询 O(n²)，空间 O(n²)，占用的主元行按实际长度存储。n=0 允许，空向量不增秩。只维护张成空间，不返回具体线性组合系数。}
\Needspace{6\baselineskip}\textbf{完整代码}
\begin{lstlisting}[firstnumber=1]
// Linear_Basis<mint> B; B.setN(n); B.insert(x); B.contains(x);
// T must be an exact field type; coordinates are 0..n-1
template<typename T> struct Linear_Basis{
static_assert(!is_integral_v<T>&&!is_floating_point_v<T>,"use an exact field type");
int n=0,rank=0;
vector<vector<T>> a;
void setN(int _n){assert(_n>=0);n=_n;rank=0;a.assign(n,{});}
bool insert(vector<T> x){
	assert(x.size()==(size_t)n);
	for(int i=n-1;i>=0;i--)if(x[i]){
		if(a[i].empty()){
			T v=T(1)/x[i];for(int j=0;j<i;j++)x[j]*=v;
			x[i]=T(1);x.resize(i+1);a[i]=move(x);rank++;return true;
		}
		T v=x[i];for(int j=0;j<i;j++)x[j]-=v*a[i][j];
		x[i]=T(0);
	}
	return false;
}
bool contains(vector<T> x)const{
	assert(x.size()==(size_t)n);
	for(int i=n-1;i>=0;i--)if(x[i]){
		if(a[i].empty())return false;
		T v=x[i];for(int j=0;j<i;j++)x[j]-=v*a[i][j];
		x[i]=T(0);
	}
	return true;
}
};
\end{lstlisting}
\par\vspace{1em}
\Needspace{10\baselineskip}
\subsection{\texttt{math/linear basis/linear-prefix.cpp}}
\textbf{前缀模意义线性基}
\TemplApi{先抄 ModInt。\texttt{Prefix\_Linear\_Basis<mint> B; B.setN(n)}，依次 \texttt{insert(x,id)}。id 必须为严格递增的正整数；增秩返回 0，否则返回最早可替换元素的编号。规则与前缀异或版相同，只是异或消元换成域上的线性消元。}
\textbf{最小示例（放入 main）：}
\begin{lstlisting}
Prefix_Linear_Basis<mint> B; B.setN(2);
assert(B.insert({1,0},1)==0);
assert(B.insert({0,1},2)==0);
assert(B.insert({0,2},3)==2);
assert(B.insert({1,2},4)==1);
assert(B.get_rank(3)==2 && B.contains({1,0},3));
assert(!B.contains({1,0},4));
assert(B.insert({0,0},5)==5);
\end{lstlisting}
\TemplNote{处理到编号 r 时，\texttt{get\_rank(l)} 查询 [l,r] 的秩，\texttt{contains(x,l)} 判断 x 是否在该区间张成空间中。l 默认 1，超出 r 时为空基。插入和可表示性查询 O(n²)，秩查询 O(n)，空间 O(n²)。只用于精确域，通常用质数模数 mint；setN 清空所有状态。pos 保留原编号，a 是消元后的向量。不是任意删除或持久化结构。}
\Needspace{6\baselineskip}\textbf{完整代码}
\begin{lstlisting}[firstnumber=1]
// Prefix_Linear_Basis<mint> B; B.setN(n); int removed=B.insert(x,id);
// id must increase; 0: rank increased, otherwise removed id
template<typename T> struct Prefix_Linear_Basis{
static_assert(!is_integral_v<T>&&!is_floating_point_v<T>,"use an exact field type");
int n=0,rank=0,last=0;
vector<vector<T>> a;
vector<int> pos;
void setN(int _n){assert(_n>=0);n=_n;rank=last=0;a.assign(n,{});pos.assign(n,0);}
int insert(vector<T> x,int id){
	assert(x.size()==(size_t)n&&id>last);last=id;
	for(int i=n-1;i>=0;i--)if(x[i]){
		if(a[i].empty()){
			T v=T(1)/x[i];for(int j=0;j<i;j++)x[j]*=v;
			x[i]=T(1);x.resize(i+1);a[i]=move(x);pos[i]=id;rank++;return 0;
		}
		if(pos[i]<id){
			x.resize(i+1);swap(a[i],x);swap(pos[i],id);
			T v=T(1)/a[i][i];for(int j=0;j<i;j++)a[i][j]*=v;
			a[i][i]=T(1);
		}
		T v=x[i];for(int j=0;j<i;j++)x[j]-=v*a[i][j];
		x[i]=T(0);
	}
	return id;
}
int get_rank(int l=1)const{
	assert(l>=1);int ans=0;for(int p:pos)ans+=p>=l;return ans;
}
bool contains(vector<T> x,int l=1)const{
	assert(x.size()==(size_t)n&&l>=1);
	for(int i=n-1;i>=0;i--)if(x[i]){
		if(pos[i]<l)return false;
		T v=x[i];for(int j=0;j<i;j++)x[j]-=v*a[i][j];
		x[i]=T(0);
	}
	return true;
}
};
\end{lstlisting}
\par\vspace{1em}
\clearpage
\section{计算几何}
\Needspace{10\baselineskip}
\subsection{\texttt{geometry/geometry-main.cpp}}
\textbf{整数二维几何基础}

\TemplUsage{点坐标和叉积均为 ll，提供凸包、面积与基本相交判断；需要另定义 ld=long double。}
\TemplApi{\begin{itemize}
\item \texttt{Point(x,y)}；点向量加减；\texttt{a*b} 是叉积，\texttt{a\textasciicircum{}b} 是点积。
\item \texttt{convex(points)} 求凸包，\texttt{convec\_area(hull)} 返回两倍面积。
\item \texttt{on\_segment}、strict\_intersect、nonstrict\_intersect 分别判断在线段上、严格相交、含端点相交。
\end{itemize}}
\textbf{粘贴本文件之前：}
\begin{lstlisting}
using ld=long double;
\end{lstlisting}
\textbf{最小示例（放入 main）：}
\begin{lstlisting}
vector<Point> a{{0,0},{2,0},{2,2},{0,2}};
auto h=convex(a); assert(convec_area(h)==8);
assert(on_segment(Point(1,0),Point(0,0),Point(2,0)));
\end{lstlisting}
\TemplNote{convex 自动排序去重。min\_cross 要求顶点按凸包环序排列且没有共线中间点，支持旋转起点、顺逆时针，O(log n)。射线包含端点，零方向按单点处理；退化三角形按三条边处理。坐标叉积、点积仍须保证不溢出 ll。Point 输出两个坐标，中间有空格；output 用于调试。}

\Needspace{6\baselineskip}\textbf{完整代码}
\begin{lstlisting}[firstnumber=1]
//
// template for 2D geometry
// version 1.0 (Last Update Jun 16th, 2026)
//
// usage:
//   Point A,B; convex(hull); convec_area(hull);
//   needs bits/stdc++.h, typedef ll/ld
//
const ld pi=acos(-1),dx=1e-9;
struct Point{
	ll x,y;
	friend istream& operator >>(istream &i,Point &P){i>>P.x>>P.y;return i;}
	friend ostream& operator <<(ostream &o,const Point &P){o<<P.x<<' '<<P.y;return o;}
	// void read(){cin>>x>>y;}
	// void output(){cout<<x<<' '<<y<<'\n';}
	void output(string s=""){cerr<<"Point "<<s<<": ("<<x<<","<<y<<")"<<endl;}
	double arg(){return atan2(y,x);}
	Point (ll _x=0,ll _y=0){x=_x,y=_y;}
};

inline bool operator ==(const Point &x,const Point &y){return x.x==y.x&&x.y==y.y;}
inline bool operator !=(const Point &x,const Point &y){return x.x!=y.x||x.y!=y.y;}

inline bool operator <(const Point &x,const Point &y){return (x.x!=y.x?x.x<y.x:x.y<y.y);}
inline Point operator -(const Point &x,const Point &y){return Point(x.x-y.x,x.y-y.y);}
inline Point operator +(const Point &x,const Point &y){return Point(x.x+y.x,x.y+y.y);}
inline void operator -=(Point &x,const Point &y){x=x-y;}
inline void operator +=(Point &x,const Point &y){x=x+y;}
inline Point operator *(const Point &x,ll y){return Point(x.x*y,x.y*y);}
inline void operator *=(Point &x,ll y){x=x*y;}
inline Point operator *(ll y,const Point &x){return Point(x.x*y,x.y*y);}
inline void operator *=(ll y,Point &x){x=x*y;}

inline ll operator *(const Point &x,const Point &y){
	// return OX · OY · sin(XOY) = 2 * area(OXY)
	return x.x*y.y-x.y*y.x;
}
inline ll operator ^(const Point &x,const Point &y){
	// return OX · OY · cos(XOY)
	return x.x*y.x+x.y*y.y;
}

int sgn(ll x){return (x==0?0:(x>0?1:-1));}

inline int sgn(const Point &P){
	return P.y<0?-1:P.y>0||P.x<0?1:0;
}
inline bool cmp(const Point &x,const Point &y){
	if(sgn(x)!=sgn(y)) return sgn(x)<sgn(y);
	return x*y>0;
}
inline ll area(const Point &x,const Point &y,const Point &z){
	// return triangle XYZ area **times 2**
	return abs((x-y)*(x-z));
}

inline ld dist(const Point &x){return sqrt((ld)x.x*x.x+(ld)x.y*x.y);}
inline ll dist2(const Point &x){return x.x*x.x+x.y*x.y;}

inline int dir(const Point &x,const Point &y,const Point &z){
	// return is YX clockwise direction to YZ 
	return sgn((x-y)*(y-z));
}
inline int suf(int x,int n){return (x+1==n?0:x+1);}
inline int pre(int x,int n){return (x==0?n-1:x-1);}
ll convec_area(const vector<Point> &a){
	// return convex A area **times 2**
	ll s=0;int n=a.size();
	for(int i=0;i<n;i++) s+=area(a[0],a[i],a[suf(i,n)]);
	return s;
}
vector<Point> convex(vector<Point> a){
	sort(a.begin(),a.end());a.erase(unique(a.begin(),a.end()),a.end());
	int n=a.size();
	if(n<=2)return a;
	vector<Point> b;int s=0,t=0;
	for(int i=2;i--;s=t,reverse(a.begin(),a.end())){
		for(Point p:a){
			while(t>=s+2&&dir(b[t-2],b[t-1],p)>=0) t--,b.pop_back();
			b.push_back(p);t++;
		}
		b.pop_back();t--;
	} 
	return b;
}
inline bool on_segment(const Point &X,const Point &A,const Point &B){
	// check if X on segment AB
	if((A-B)*(X-B)!=0) return false;
	return min(A.x,B.x)<=X.x&&X.x<=max(A.x,B.x)&&
		min(A.y,B.y)<=X.y&&X.y<=max(A.y,B.y);
}
inline bool strict_intersect(const Point &A,const Point &B,
	const Point &C,const Point &D){
	// check if AB,CD (except endpoint)strictly intersect
	return dir(A,C,D)*dir(B,C,D)==-1&&dir(C,A,B)*dir(D,A,B)==-1;
}
inline bool nonstrict_intersect(const Point &A,const Point &B,
	const Point &C,const Point &D){
	// check if AB,CD (include endpoints) intersect
	return (dir(A,C,D)*dir(B,C,D)==-1&&dir(C,A,B)*dir(D,A,B)==-1)||
	on_segment(A,C,D)||on_segment(B,C,D)||on_segment(C,A,B)||on_segment(D,A,B);
}
inline bool parallel(Point A,Point B,
	Point C,Point D){
	// check if AB // CD
	B-=A,D-=C;A-=C;C=-1*D;
	return sgn(B*C)==0&&sgn(A*B)!=0;
}
inline bool ray_intersect(Point A,Point B,Point C,Point D){
	// closed rays; a zero direction is a single point
	Point u=B-A,v=D-C,w=C-A;
	if(u==Point())return v==Point()?A==C:((A-C)*v==0&&((A-C)^v)>=0);
	if(v==Point())return w*u==0&&(w^u)>=0;
	ll d=u*v;
	if(!d)return w*u==0&&((u^v)>0||(w^u)>=0);
	return sgn(w*v)*sgn(d)>=0&&sgn(w*u)*sgn(d)>=0;
}
inline bool in_triangle(const Point &P,const Point &A,
	const Point &B,const Point &C){
	// check if P in triangle ABC
	// P may lie on bound
	if(!area(A,B,C))return on_segment(P,A,B)||on_segment(P,B,C)||on_segment(P,C,A);
	return area(A,B,C)==area(P,A,B)+area(P,B,C)+area(P,C,A);
}
template<typename F> int convex_min(int n,F val){
	assert(n>0);
	auto low=[&](int i){return val(i)<=val((i+n-1)%n)&&val(i)<=val((i+1)%n);};
	if(low(0))return 0;
	int l=0,r=n;
	while(l+1<r){
		int mid=(l+r)>>1;
		if(low(mid))return mid;
		bool a=val((l+1)%n)>=val(l),b=val((mid+1)%n)>=val(mid);
		if(a!=b){if(a)l=mid;else r=mid;}
		else if((val(mid)>val(l))==a)l=mid;
		else r=mid;
	}
	return r%n;
}
inline ll min_cross(const vector<Point> &I,Point X){
	// convex polygon in cyclic order, without collinear intermediate points
	auto val=[&](int i){return I[i]*X;};
	return val(convex_min(I.size(),val));
}
inline ld min_cross(const vector<Point> &I,ld x,ld y){
	auto val=[&](int i){return I[i].x*x+I[i].y*y;};
	return val(convex_min(I.size(),val));
}
\end{lstlisting}
\par\vspace{1em}

\end{document}
